\documentclass[12pt,a4paper]{article}
\usepackage[utf8]{inputenc}
\usepackage[T1]{fontenc}
\usepackage{hyperref}
\usepackage{color}
\usepackage{fancyvrb}
\newcommand{\VerbBar}{|}
\newcommand{\VERB}{\Verb[commandchars=\\\{\}]}
\DefineVerbatimEnvironment{Highlighting}{Verbatim}{commandchars=\\\{\}}
% Add ',fontsize=\small' for more characters per line
\newenvironment{Shaded}{}{}
\newcommand{\AlertTok}[1]{\textcolor[rgb]{1.00,0.00,0.00}{\textbf{#1}}}
\newcommand{\AnnotationTok}[1]{\textcolor[rgb]{0.38,0.63,0.69}{\textbf{\textit{#1}}}}
\newcommand{\AttributeTok}[1]{\textcolor[rgb]{0.49,0.56,0.16}{#1}}
\newcommand{\BaseNTok}[1]{\textcolor[rgb]{0.25,0.63,0.44}{#1}}
\newcommand{\BuiltInTok}[1]{\textcolor[rgb]{0.00,0.50,0.00}{#1}}
\newcommand{\CharTok}[1]{\textcolor[rgb]{0.25,0.44,0.63}{#1}}
\newcommand{\CommentTok}[1]{\textcolor[rgb]{0.38,0.63,0.69}{\textit{#1}}}
\newcommand{\CommentVarTok}[1]{\textcolor[rgb]{0.38,0.63,0.69}{\textbf{\textit{#1}}}}
\newcommand{\ConstantTok}[1]{\textcolor[rgb]{0.53,0.00,0.00}{#1}}
\newcommand{\ControlFlowTok}[1]{\textcolor[rgb]{0.00,0.44,0.13}{\textbf{#1}}}
\newcommand{\DataTypeTok}[1]{\textcolor[rgb]{0.56,0.13,0.00}{#1}}
\newcommand{\DecValTok}[1]{\textcolor[rgb]{0.25,0.63,0.44}{#1}}
\newcommand{\DocumentationTok}[1]{\textcolor[rgb]{0.73,0.13,0.13}{\textit{#1}}}
\newcommand{\ErrorTok}[1]{\textcolor[rgb]{1.00,0.00,0.00}{\textbf{#1}}}
\newcommand{\ExtensionTok}[1]{#1}
\newcommand{\FloatTok}[1]{\textcolor[rgb]{0.25,0.63,0.44}{#1}}
\newcommand{\FunctionTok}[1]{\textcolor[rgb]{0.02,0.16,0.49}{#1}}
\newcommand{\ImportTok}[1]{\textcolor[rgb]{0.00,0.50,0.00}{\textbf{#1}}}
\newcommand{\InformationTok}[1]{\textcolor[rgb]{0.38,0.63,0.69}{\textbf{\textit{#1}}}}
\newcommand{\KeywordTok}[1]{\textcolor[rgb]{0.00,0.44,0.13}{\textbf{#1}}}
\newcommand{\NormalTok}[1]{#1}
\newcommand{\OperatorTok}[1]{\textcolor[rgb]{0.40,0.40,0.40}{#1}}
\newcommand{\OtherTok}[1]{\textcolor[rgb]{0.00,0.44,0.13}{#1}}
\newcommand{\PreprocessorTok}[1]{\textcolor[rgb]{0.74,0.48,0.00}{#1}}
\newcommand{\RegionMarkerTok}[1]{#1}
\newcommand{\SpecialCharTok}[1]{\textcolor[rgb]{0.25,0.44,0.63}{#1}}
\newcommand{\SpecialStringTok}[1]{\textcolor[rgb]{0.73,0.40,0.53}{#1}}
\newcommand{\StringTok}[1]{\textcolor[rgb]{0.25,0.44,0.63}{#1}}
\newcommand{\VariableTok}[1]{\textcolor[rgb]{0.10,0.09,0.49}{#1}}
\newcommand{\VerbatimStringTok}[1]{\textcolor[rgb]{0.25,0.44,0.63}{#1}}
\newcommand{\WarningTok}[1]{\textcolor[rgb]{0.38,0.63,0.69}{\textbf{\textit{#1}}}}
\usepackage{enumitem}
\usepackage{longtable}
\usepackage{booktabs}
\usepackage[margin=1in]{geometry}

\date{\today}

\begin{document}

\section{📘 KONU 1: Cargo ve Proje Yapısı}\label{konu-1-cargo-ve-proje-yapısı}

\textbf{Kategori:} Araçlar \textbar{} \textbf{Sürüm:} Rust 1.96.0

\begin{center}\rule{0.5\linewidth}{0.5pt}\end{center}

\subsection{🟢 KATMAN 1 --- ILK OKUMA}\label{katman-1-—-ilk-okuma}

\begin{Shaded}
\begin{Highlighting}[]
\CommentTok{\# Terminalde (proje klasörünün DIŞINDA) çalıştır:}
\NormalTok{cargo new hello\_cargo     }\CommentTok{\# yeni proje oluşturur}
\BuiltInTok{cd}\NormalTok{ hello\_cargo            }\CommentTok{\# Cargo.toml'un olduğu klasöre gir}
\NormalTok{cargo run                 }\CommentTok{\# derle + çalıştır}
\end{Highlighting}
\end{Shaded}

\begin{Shaded}
\begin{Highlighting}[]
\CommentTok{// src/main.rs — cargo new bunu otomatik oluşturur}
\KeywordTok{fn} \FunctionTok{main}\OperatorTok{()} \OperatorTok{\{}
\NormalTok{    println}\OperatorTok{!(}\StringTok{"Hello, Cargo!"}\OperatorTok{);}   \CommentTok{// ekrana yazar}
\OperatorTok{\}}
\CommentTok{// Çıktı: Hello, Cargo!}
\end{Highlighting}
\end{Shaded}

\texttt{cargo\ new} komutu bir package (paket) oluşturur. İçinde manifest (bildirim dosyası) olan \texttt{Cargo.toml} ve kaynak kod olan \texttt{src/main.rs} bulunur. \texttt{cargo\ run} önce compiler (derleyici) ile build (derleme) yapar, sonra oluşan binary'yi (çalıştırılabilir dosya) çalıştırır. Çıktılar \texttt{target/debug/} altına yazılır ve bu klasörü elle düzenlemezsin. \texttt{cargo\ check} binary üretmeden sadece hataları kontrol eder, bu yüzden \texttt{cargo\ build}`den hızlıdır.

\textbf{1.96.0 notu:} x86\_64 Linux'ta LLD artık varsayılan linker (bağlayıcı); özellikle büyük binary'lerde ve incremental rebuild'lerde linking daha hızlı. Komutların kendisi değişmedi.

\begin{center}\rule{0.5\linewidth}{0.5pt}\end{center}

\subsection{🔵 KATMAN 2 --- RECALL KARTLARI}\label{katman-2-—-recall-kartlari}

\begin{itemize}
\tightlist
\item
  \texttt{cargo\ new\ hello\_cargo} → yeni klasör + \texttt{Cargo.toml} + \texttt{src/main.rs} + \texttt{.gitignore} + git repo
  İlgili: \texttt{cargo\ init}
\item
  \texttt{cargo\ init} → \textbf{mevcut} klasörü Cargo projesi yapar
  Hata: \texttt{error:\ cargo\ init\ cannot\ be\ run\ on\ existing\ Cargo\ packages}
\item
  \texttt{cargo\ check} → binary üretmez, sadece hata kontrolü (en hızlısı)
  İlgili: \texttt{cargo\ build}
\item
  \texttt{cargo\ build} → \texttt{target/debug/hello\_cargo} üretir; ilk build'de \texttt{Cargo.lock} oluşur
\item
  \texttt{cargo\ run} → build (gerekirse) + çalıştır
  İlgili: \texttt{cargo\ run\ --bin\ isim}
\item
  \texttt{cargo\ build\ --release} → \texttt{target/release/} altına optimize binary
  İlgili: profile (profil), \texttt{opt-level}
\item
  \texttt{Cargo.toml} iskeleti:

\begin{verbatim}
[package]
name = "hello_cargo"
version = "0.1.0"
edition = "2024"    # edition (sürüm dönemi)

[dependencies]
\end{verbatim}
\item
  \texttt{cargo\ add\ rand} → \texttt{{[}dependencies{]}} altına dependency (bağımlılık) ekler
  İlgili: \texttt{cargo\ remove\ rand}
\item
  \texttt{cargo\ clean} → \texttt{target/} klasörünü siler
\item
  Hata: \texttt{could\ not\ find\ Cargo.toml\ in\ current\ directory\ or\ any\ parent\ directory} → yanlış klasördesin
\item
  Hata: \texttt{E0601} → crate (kutu) içinde \texttt{fn\ main} yok
\end{itemize}

\textbf{GoT bağlantı haritası:}

\begin{verbatim}
cargo new ──┬── Cargo.toml ──┬── [package] ──► edition
            │                └── [dependencies] ◄── cargo add
            └── src/main.rs ──► cargo check ──► cargo build ──► cargo run
                                                     │
                                                     └──► --release
\end{verbatim}

\begin{center}\rule{0.5\linewidth}{0.5pt}\end{center}

\subsection{🟣 KATMAN 3 --- DERIN TEKNIK NOT}\label{katman-3-—-derin-teknik-not}

\emph{(Bu konuda stack/heap yok, bu yüzden ``bellek'' adımı diske ve process'e uyarlandı.)}

\textbf{Adım 1 --- Disk ve Process Seviyesi}

\begin{itemize}
\tightlist
\item
  Önce soru: ``Diskte ne oluşuyor?'' Cevap: \texttt{target/debug/} altında binary, \texttt{deps/} (ara dosyalar), \texttt{incremental/} (artımlı derleme cache'i) ve \texttt{.fingerprint/}.
\item
  \texttt{Cargo.lock} kaynak koduyla birlikte tutulur ve tüm dependency'lerin tam sürümlerini kilitler. Böylece reproducible build (tekrarlanabilir derleme) mümkün olur.
\item
  Binary çalışınca OS bir process açar ve stack'i kurar. \texttt{fn\ main}`i doğrudan OS değil, Rust'ın std runtime başlangıç kodu (\texttt{lang\_start}) çağırır.
\item
  Debug ve release farkı:
\end{itemize}

{\def\LTcaptype{none} % do not increment counter
\begin{longtable}[]{@{}lll@{}}
\toprule\noalign{}
& debug (\texttt{cargo\ build}) & release (\texttt{--release}) \\
\midrule\noalign{}
\endhead
\bottomrule\noalign{}
\endlastfoot
\texttt{opt-level} & 0 & 3 \\
debug-assertions & açık & kapalı \\
overflow-checks & açık (integer overflow (tam sayı taşması) → panic) & kapalı (wrap eder) \\
Derleme süresi & kısa & uzun \\
Çalışma hızı & yavaş & hızlı \\
\end{longtable}
}

\textbf{Adım 2 --- Compiler Davranışı}

\begin{itemize}
\tightlist
\item
  Önce soru: ``Kim derliyor?'' Cevap: Cargo derlemez, her crate için \texttt{rustc} çağırır.
\item
  \texttt{cargo\ build\ -v} ile çağrılan \texttt{rustc} komutunu görebilirsin (\texttt{--crate-name}, \texttt{--edition=2024}, \texttt{--crate-type\ bin} gibi bayraklar).
\item
  \texttt{cargo\ check}, \texttt{rustc}`ı sadece metadata üretecek şekilde çalıştırır; code generation ve linking atlanır. Hızı buradan gelir.
\item
  Dosya değişmediyse Cargo yeniden derlemez (fingerprint kontrolü).
\item
  MIR görmek için: \texttt{cargo\ rustc\ --\ --emit=mir} (çıktı \texttt{target/debug/deps/} altına yazılır).
\end{itemize}

\textbf{Adım 3 --- Hata Kodları}

\begin{itemize}
\tightlist
\item
  \textbf{E0601:} binary crate'te \texttt{main} fonksiyonu yok. \texttt{fn\ main()\ \{\}} ekle veya dosya adını/konumunu kontrol et. Detay için \texttt{rustc\ --explain\ E0601}.
\item
  \textbf{Cargo hatası} (E kodu yok): \texttt{no\ targets\ specified\ in\ the\ manifest} → \texttt{src/main.rs} veya \texttt{src/lib.rs} eksik.
\item
  \textbf{Cargo hatası:} \texttt{feature\ edition2024\ is\ required} → Cargo eski. \texttt{rustup\ update} ile çöz.
\end{itemize}

\textbf{Adım 4 --- ``Neden Böyle?''}

\begin{itemize}
\tightlist
\item
  Çözülen problem: build, dependency yönetimi ve proje düzeni için tek standart araç.
\item
  C/C++`ta Make/CMake yazarsın, kütüphaneleri elle bulur ve bağlarsın. Rust'ta \texttt{Cargo.toml} + \texttt{cargo\ build} yeter.
\item
  Trade-off: Rust'ta derleme yavaş, çalışma hızlı. Bu yüzden iki profil var: geliştirirken debug, yayınlarken release.
\end{itemize}

\textbf{Adım 5 --- Rust 1.96.0 Farkı}

\begin{itemize}
\tightlist
\item
  Rust 1.96.0'daki değişiklikler (yeni Range* tipleri, \texttt{assert\_matches!} makrosu, Wasm linker güncellemeleri) temel Cargo proje yapısını değiştirmiyor. Wasm hedeflerinde linker davranışı güncellendi.
\item
  Yeni Range* tipleri ve \texttt{assert\_matches!} makrosu \texttt{cargo\ publish\ --workspace} kullanımını değiştirmiyor.
\item
  x86\_64-apple-darwin, Tier 1'den Tier 2'ye indirildi; derlenmiş sürümler dağıtılmaya devam ediyor.
\end{itemize}

\begin{center}\rule{0.5\linewidth}{0.5pt}\end{center}

\subsection{🔴 KATMAN 4 --- SIK UNUTULANLAR + DERINLEŞTIRME}\label{katman-4-—-sik-unutulanlar-derinleştirme}

\textbf{Sık Unutulan 1: \texttt{check} / \texttt{build} / \texttt{run} farkı}

\begin{itemize}
\tightlist
\item
  Neden unutuluyor: üçü de ``derliyor'' gibi duruyor.
\item
  Yol 1 (analoji): check = imla kontrolü, build = kitabı bas, run = bas + oku.
\item
  Yol 2 (kod): \texttt{cargo\ check} sonrası \texttt{target/debug/hello\_cargo} oluşmaz.
\item
  Yol 3 (hata): \texttt{./target/debug/hello\_cargo} → \texttt{No\ such\ file\ or\ directory} (sadece check yaptıysan)
\end{itemize}

\textbf{Sık Unutulan 2: Komut nerede çalışıyor?}

\begin{itemize}
\tightlist
\item
  Neden unutuluyor: \texttt{cargo\ new} projenin dışında, \texttt{cargo\ run} içinde çalışır.
\item
  Yol 1 (görsel): \texttt{Cargo.toml}`u gören klasör = çalışma klasörü.
\item
  Yol 2 (kod): \texttt{cd\ hello\_cargo\ \&\&\ cargo\ run}
\item
  Yol 3 (hata): \texttt{could\ not\ find\ Cargo.toml\ in\ current\ directory\ or\ any\ parent\ directory}
\end{itemize}

\textbf{Sık Unutulan 3: Release binary nerede?}

\begin{itemize}
\tightlist
\item
  Neden unutuluyor: \texttt{cargo\ run\ --release} çalıştırır, ama yolunu söylemez.
\item
  Yol 1 (analoji): \texttt{debug} = taslak klasörü, \texttt{release} = baskı klasörü.
\item
  Yol 2 (kod): \texttt{./target/release/hello\_cargo}
\item
  Yol 3 (hata): \texttt{--release} olmadan \texttt{target/release/} boştur.
\end{itemize}

\textbf{Hazır Derinleştirme Prompt'u:}

\begin{verbatim}
Rust 1.96.0'da Cargo ve proje yapısı konusunu en detaylı şekilde anlat. Şunları kapsa:

1. DİSK/PROCESS SEVİYESİ: target/ klasöründe ne var, Cargo.lock ne işe yarar,
   binary çalışınca process'te ne oluyor?

2. COMPILER DAVRANIŞI: Cargo, rustc'ı nasıl çağırıyor? cargo build -v çıktısını
   satır satır açıkla. check, build, release farkı nedir? Incremental compilation
   nasıl çalışıyor?

3. HATA KODLARI: E0601 ve sık görülen Cargo hataları. Her biri için:
   ne zaman, neden, nasıl çözülür. `rustc --explain E0601` çıktısını da ekle.

4. C/C++ KARŞILAŞTIRMASI: Make/CMake ile aynı iş nasıl yapılır?
   Cargo hangi problemi çözüyor?

5. ANALOJİ: Günlük hayattan bir benzetme yap (örn: mutfak, fabrika hattı).

6. RUST 1.96.0: LLD varsayılan linker değişikliği ve cargo publish --workspace
   gibi yeniliklerin günlük kullanıma etkisi.

7. PRATİK: 3 tane örnek ver — kolay, orta, zor.

Açıklama Türkçe, teknik terimler İngilizce, kod yorumları Türkçe olsun.
Her bölümü ayrı başlıkla ver. Komutları ve kodu çalıştırılabilir yaz.
\end{verbatim}

\begin{center}\rule{0.5\linewidth}{0.5pt}\end{center}

\subsection{📝 GÖREV (Boş terminal + boş dosya)}\label{görev-boş-terminal-boş-dosya}

\textbf{🟢 Kolay}

\begin{enumerate}
\def\labelenumi{\arabic{enumi}.}
\tightlist
\item
  \texttt{cargo\ new\ rust\_notes} aç, içine gir, \texttt{cargo\ run} ile çalıştır.
\item
  \texttt{Cargo.toml} içindeki \texttt{name}, \texttt{version}, \texttt{edition} satırlarını bul.
\end{enumerate}

\textbf{🟡 Orta}
3. \texttt{main.rs}`e 2 tane \texttt{println!} ekle. Sırayla \texttt{cargo\ check}, \texttt{cargo\ build}, \texttt{cargo\ run} çalıştır.
4. \texttt{cargo\ add\ rand} ile dependency ekle, \texttt{Cargo.toml}`da gör, sonra \texttt{cargo\ remove\ rand} ile sil.

\textbf{🔴 Zor}
5. \texttt{src/bin/second.rs} oluştur (içinde \texttt{fn\ main}), \texttt{cargo\ run\ --bin\ second} ile çalıştır. Sonra \texttt{cargo\ build\ --release} yap ve debug/release binary boyutlarını \texttt{ls\ -lh\ target/debug\ target/release} ile karşılaştır.

\textbf{Başarı kriteri:} Hepsini hatasız yaparsan → öğrendin ✅

\subsubsection{✅ Kendini Test Et (2-3 saat sonra, sıfırdan)}\label{kendini-test-et-2-3-saat-sonra-sıfırdan}

\begin{enumerate}
\def\labelenumi{\arabic{enumi}.}
\tightlist
\item
  Proje oluştur ve çalıştır.
\item
  \texttt{Cargo.toml} iskeletini ezberden yaz.
\item
  \texttt{check} / \texttt{build} / \texttt{run} / \texttt{--release} farkını söyle ve binary yollarını yaz.
\end{enumerate}

\begin{center}\rule{0.5\linewidth}{0.5pt}\end{center}

\subsection{📚 KAYNAKÇA}\label{kaynakça}

\subsubsection{Resmi Kaynaklar}\label{resmi-kaynaklar}

\begin{itemize}
\tightlist
\item
  \href{https://doc.rust-lang.org/book/ch01-03-hello-cargo.html}{Rust Book --- Hello, Cargo!}
\item
  \href{https://doc.rust-lang.org/cargo/guide/project-layout.html}{Cargo Book --- Package Layout}
\item
  \href{https://doc.rust-lang.org/cargo/guide/creating-a-new-project.html}{Cargo Book --- Creating a New Package}
\item
  \href{https://doc.rust-lang.org/cargo/reference/profiles.html}{Cargo Book --- Profiles}
\end{itemize}

\subsubsection{Sürüm Notları}\label{sürüm-notları}

\begin{itemize}
\tightlist
\item
  \href{https://blog.rust-lang.org/2026/05/28/Rust-1.96.0/}{Rust Blog --- Announcing 1.96.0}
\item
  \href{https://doc.rust-lang.org/stable/releases.html\#version-1960-2026-05-28}{Rust 1.96.0 Release Notes}
\item
  \href{https://blog.rust-lang.org/2026/05/28/Rust-1.96.0/}{Rust Blog --- LLD on 1.96.0 stable}
\end{itemize}

\subsubsection{Hata Kodları}\label{hata-kodları}

\begin{itemize}
\tightlist
\item
  \href{https://doc.rust-lang.org/error_codes/E0601.html}{E0601}
\end{itemize}

\subsubsection{Ek Kaynaklar}\label{ek-kaynaklar}

\begin{itemize}
\tightlist
\item
  \href{https://doc.rust-lang.org/rust-by-example/cargo.html}{Rust By Example --- Cargo}
\end{itemize}

\emph{Not: Sürüm notu kaynakları bu oturumda aramayla doğrulandı. Rust Book, Cargo Book ve E0601 bağlantıları standart doc.rust-lang.org adresleridir. Kodları ve komutları kendi makinende bir kez çalıştırarak kontrol et.}

\begin{center}\rule{0.5\linewidth}{0.5pt}\end{center}

\section{📘 KONU 2: println!() ve Formatlama}\label{konu-2-println-ve-formatlama}

\textbf{Kategori:} Temel \textbar{} \textbf{Sürüm:} Rust 1.96.0

\begin{center}\rule{0.5\linewidth}{0.5pt}\end{center}

\subsection{🟢 KATMAN 1 --- ILK OKUMA}\label{katman-1-—-ilk-okuma-1}

\begin{Shaded}
\begin{Highlighting}[]
\KeywordTok{fn} \FunctionTok{main}\OperatorTok{()} \OperatorTok{\{}
    \KeywordTok{let}\NormalTok{ name }\OperatorTok{=} \StringTok{"Ferris"}\OperatorTok{;}
    \KeywordTok{let}\NormalTok{ age }\OperatorTok{=} \DecValTok{7}\OperatorTok{;}
    \KeywordTok{let}\NormalTok{ pi }\OperatorTok{=} \FloatTok{3.14159}\OperatorTok{;}

\NormalTok{    println}\OperatorTok{!(}\StringTok{"Hello, \{\}!"}\OperatorTok{,}\NormalTok{ name}\OperatorTok{);}        \CommentTok{// positional: sırayla doldurur}
\NormalTok{    println}\OperatorTok{!(}\StringTok{"\{name\} is \{age\} years old"}\OperatorTok{);} \CommentTok{// inline: değişken adını direkt yaz}
\NormalTok{    println}\OperatorTok{!(}\StringTok{"pi = \{pi:.2\}"}\OperatorTok{);}            \CommentTok{// virgülden sonra 2 basamak}
\NormalTok{    println}\OperatorTok{!(}\StringTok{"[\{:\textgreater{}6\}]"}\OperatorTok{,}\NormalTok{ age}\OperatorTok{);}            \CommentTok{// sağa hizala}
\NormalTok{    println}\OperatorTok{!(}\StringTok{"[\{:\textless{}6\}]"}\OperatorTok{,}\NormalTok{ age}\OperatorTok{);}            \CommentTok{// sola hizala}
\NormalTok{    println}\OperatorTok{!(}\StringTok{"[\{:\^{}6\}]"}\OperatorTok{,}\NormalTok{ age}\OperatorTok{);}            \CommentTok{// ortala}
\NormalTok{    println}\OperatorTok{!(}\StringTok{"[\{:06\}]"}\OperatorTok{,}\NormalTok{ age}\OperatorTok{);}            \CommentTok{// sıfırla doldur}
\NormalTok{    println}\OperatorTok{!(}\StringTok{"\{:?\}"}\OperatorTok{,}\NormalTok{ name}\OperatorTok{);}              \CommentTok{// Debug formatı}
\NormalTok{    println}\OperatorTok{!(}\StringTok{"\{\{\}\} are braces"}\OperatorTok{);}         \CommentTok{// süslü parantezi yazdırma}
\OperatorTok{\}}
\CommentTok{// Çıktı:}
\CommentTok{// Hello, Ferris!}
\CommentTok{// Ferris is 7 years old}
\CommentTok{// pi = 3.14}
\CommentTok{// [     7]}
\CommentTok{// [7     ]}
\CommentTok{// [  7   ]}
\CommentTok{// [000007]}
\CommentTok{// "Ferris"}
\CommentTok{// \{\} are braces}
\end{Highlighting}
\end{Shaded}

\texttt{println!} bir macro (makro) olduğu için sonunda \texttt{!} var ve ekrana yazıp satır sonu ekler. Metnin içindeki \texttt{\{\}} işaretleri placeholder (yer tutucu) görevi görür ve argümanlarla sırayla doldurulur. \texttt{\{name\}} yazarsan aynı isimli değişken otomatik alınır. \texttt{:} işaretinden sonra width (genişlik), alignment (hizalama) ve precision (hassasiyet) gibi ayarlar yazılır. \texttt{\{:?\}} Debug biçimidir ve string'i tırnaklı gösterir.

\textbf{1.96.0 notu:} Formatlama kuralları aynı; 1.96.0'ın öne çıkan değişiklikleri (LLD, workspace publish) bu konuyu etkilemiyor.

\begin{center}\rule{0.5\linewidth}{0.5pt}\end{center}

\subsection{🔵 KATMAN 2 --- RECALL KARTLARI}\label{katman-2-—-recall-kartlari-1}

\begin{itemize}
\item
  \texttt{println!("\{\}",\ x)} → Display biçimi, sonuna \texttt{\textbackslash{}n} ekler
  İlgili: \texttt{print!}
\item
  \texttt{print!("no\ newline")} → satır sonu eklemez
  Hata: çıktı geç görünür → stdout tamponlanır, \texttt{flush} gerekir (aşağıya bak)
\item
  \texttt{eprintln!("error")} → stderr'e yazar
  İlgili: \texttt{eprint!}
\item
  \texttt{let\ s\ =\ format!("\{a\}-\{b\}");} → ekrana yazmaz, \texttt{String} döndürür
  İlgili: heap allocation
\item
  \texttt{\{0\}\ \{1\}\ \{0\}} → numaralı argüman; \texttt{\{n\}\ \{m\}} → isimli argüman

\begin{Shaded}
\begin{Highlighting}[]
\NormalTok{println}\OperatorTok{!(}\StringTok{"\{0\} \{1\} \{0\}"}\OperatorTok{,} \StringTok{"a"}\OperatorTok{,} \StringTok{"b"}\OperatorTok{);}        \CommentTok{// a b a}
\NormalTok{println}\OperatorTok{!(}\StringTok{"\{x\} \{y\}"}\OperatorTok{,}\NormalTok{ x }\OperatorTok{=} \DecValTok{1}\OperatorTok{,}\NormalTok{ y }\OperatorTok{=} \DecValTok{2}\OperatorTok{);}        \CommentTok{// 1 2}
\end{Highlighting}
\end{Shaded}
\item
  Hizalama: \texttt{\{:\textless{}8\}} sol, \texttt{\{:\textgreater{}8\}} sağ, \texttt{\{:\^{}8\}} orta, \texttt{\{:*\^{}9\}} yıldızla doldurarak orta
\item
  Sayı ayarları: \texttt{\{:+\}} işaret, \texttt{\{:05\}} sıfır dolgu, \texttt{\{:.3\}} 3 basamak, \texttt{\{:8.2\}} genişlik 8 + 2 basamak
\item
  Taban ve bilimsel: \texttt{\{:b\}} binary, \texttt{\{:x\}} hex, \texttt{\{:\#x\}} \texttt{0x} önekli hex, \texttt{\{:o\}} octal, \texttt{\{:e\}} bilimsel
\item
  Dinamik genişlik: \texttt{println!("\{:\textgreater{}w\$\}",\ x,\ w\ =\ 8);} veya \texttt{println!("\{x:\textgreater{}w\$.p\$\}");}
\item
  Debug: \texttt{\{:?\}} tek satır, \texttt{\{:\#?\}} pretty-print (okunaklı yazdırma)

\begin{Shaded}
\begin{Highlighting}[]
\NormalTok{\#}\OperatorTok{[}\FunctionTok{derive}\OperatorTok{(}\NormalTok{Debug}\OperatorTok{)]}
\KeywordTok{struct}\NormalTok{ Point }\OperatorTok{\{}\NormalTok{ x}\OperatorTok{:} \DataTypeTok{i32}\OperatorTok{,}\NormalTok{ y}\OperatorTok{:} \DataTypeTok{i32} \OperatorTok{\}}
\CommentTok{// \{:?\}  → Point \{ x: 1, y: 2 \}}
\end{Highlighting}
\end{Shaded}


  Hata: \texttt{E0277}: \texttt{Point} doesn't implement \texttt{std::fmt::Display} (\texttt{\{\}} kullandın) veya \texttt{Debug} (derive unuttun)
\item
  \texttt{dbg!(x)} → stderr'e dosya:satır ile yazar ve değeri geri döndürür
\item
  \texttt{flush} için: \texttt{use\ std::io::Write;\ std::io::stdout().flush().unwrap();}
  Hata: \texttt{E0599}: no method named \texttt{flush} found → \texttt{use\ std::io::Write;} eksik
\item
  Hata: \texttt{\{foo\}} yazdın ama \texttt{foo} yok → \texttt{E0425}: cannot find value \texttt{foo} in this scope
\end{itemize}

\textbf{GoT bağlantı haritası:}

\begin{verbatim}
println! ──┬── {}        ──► Display
           ├── {:?}      ──► Debug  ──► #[derive(Debug)]
           │     └─ {:#?} ──► pretty-print
           ├── {name}    ──► inline capture
           └── {:>8.2}   ──► width / precision
print!    ──► flush gerekebilir
eprintln! ──► stderr
format!   ──► String (heap)
\end{verbatim}

\begin{center}\rule{0.5\linewidth}{0.5pt}\end{center}

\subsection{🟣 KATMAN 3 --- DERIN TEKNIK NOT}\label{katman-3-—-derin-teknik-not-1}

\textbf{Adım 1 --- Bellek Seviyesi}

\begin{itemize}
\tightlist
\item
  Önce soru: ``Bellekte ne oluyor?'' Cevap: \texttt{println!} argümanları \textbf{taşımaz}, referans alır.

\begin{Shaded}
\begin{Highlighting}[]
\KeywordTok{let}\NormalTok{ s }\OperatorTok{=} \DataTypeTok{String}\OperatorTok{::}\FunctionTok{from}\OperatorTok{(}\StringTok{"hi"}\OperatorTok{);}
\NormalTok{println}\OperatorTok{!(}\StringTok{"\{s\}"}\OperatorTok{);}
\NormalTok{println}\OperatorTok{!(}\StringTok{"\{s\}"}\OperatorTok{);} \CommentTok{// OK: s hâlâ geçerli (move olmadı)}
\end{Highlighting}
\end{Shaded}
\item
  \texttt{format\_args!} stack'te küçük bir \texttt{fmt::Arguments} yapısı kurar. İçinde sabit metin parçalarına referanslar ve her argüman için (pointer + formatlayıcı fonksiyon pointer'ı) çifti vardır. Bu aşamada heap kullanılmaz.
\item
  Çıktı stdout'a gider. Stdout'un arkasında küçük bir buffer (tampon) vardır ve \texttt{\textbackslash{}n} görünce flush edilir. \texttt{print!} satır sonu içermediği için çıktı hemen görünmeyebilir.
\item
  \texttt{format!} ise sonucu heap'te büyüyen bir \texttt{String}`e yazar. Stack'te (ptr, len, cap), heap'te metin tutulur.
\end{itemize}

\textbf{Adım 2 --- Compiler Davranışı}

\begin{itemize}
\tightlist
\item
  Önce soru: ``Compiler ne kontrol ediyor?'' Cevap: format string \textbf{compile time'da} parse edilir.

  \begin{itemize}
  \tightlist
  \item
    Placeholder sayısı ile argüman sayısı eşleşmezse hata verir.
  \item
    Kullanılmayan argüman hata olur.
  \item
    Her placeholder için ilgili trait kontrol edilir (\texttt{\{\}} → \texttt{Display}, \texttt{\{:?\}} → \texttt{Debug}, \texttt{\{:x\}} → \texttt{LowerHex}).
  \end{itemize}
\item
  Format string \textbf{string literal} olmalıdır; \texttt{println!(s)} derlenmez.
\item
  Inline capture sadece \textbf{identifier} alır; \texttt{\{x.y\}} ya da \texttt{\{a\ +\ b\}} yazamazsın. Önce değişkene ata.
\item
  Macro'nun açılmış hâlini görmek için nightly gerekir: \texttt{cargo\ +nightly\ rustc\ --\ -Zunpretty=expanded}. MIR için: \texttt{cargo\ rustc\ --\ --emit=mir}.
\end{itemize}

\textbf{Adım 3 --- Hata Kodları ve Mesajlar}

\begin{itemize}
\tightlist
\item
  \textbf{E0277:} \texttt{\{\}} ile yazdırdığın tipte \texttt{Display} yok. Çözüm: \texttt{\{:?\}} + \texttt{\#{[}derive(Debug){]}} kullan veya \texttt{impl\ fmt::Display} yaz. Detay: \texttt{rustc\ --explain\ E0277}.
\item
  \textbf{E0599:} \texttt{flush} çağırıyorsun ama \texttt{Write} trait'i kapsamda değil. Çözüm: \texttt{use\ std::io::Write;}.
\item
  \textbf{E0425:} \texttt{\{foo\}} içindeki \texttt{foo} tanımlı değil. Çözüm: değişkeni tanımla veya \texttt{foo\ =\ ...} ile isimli argüman ver.
\item
  \textbf{Kodsuz hatalar:}

  \begin{itemize}
  \tightlist
  \item
    \texttt{2\ positional\ arguments\ in\ format\ string,\ but\ there\ is\ 1\ argument}
  \item
    \texttt{argument\ never\ used}
  \item
    \texttt{format\ argument\ must\ be\ a\ string\ literal}
  \item
    \texttt{invalid\ format\ string:\ expected\ '\}'} (süslü parantezi \texttt{\{\{\ \}\}} ile kaçırmadın)
  \end{itemize}
\end{itemize}

\textbf{Adım 4 --- ``Neden Böyle?''}

\begin{itemize}
\tightlist
\item
  Çözülen problem: C'deki \texttt{printf("\%d",\ x)} tip güvensizliği. \texttt{\%d} yazıp \texttt{char*} verirsen derleyici çoğu zaman susar ve davranış tanımsız olur (undefined behavior).
\item
  Rust'ta \texttt{println!("\{\}",\ x)} tip kontrolünü compile time'da yapar; runtime'da format string yorumlanmaz.
\item
  Trade-off: macro olduğu için format string çalışma zamanında değiştirilemez. Dinamik biçim gerekiyorsa \texttt{\{:w\$\}} gibi parametreleri kullan.
\end{itemize}

\textbf{Adım 5 --- Rust 1.96.0 Farkı}

\begin{itemize}
\tightlist
\item
  Rust 1.96.0'daki değişiklikler (yeni Range* tipleri, \texttt{assert\_matches!} makrosu, Wasm linker güncellemeleri) formatlama macro'larını etkilemiyor.
\item
  Format args capture, Rust 1.58 duyurusunda ``captured identifiers in format strings'' başlığıyla yer alır. Bu yüzden \texttt{\{name\}} güvenle kullanılabilir.
\item
  Width ve precision parametreleri de capture edilebilir (\texttt{\{x:width\$.precision\$\}}); aynı isimli açık bir named argument varsa o önceliklidir.
\end{itemize}

\begin{center}\rule{0.5\linewidth}{0.5pt}\end{center}

\subsection{🔴 KATMAN 4 --- SIK UNUTULANLAR + DERINLEŞTIRME}\label{katman-4-—-sik-unutulanlar-derinleştirme-1}

\textbf{Sık Unutulan 1: \texttt{\{\}} mi \texttt{\{:?\}} mi?}

\begin{itemize}
\tightlist
\item
  Neden unutuluyor: ikisi de ``yazdır'' gibi görünüyor.
\item
  Yol 1 (görsel): \texttt{?} = soru işareti = ``içini göster, debug için''.
\item
  Yol 2 (kod): \texttt{println!("\{:?\}",\ vec!{[}1,\ 2{]});} → \texttt{{[}1,\ 2{]}}, ama \texttt{println!("\{\}",\ vec!{[}1,\ 2{]});} derlenmez.
\item
  Yol 3 (hata): \texttt{E0277:\ Vec\textless{}i32\textgreater{}\ doesn't\ implement\ std::fmt::Display}
\end{itemize}

\textbf{Sık Unutulan 2: Süslü parantezi yazdırmak}

\begin{itemize}
\tightlist
\item
  Neden unutuluyor: \texttt{\{} normalde placeholder açar.
\item
  Yol 1 (analoji): iki kapı kilidi, \texttt{\{\{} ve \texttt{\}\}} ikişer kez yazılır.
\item
  Yol 2 (kod): \texttt{println!("\{\{\}\}");} → \texttt{\{\}}
\item
  Yol 3 (hata): \texttt{println!("\{");} → \texttt{invalid\ format\ string:\ expected\ '\}'\ but\ string\ was\ terminated}
\end{itemize}

\textbf{Sık Unutulan 3: Sözdizimi sırası}

\begin{itemize}
\tightlist
\item
  Neden unutuluyor: \texttt{fill\ align\ sign\ \#\ 0\ width\ .precision\ type} sırası uzun.
\item
  Yol 1 (görsel): önce \textbf{nasıl dolsun} (fill+align), sonra \textbf{ne kadar geniş} (width), sonra \textbf{kaç basamak} (.precision), en sonda \textbf{hangi tür} (x, b, e, ?).
\item
  Yol 2 (kod): \texttt{\{:*\^{}10.2\}} → fill \texttt{*}, orta, genişlik 10, 2 basamak.
\item
  Yol 3 (hata): \texttt{\{:.2\textgreater{}8\}} yanlış sıra → \texttt{invalid\ format\ string}
\end{itemize}

\textbf{Hazır Derinleştirme Prompt'u:}

\begin{verbatim}
Rust 1.96.0'da println!() ve formatlama konusunu en detaylı şekilde anlat. Şunları kapsa:

1. BELLEK SEVİYESİ: println! ve format! çağrısında stack'te ne var, heap'te ne var?
   fmt::Arguments yapısı nedir? Argümanlar neden move edilmez?

2. COMPILER DAVRANIŞI: Format string compile time'da nasıl parse ediliyor?
   Display, Debug, LowerHex trait'leri nasıl seçiliyor? Macro açılımı nasıl görünür?
   (-Zunpretty=expanded, --emit=mir)

3. HATA KODLARI: E0277, E0599, E0425 ve format string hata mesajları.
   Her biri için: ne zaman, neden, nasıl çözülür.
   `rustc --explain EXXXX` çıktısını da ekle.

4. C/C++ KARŞILAŞTIRMASI: printf ve std::cout ile aynı işlem nasıl yapılır?
   Rust neden farklı? Hangi problemi çözüyor?

5. ANALOJİ: Günlük hayattan bir benzetme yap (örn: şablon mektup, doldurulacak form).

6. RUST 1.96.0: Bu konuda son sürümde değişen bir şey var mı?

7. PRATİK: 3 tane kod örneği ver — kolay, orta, zor.

Açıklama Türkçe, teknik terimler İngilizce, kod yorumları Türkçe olsun.
Her bölümü ayrı başlıkla ver. Kod bloklarını çalıştırılabilir yaz.
\end{verbatim}

\begin{center}\rule{0.5\linewidth}{0.5pt}\end{center}

\subsection{📝 GÖREV (Boş dosyaya yaz)}\label{görev-boş-dosyaya-yaz}

\textbf{🟢 Kolay}

\begin{enumerate}
\def\labelenumi{\arabic{enumi}.}
\tightlist
\item
  \texttt{name} ve \texttt{age} değişkenlerini tanımla; \texttt{println!} ile inline capture kullanarak ``Ali is 20 years old'' yazdır.
\item
  \texttt{pi\ =\ 3.14159} değerini 2 basamakla yazdır.
\end{enumerate}

\textbf{🟡 Orta}
3. \texttt{\#{[}derive(Debug){]}\ struct\ Point\ \{\ x:\ i32,\ y:\ i32\ \}} tanımla; \texttt{\{:?\}} ve \texttt{\{:\#?\}} ile yazdır.
4. \texttt{255} sayısını decimal, binary (\texttt{\{:b\}}), hex (\texttt{\{:\#x\}}) ve octal olarak tek tek yazdır.

\textbf{🔴 Zor}
5. \texttt{\{\{} \texttt{\}\}} kullanarak \texttt{\{value\}} metnini harfi harfine yazdır. Sonra \texttt{w\ =\ 10} değişkeniyle dinamik genişlikte (\texttt{\{:\textgreater{}w\$\}}) sağa hizalı bir sayı yazdır ve \texttt{print!} + \texttt{flush} ile aynı satırda ``Loading\ldots{}'' yazıp devam et.

\textbf{Başarı kriteri:} Hepsini hatasız yazarsan → öğrendin ✅

\subsubsection{✅ Kendini Test Et (2-3 saat sonra, sıfırdan)}\label{kendini-test-et-2-3-saat-sonra-sıfırdan-1}

\begin{enumerate}
\def\labelenumi{\arabic{enumi}.}
\tightlist
\item
  \texttt{\{\}} ve \texttt{\{:?\}} farkını gösteren iki satır yaz.
\item
  \texttt{\{:*\^{}9.2\}} ile bir \texttt{f64} yazdır.
\item
  \texttt{format!} ile bir \texttt{String} üret ve yazdır.
\end{enumerate}

\begin{center}\rule{0.5\linewidth}{0.5pt}\end{center}

\subsection{📚 KAYNAKÇA}\label{kaynakça-1}

\subsubsection{Resmi Kaynaklar}\label{resmi-kaynaklar-1}

\begin{itemize}
\tightlist
\item
  \href{https://doc.rust-lang.org/book/ch01-02-hello-world.html}{Rust Book --- Hello, World!}
\item
  \href{https://doc.rust-lang.org/std/fmt/index.html}{std::fmt}
\item
  \href{https://doc.rust-lang.org/std/macro.println.html}{std --- println!}
\end{itemize}

\subsubsection{Sürüm Notları}\label{sürüm-notları-1}

\begin{itemize}
\tightlist
\item
  \href{https://blog.rust-lang.org/2022/01/13/Rust-1.58.0.html}{Rust 1.58.0 --- Captured identifiers in format strings}
\item
  \href{https://blog.rust-lang.org/2026/05/28/Rust-1.96.0/}{Rust Blog --- Announcing 1.96.0}
\end{itemize}

\subsubsection{Hata Kodları}\label{hata-kodları-1}

\begin{itemize}
\tightlist
\item
  \href{https://doc.rust-lang.org/error_codes/E0277.html}{E0277}
\item
  \href{https://doc.rust-lang.org/error_codes/E0599.html}{E0599}
\item
  \href{https://doc.rust-lang.org/error_codes/E0425.html}{E0425}
\end{itemize}

\subsubsection{Ek Kaynaklar}\label{ek-kaynaklar-1}

\begin{itemize}
\tightlist
\item
  \href{https://doc.rust-lang.org/rust-by-example/hello/print.html}{Rust By Example --- Formatted print}
\item
  \href{https://rust.googlesource.com/rust/+/HEAD/tests/ui/fmt/format-args-capture.rs}{rustc testleri --- format-args-capture}
\end{itemize}

\emph{Not: Inline capture ve width/precision capture bilgileri aramada doğrulandı. Kitap, std ve hata kodu bağlantıları standart doc.rust-lang.org adresleridir. Çıktı satırlarını ve hata mesajlarını kendi makinende bir kez çalıştırarak kontrol et.}

\begin{center}\rule{0.5\linewidth}{0.5pt}\end{center}

\section{📘 KONU 3: Variables (let, mut)}\label{konu-3-variables-let-mut}

\textbf{Kategori:} Temel \textbar{} \textbf{Sürüm:} Rust 1.96.0

\begin{center}\rule{0.5\linewidth}{0.5pt}\end{center}

\subsection{🟢 KATMAN 1 --- ILK OKUMA}\label{katman-1-—-ilk-okuma-2}

\begin{Shaded}
\begin{Highlighting}[]
\KeywordTok{const}\NormalTok{ MAX\_POINTS}\OperatorTok{:} \DataTypeTok{u32} \OperatorTok{=} \DecValTok{100\_}\DecValTok{000}\OperatorTok{;}      \CommentTok{// const: tür yazmak zorunlu}

\KeywordTok{fn} \FunctionTok{main}\OperatorTok{()} \OperatorTok{\{}
    \KeywordTok{let}\NormalTok{ x }\OperatorTok{=} \DecValTok{5}\OperatorTok{;}                        \CommentTok{// immutable}
    \KeywordTok{let} \KeywordTok{mut}\NormalTok{ y }\OperatorTok{=} \DecValTok{10}\OperatorTok{;}                   \CommentTok{// mutable}
\NormalTok{    y }\OperatorTok{+=} \DecValTok{1}\OperatorTok{;}                           \CommentTok{// OK: y değişebilir}
    \KeywordTok{let}\NormalTok{ x }\OperatorTok{=}\NormalTok{ x }\OperatorTok{+} \DecValTok{1}\OperatorTok{;}                    \CommentTok{// shadowing: yeni bir x (6)}

    \OperatorTok{\{}
        \KeywordTok{let}\NormalTok{ x }\OperatorTok{=}\NormalTok{ x }\OperatorTok{*} \DecValTok{2}\OperatorTok{;}                \CommentTok{// iç blokta başka bir x (12)}
\NormalTok{        println}\OperatorTok{!(}\StringTok{"inner x = \{x\}"}\OperatorTok{);}
    \OperatorTok{\}}                                 \CommentTok{// iç x burada biter}

    \KeywordTok{let}\NormalTok{ spaces }\OperatorTok{=} \StringTok{"   "}\OperatorTok{;}               \CommentTok{// \&str}
    \KeywordTok{let}\NormalTok{ spaces }\OperatorTok{=}\NormalTok{ spaces}\OperatorTok{.}\FunctionTok{len}\OperatorTok{();}        \CommentTok{// shadowing: tür usize oldu}

\NormalTok{    println}\OperatorTok{!(}\StringTok{"x = \{x\}, y = \{y\}, spaces = \{spaces\}, max = \{MAX\_POINTS\}"}\OperatorTok{);}
\OperatorTok{\}}
\CommentTok{// Çıktı:}
\CommentTok{// inner x = 12}
\CommentTok{// x = 6, y = 11, spaces = 3, max = 100000}
\end{Highlighting}
\end{Shaded}

Rust'ta \texttt{let} ile tanımlanan variable (değişken) varsayılan olarak immutable (değiştirilemez) olur. Değiştirmek istiyorsan mutable (değiştirilebilir) yapmak için \texttt{mut} yazman gerekir. Aynı ismi \texttt{let} ile tekrar yazarsan shadowing (gölgeleme) olur: yeni bir variable oluşur ve tür bile değişebilir. \texttt{const} ise constant (sabit) değerdir; türü yazmak zorunludur ve isim büyük harfle yazılır. Süslü parantezle açılan her block (blok) yeni bir scope (kapsam) oluşturur ve içeride yapılan shadowing blok bitince sona erer.

\textbf{1.96.0 notu:} Variables konusunda 1.96.0'da değişiklik yok; kurallar aynı.

\begin{center}\rule{0.5\linewidth}{0.5pt}\end{center}

\subsection{🔵 KATMAN 2 --- RECALL KARTLARI}\label{katman-2-—-recall-kartlari-2}

\begin{itemize}
\tightlist
\item
  \texttt{let\ x\ =\ 5;} → immutable
  Hata: \texttt{E0384}: cannot assign twice to immutable variable \texttt{x} (\texttt{x\ =\ 6;} yazdın)
  İlgili: \texttt{mut}
\item
  \texttt{let\ mut\ y\ =\ 10;\ y\ =\ 11;} → değeri değiştirir, \textbf{tür aynı kalmalı}
  İlgili: \texttt{unused\_mut} uyarısı (değiştirmiyorsan \texttt{mut} gereksiz)
\item
  \texttt{let\ x\ =\ x\ +\ 1;} → shadowing: eski \texttt{x} bir kenara bırakılır, yeni variable oluşur
  İlgili: scope
\item
  Tür değiştirme:

\begin{Shaded}
\begin{Highlighting}[]
\KeywordTok{let}\NormalTok{ s }\OperatorTok{=} \StringTok{"abc"}\OperatorTok{;}
\KeywordTok{let}\NormalTok{ s }\OperatorTok{=}\NormalTok{ s}\OperatorTok{.}\FunctionTok{len}\OperatorTok{();}      \CommentTok{// OK: shadowing, tür değişir}
\KeywordTok{let} \KeywordTok{mut}\NormalTok{ t }\OperatorTok{=} \StringTok{"abc"}\OperatorTok{;}
\NormalTok{t }\OperatorTok{=}\NormalTok{ t}\OperatorTok{.}\FunctionTok{len}\OperatorTok{();}          \CommentTok{// E0308: mismatched types}
\end{Highlighting}
\end{Shaded}
\item
  \texttt{let\ z:\ i64\ =\ 10;} → type annotation (tür belirtme); yazmazsan type inference (tür çıkarımı) devreye girer
\item
  \texttt{const\ SECS:\ u32\ =\ 60\ *\ 60\ *\ 3;} → tür zorunlu, \texttt{UPPER\_SNAKE\_CASE}, \texttt{mut} yok
  Hata: \texttt{const} ile \texttt{mut} birleşmez (\texttt{const\ globals\ cannot\ be\ mutable})
  İlgili: \texttt{static}
\item
  \texttt{let\ n\ =\ 5;\ const\ C:\ i32\ =\ n;} → \texttt{E0435}: attempt to use a non-constant value in a constant
\item
  \texttt{let\ a;\ a\ =\ 1;} → sonradan değer verme (deferred initialization) OK
  Hata: \texttt{E0381}: değer atamadan \texttt{a}`yı kullandın
\item
  \texttt{let\ \_unused\ =\ 5;} → \texttt{\_} öneki ``kullanmıyorum'' der, uyarı gelmez
  Hata/uyarı: \texttt{warning:\ unused\ variable}
\item
  \texttt{let\ (a,\ b)\ =\ (1,\ 2);} → destructuring ile iki variable tanımlar
\end{itemize}

\textbf{GoT bağlantı haritası:}

\begin{verbatim}
let ──┬── mut ───────► aynı tür, aynı variable
      ├── shadowing ─► yeni variable, tür değişebilir
      ├── tür: annotation / inference
      └── deferred init ──► E0381
const ──► tür zorunlu ──► UPPER_SNAKE_CASE ──► compile time değeri
\end{verbatim}

\begin{center}\rule{0.5\linewidth}{0.5pt}\end{center}

\subsection{🟣 KATMAN 3 --- DERIN TEKNIK NOT}\label{katman-3-—-derin-teknik-not-2}

\textbf{Adım 1 --- Bellek Seviyesi}

\begin{itemize}
\tightlist
\item
  Önce soru: ``Bellekte ne oluyor?'' Cevap: yerel variable'lar fonksiyonun stack frame'inde yaşar. \texttt{i32} 4 byte, \texttt{\&str} 16 byte (ptr + len) tutar.
\item
  \texttt{mut} bellek düzenini \textbf{değiştirmez}. Immutability tamamen compile time kavramıdır ve runtime maliyeti yoktur.
\item
  Shadowing yeni bir variable demektir (derleyici için ayrı bir local). Eski değer scope sonuna kadar yaşar. Örneğin bir \texttt{String} shadow edilince hemen drop edilmez, blok bitince edilir.
\item
  \texttt{const} bir bellek adresine sahip değildir; değer kullanıldığı her yere inline edilir. Sabit adres isteyen \texttt{static} kullanır.
\item
  Kontrol için: \texttt{std::mem::size\_of::\textless{}i32\textgreater{}()} → \texttt{4}
\end{itemize}

\textbf{Adım 2 --- Compiler Davranışı}

\begin{itemize}
\tightlist
\item
  Önce soru: ``Hangi kontrol yapılıyor?'' Cevap: \texttt{E0384} borrow checker aşamasında yakalanır; \texttt{E0308} type checker aşamasında.
\item
  Lint'ler: \texttt{unused\_variables} ve \texttt{unused\_mut} uyarı verir, hata değil.
\item
  Tür çıkarımı: \texttt{let\ x\ =\ 5;} → varsayılan \texttt{i32}. Sonradan \texttt{let\ y:\ u8\ =\ x;} yazarsan çıkarım \texttt{u8} olur.
\item
  MIR'de her shadow edilen variable ayrı local olarak (\texttt{\_1}, \texttt{\_3} gibi) görünür. Görmek için: \texttt{cargo\ rustc\ --\ --emit=mir}.
\end{itemize}

\textbf{Adım 3 --- Hata Kodları}

\begin{itemize}
\tightlist
\item
  \textbf{E0384:} immutable variable'a ikinci kez değer atadın. Çözüm: \texttt{let\ mut} yap veya shadowing kullan. Rust Book bu hatayı ilk örnek olarak gösterir. Detay: \texttt{rustc\ --explain\ E0384}.
\item
  \textbf{E0308:} \texttt{mut} variable'ın türünü değiştirmeye çalıştın. Çözüm: shadowing kullan veya doğru türde değer ata.
\item
  \textbf{E0381:} değer atanmamış variable'ı okudun. Çözüm: tanımla ve ata, ya da her kolda atama yap.
\item
  \textbf{E0435:} \texttt{const} içinde runtime değeri (\texttt{let} variable'ı) kullandın. Çözüm: sadece compile time'da bilinen değer kullan.
\item
  \textbf{Kodsuz:} \texttt{missing\ type\ for\ const\ item} → \texttt{const} için tür yazmadın.
\end{itemize}

\textbf{Adım 4 --- ``Neden Böyle?''}

\begin{itemize}
\tightlist
\item
  Çözülen problem: ``Bu değer bir yerde sessizce değişti'' hataları.
\item
  C/C++`ta \texttt{int\ x\ =\ 5;} değişkendir, \texttt{const\ int} ise sonradan istenir. Rust bunu tersine çevirir: varsayılan sabit, değişim \texttt{mut} ile açıkça belirtilir.
\item
  Immutable varsayılan olunca kod okurken ``bu değer değişmez'' garantisi alırsın; eşzamanlılıkta (concurrency) da paylaşım güvenli olur.
\item
  Seçim: değeri aynı türde yerinde güncelleyeceksen \texttt{mut}, dönüştürüp yeni isimle devam edeceksen shadowing.
\end{itemize}

\textbf{Adım 5 --- Rust 1.96.0 Farkı}

\begin{itemize}
\tightlist
\item
  Rust 1.96.0'daki değişiklikler (yeni Range* tipleri, \texttt{assert\_matches!} makrosu, Wasm linker güncellemeleri) variables, shadowing ve \texttt{const} kurallarını etkilemiyor.
\item
  LLD linker ve \texttt{cargo\ publish\ --workspace} değişiklikleri bu konuyu etkilemez.
\end{itemize}

\begin{center}\rule{0.5\linewidth}{0.5pt}\end{center}

\subsection{🔴 KATMAN 4 --- SIK UNUTULANLAR + DERINLEŞTIRME}\label{katman-4-—-sik-unutulanlar-derinleştirme-2}

\textbf{Sık Unutulan 1: Shadowing ile \texttt{mut} farkı}

\begin{itemize}
\tightlist
\item
  Neden unutuluyor: ikisi de ``değer değişti'' gibi görünüyor.
\item
  Yol 1 (analoji): \texttt{mut} = aynı kutunun içini değiştir; shadowing = yeni kutu al, eskisini kenara koy.
\item
  Yol 2 (kod): \texttt{let\ s\ =\ "a";\ let\ s\ =\ s.len();} OK; \texttt{let\ mut\ s\ =\ "a";\ s\ =\ s.len();} hata.
\item
  Yol 3 (hata): \texttt{E0308:\ mismatched\ types}
\end{itemize}

\textbf{Sık Unutulan 2: \texttt{const} için tür yazmak}

\begin{itemize}
\tightlist
\item
  Neden unutuluyor: \texttt{let}`te tür yazmak zorunlu değil.
\item
  Yol 1 (görsel): \texttt{const} = tabelaya yazılan sabit, tür de tabelada olmalı.
\item
  Yol 2 (kod): \texttt{const\ MAX:\ u32\ =\ 100;}
\item
  Yol 3 (hata): \texttt{missing\ type\ for\ const\ item}
\end{itemize}

\textbf{Sık Unutulan 3: Shadowing scope sonunda biter}

\begin{itemize}
\tightlist
\item
  Neden unutuluyor: iç blokta \texttt{x} değişmiş gibi görünür.
\item
  Yol 1 (analoji): iç blok geçici bir maske; çıkınca yüz eski hâline döner.
\item
  Yol 2 (kod): iç blokta \texttt{let\ x\ =\ 12;}, dışarıda \texttt{x} hâlâ \texttt{6}.
\item
  Yol 3 (hata): iç bloktaki \texttt{y} dışarıda kullanılırsa \texttt{E0425}: cannot find value \texttt{y} in this scope
\end{itemize}

\textbf{Hazır Derinleştirme Prompt'u:}

\begin{verbatim}
Rust 1.96.0'da Variables (let, mut, shadowing, const) konusunu en detaylı şekilde anlat. Şunları kapsa:

1. BELLEK SEVİYESİ: Stack'te ne var? Kaç byte? mut bellek düzenini değiştiriyor mu?
   Shadowing ile eski değere ne oluyor? const ile static farkı nedir?

2. COMPILER DAVRANIŞI: Borrow checker E0384'ü nasıl buluyor? Tür çıkarımı nasıl
   çalışıyor? unused_variables ve unused_mut lint'leri. MIR çıktısında
   shadow edilen variable'lar nasıl görünür? (--emit=mir)

3. HATA KODLARI: E0384, E0308, E0381, E0435. Her biri için: ne zaman, neden,
   nasıl çözülür. `rustc --explain EXXXX` çıktısını da ekle.

4. C/C++ KARŞILAŞTIRMASI: int x, const int, constexpr ile karşılaştır.
   Rust varsayılanı neden immutable yapıyor?

5. ANALOJİ: Günlük hayattan bir benzetme yap (örn: kalemle yazı / kurşun kalem,
   kutular, etiketler).

6. RUST 1.96.0: Bu konuda son sürümde değişen bir şey var mı?

7. PRATİK: 3 tane kod örneği ver — kolay, orta, zor.

Açıklama Türkçe, teknik terimler İngilizce, kod yorumları Türkçe olsun.
Her bölümü ayrı başlıkla ver. Kod bloklarını çalıştırılabilir yaz.
\end{verbatim}

\begin{center}\rule{0.5\linewidth}{0.5pt}\end{center}

\subsection{📝 GÖREV (Boş dosyaya yaz)}\label{görev-boş-dosyaya-yaz-1}

\textbf{🟢 Kolay}

\begin{enumerate}
\def\labelenumi{\arabic{enumi}.}
\tightlist
\item
  \texttt{let\ x\ =\ 5;} yaz ve yazdır. Sonra \texttt{let\ mut} yapıp \texttt{x}`i \texttt{10} yap ve yazdır.
\item
  \texttt{const\ MAX\_USERS:\ u32\ =\ 50;} tanımla ve yazdır.
\end{enumerate}

\textbf{🟡 Orta}
3. \texttt{let\ x\ =\ 5;} ile başla; shadowing ile \texttt{x}`i \texttt{x\ +\ 1} yap. Bir iç blokta \texttt{x\ *\ 2} olarak tekrar shadow et. Üç değeri de yazdır.
4. \texttt{let\ spaces\ =\ "\ \ \ ";} sonra shadowing ile \texttt{spaces.len()} değerine çevir ve yazdır.

\textbf{🔴 Zor}
5. \texttt{let\ a;\ } ile tanımla, \texttt{if} ile iki farklı kolda \texttt{a}`ya değer ata, sonra yazdır. Sonra \texttt{a}`yı değer atamadan yazdırmayı dene ve \texttt{E0381}`i gör. Son olarak \texttt{let\ (p,\ q)\ =\ (1,\ 2);} ile destructuring yap.

\textbf{Başarı kriteri:} Hepsini hatasız yazarsan → öğrendin ✅

\subsubsection{✅ Kendini Test Et (2-3 saat sonra, sıfırdan)}\label{kendini-test-et-2-3-saat-sonra-sıfırdan-2}

\begin{enumerate}
\def\labelenumi{\arabic{enumi}.}
\tightlist
\item
  \texttt{let}, \texttt{let\ mut} ve shadowing ile birer örnek yaz.
\item
  \texttt{const} tanımla (tür dahil).
\item
  \texttt{E0384} ve \texttt{E0308} hatalarını bilerek üret ve nedenini söyle.
\end{enumerate}

\begin{center}\rule{0.5\linewidth}{0.5pt}\end{center}

\subsection{📚 KAYNAKÇA}\label{kaynakça-2}

\subsubsection{Resmi Kaynaklar}\label{resmi-kaynaklar-2}

\begin{itemize}
\tightlist
\item
  \href{https://doc.rust-lang.org/book/ch03-01-variables-and-mutability.html}{Rust Book --- Variables and Mutability}
\item
  \href{https://doc.rust-lang.org/reference/items/constant-items.html}{Rust Reference --- Constant items}
\item
  \href{https://doc.rust-lang.org/std/mem/fn.size_of.html}{std --- mem::size\_of}
\end{itemize}

\subsubsection{Sürüm Notları}\label{sürüm-notları-2}

\begin{itemize}
\tightlist
\item
  \href{https://blog.rust-lang.org/2026/05/28/Rust-1.96.0/}{Rust Blog --- Announcing 1.96.0}
\end{itemize}

\subsubsection{Hata Kodları}\label{hata-kodları-2}

\begin{itemize}
\tightlist
\item
  \href{https://doc.rust-lang.org/error_codes/E0384.html}{E0384}
\item
  \href{https://doc.rust-lang.org/error_codes/E0308.html}{E0308}
\item
  \href{https://doc.rust-lang.org/error_codes/E0381.html}{E0381}
\item
  \href{https://doc.rust-lang.org/error_codes/E0435.html}{E0435}
\end{itemize}

\subsubsection{Ek Kaynaklar}\label{ek-kaynaklar-2}

\begin{itemize}
\tightlist
\item
  \href{https://doc.rust-lang.org/rust-by-example/variable_bindings.html}{Rust By Example --- Variable Bindings}
\item
  \href{https://education.web3.foundation/docs/Rust/section2/variables-mutability}{Web3 Foundation --- Variables \& Mutability}
\end{itemize}

\emph{Not: \texttt{E0384} hata metni, \texttt{const} için tür zorunluluğu ve shadowing davranışı aramada Rust Book üzerinden doğrulandı. Diğer bağlantılar standart doc.rust-lang.org adresleridir. Kodları kendi makinende bir kez çalıştırarak kontrol et.}

\begin{center}\rule{0.5\linewidth}{0.5pt}\end{center}

\section{📘 KONU 4: Data Types (int, float, bool, char)}\label{konu-4-data-types-int-float-bool-char}

\textbf{Kategori:} Temel \textbar{} \textbf{Sürüm:} Rust 1.96.0

\begin{center}\rule{0.5\linewidth}{0.5pt}\end{center}

\subsection{🟢 KATMAN 1 --- ILK OKUMA}\label{katman-1-—-ilk-okuma-3}

\begin{Shaded}
\begin{Highlighting}[]
\KeywordTok{fn} \FunctionTok{main}\OperatorTok{()} \OperatorTok{\{}
    \KeywordTok{let}\NormalTok{ a}\OperatorTok{:} \DataTypeTok{i32} \OperatorTok{=} \OperatorTok{-}\DecValTok{42}\OperatorTok{;}           \CommentTok{// işaretli tam sayı}
    \KeywordTok{let}\NormalTok{ b}\OperatorTok{:} \DataTypeTok{u8} \OperatorTok{=} \DecValTok{255}\OperatorTok{;}            \CommentTok{// işaretsiz, max 255}
    \KeywordTok{let}\NormalTok{ c }\OperatorTok{=} \DecValTok{1\_}\DecValTok{000\_}\DecValTok{000}\OperatorTok{;}          \CommentTok{// tür yazmadık → i32}
    \KeywordTok{let}\NormalTok{ pi}\OperatorTok{:} \DataTypeTok{f64} \OperatorTok{=} \FloatTok{3.14159}\OperatorTok{;}      \CommentTok{// ondalıklı (varsayılan f64)}
    \KeywordTok{let}\NormalTok{ half }\OperatorTok{=} \FloatTok{0.5\_f}\DecValTok{32}\OperatorTok{;}         \CommentTok{// tür eki ile f32}
    \KeywordTok{let}\NormalTok{ ok}\OperatorTok{:} \DataTypeTok{bool} \OperatorTok{=} \ConstantTok{true}\OperatorTok{;}        \CommentTok{// mantıksal değer}
    \KeywordTok{let}\NormalTok{ letter}\OperatorTok{:} \DataTypeTok{char} \OperatorTok{=} \CharTok{'R'}\OperatorTok{;}     \CommentTok{// tek tırnak → char}
    \KeywordTok{let}\NormalTok{ heart }\OperatorTok{=} \CharTok{'❤'}\OperatorTok{;}            \CommentTok{// char Unicode da tutar}

\NormalTok{    println}\OperatorTok{!(}\StringTok{"\{a\} \{b\} \{c\}"}\OperatorTok{);}
\NormalTok{    println}\OperatorTok{!(}\StringTok{"\{pi\} \{half\}"}\OperatorTok{);}
\NormalTok{    println}\OperatorTok{!(}\StringTok{"\{ok\} \{letter\} \{heart\}"}\OperatorTok{);}
\NormalTok{    println}\OperatorTok{!(}\StringTok{"\{\}"}\OperatorTok{,} \DecValTok{7} \OperatorTok{/} \DecValTok{2}\OperatorTok{);}      \CommentTok{// tam sayı bölmesi}
\NormalTok{    println}\OperatorTok{!(}\StringTok{"\{\}"}\OperatorTok{,} \FloatTok{7.0} \OperatorTok{/} \FloatTok{2.0}\OperatorTok{);}  \CommentTok{// ondalıklı bölme}
\NormalTok{    println}\OperatorTok{!(}\StringTok{"\{\}"}\OperatorTok{,} \DecValTok{7} \OperatorTok{\%} \DecValTok{2}\OperatorTok{);}      \CommentTok{// kalan}
\OperatorTok{\}}
\CommentTok{// Çıktı:}
\CommentTok{// -42 255 1000000}
\CommentTok{// 3.14159 0.5}
\CommentTok{// true R ❤}
\CommentTok{// 3}
\CommentTok{// 3.5}
\CommentTok{// 1}
\end{Highlighting}
\end{Shaded}

Rust'ta her değerin bir data type (veri tipi) vardır ve dört scalar (tekil değer) tip temeldir: integer (tam sayı), floating-point (kayan noktalı sayı), bool (mantıksal değer) ve char (karakter). Tam sayıda varsayılan tip \texttt{i32}, ondalıklı sayıda \texttt{f64}`tür. \texttt{u8} gibi unsigned (işaretsiz) tipler sadece pozitif sayı tutar, \texttt{i32} gibi signed (işaretli) tipler negatif de tutar. \texttt{char} tek tırnakla yazılır ve 4 byte yer kaplar; çift tırnak ise string olur. İki tam sayıyı bölersen sonuç da tam sayı çıkar (\texttt{7\ /\ 2} = \texttt{3}).

\textbf{1.96.0 notu:} Temel tipler aynı. Float tipleri için bazı matematik fonksiyonları (\texttt{floor}, \texttt{ceil}, \texttt{round}) bu sürümde const olarak kullanılabilir hale geldi (aşağıda Katman 3, Adım 5).

\begin{center}\rule{0.5\linewidth}{0.5pt}\end{center}

\subsection{🔵 KATMAN 2 --- RECALL KARTLARI}\label{katman-2-—-recall-kartlari-3}

\begin{itemize}
\item
  Tam sayı tipleri:

\begin{verbatim}
i8 i16 i32 i64 i128 isize   (işaretli)
u8 u16 u32 u64 u128 usize   (işaretsiz)
\end{verbatim}


  \texttt{u8} → 0..=255, \texttt{i8} → -128..=127, \texttt{isize}/\texttt{usize} → işlemci genişliği (64-bit'te 8 byte)
  İlgili: \texttt{usize} index ve uzunluk için kullanılır
\item
  Literal (sabit ifade) yazımı:

\begin{Shaded}
\begin{Highlighting}[]
\KeywordTok{let}\NormalTok{ a }\OperatorTok{=} \DecValTok{1\_}\DecValTok{000}\OperatorTok{;}     \CommentTok{// ayırıcı \_}
\KeywordTok{let}\NormalTok{ b }\OperatorTok{=} \BaseNTok{0xff}\OperatorTok{;}      \CommentTok{// hex}
\KeywordTok{let}\NormalTok{ c }\OperatorTok{=} \BaseNTok{0o77}\OperatorTok{;}      \CommentTok{// octal}
\KeywordTok{let}\NormalTok{ d }\OperatorTok{=} \BaseNTok{0b1010}\OperatorTok{;}    \CommentTok{// binary}
\KeywordTok{let}\NormalTok{ e }\OperatorTok{=}\NormalTok{ b}\CharTok{'A'}\OperatorTok{;}      \CommentTok{// u8 (65)}
\KeywordTok{let}\NormalTok{ f }\OperatorTok{=} \DecValTok{10u}\DecValTok{8}\OperatorTok{;}      \CommentTok{// tür eki}
\end{Highlighting}
\end{Shaded}


  Hata: \texttt{let\ x:\ u8\ =\ 256;} → \texttt{error:\ literal\ out\ of\ range\ for\ u8} (kod yok, deny-by-default lint)
\item
  Overflow (taşma):

  \begin{itemize}
  \tightlist
  \item
    debug: panic → \texttt{attempt\ to\ add\ with\ overflow}
  \item
    release: wrapping, yani 255 + 1 → 0
  \item
    Kontrollü yollar:

\begin{Shaded}
\begin{Highlighting}[]
\DecValTok{255u}\DecValTok{8}\OperatorTok{.}\FunctionTok{checked\_add}\OperatorTok{(}\DecValTok{1}\OperatorTok{)}     \CommentTok{// None}
\DecValTok{255u}\DecValTok{8}\OperatorTok{.}\FunctionTok{wrapping\_add}\OperatorTok{(}\DecValTok{1}\OperatorTok{)}    \CommentTok{// 0}
\DecValTok{255u}\DecValTok{8}\OperatorTok{.}\FunctionTok{saturating\_add}\OperatorTok{(}\DecValTok{1}\OperatorTok{)}  \CommentTok{// 255}
\DecValTok{255u}\DecValTok{8}\OperatorTok{.}\FunctionTok{overflowing\_add}\OperatorTok{(}\DecValTok{1}\OperatorTok{)} \CommentTok{// (0, true)}
\end{Highlighting}
\end{Shaded}
  \end{itemize}
\item
  Bölme: \texttt{7\ /\ 2} → \texttt{3}, \texttt{-7\ /\ 2} → \texttt{-3} (sıfıra doğru), \texttt{7\ \%\ 2} → \texttt{1}
  Hata (runtime): sıfıra bölme → panic \texttt{attempt\ to\ divide\ by\ zero}
\item
  Float:

\begin{Shaded}
\begin{Highlighting}[]
\FloatTok{0.1} \OperatorTok{+} \FloatTok{0.2} \OperatorTok{==} \FloatTok{0.3}          \CommentTok{// false}
\KeywordTok{let}\NormalTok{ n }\OperatorTok{=} \DataTypeTok{f64}\OperatorTok{::}\NormalTok{NAN}\OperatorTok{;}
\NormalTok{n}\OperatorTok{.}\FunctionTok{is\_nan}\OperatorTok{()}                \CommentTok{// true}
\end{Highlighting}
\end{Shaded}


  Hata: \texttt{7\ /\ 2.0} → \texttt{E0277}: cannot divide \texttt{\{integer\}} by \texttt{\{float\}}
\item
  Tipleri karıştırma:

\begin{Shaded}
\begin{Highlighting}[]
\KeywordTok{let}\NormalTok{ a}\OperatorTok{:} \DataTypeTok{i32} \OperatorTok{=} \DecValTok{1}\OperatorTok{;}
\KeywordTok{let}\NormalTok{ b}\OperatorTok{:} \DataTypeTok{i64} \OperatorTok{=} \DecValTok{2}\OperatorTok{;}
\KeywordTok{let}\NormalTok{ c }\OperatorTok{=}\NormalTok{ a }\OperatorTok{+}\NormalTok{ b}\OperatorTok{;}   \CommentTok{// E0308: mismatched types}
\end{Highlighting}
\end{Shaded}
\item
  \texttt{as} ile dönüştürme:

\begin{Shaded}
\begin{Highlighting}[]
\FloatTok{3.9\_f}\DecValTok{64} \KeywordTok{as} \DataTypeTok{i32}     \CommentTok{// 3 (kesilir)}
\OperatorTok{-}\DecValTok{1\_i}\DecValTok{32} \KeywordTok{as} \DataTypeTok{u32}      \CommentTok{// 4294967295}
\DecValTok{300\_i}\DecValTok{32} \KeywordTok{as} \DataTypeTok{u8}      \CommentTok{// 44}
\end{Highlighting}
\end{Shaded}
\item
  bool:

\begin{Shaded}
\begin{Highlighting}[]
\KeywordTok{let}\NormalTok{ t}\OperatorTok{:} \DataTypeTok{bool} \OperatorTok{=} \ConstantTok{true}\OperatorTok{;}
\OperatorTok{!}\NormalTok{t }\OperatorTok{\&\&} \ConstantTok{false} \OperatorTok{||} \ConstantTok{true}    \CommentTok{// mantıksal işlemler}
\end{Highlighting}
\end{Shaded}


  Hata: \texttt{if\ 1\ \{\ \}} → \texttt{E0308}: expected \texttt{bool}, found integer
\item
  char:

\begin{Shaded}
\begin{Highlighting}[]
\CharTok{'a'} \KeywordTok{as} \DataTypeTok{u32}          \CommentTok{// 97}
\DecValTok{97u}\DecValTok{8} \KeywordTok{as} \DataTypeTok{char}        \CommentTok{// 'a'}
\end{Highlighting}
\end{Shaded}


  Hata: \texttt{65u32\ as\ char} → \texttt{E0604}: only \texttt{u8} can be cast as \texttt{char}
  Hata: \texttt{let\ c:\ char\ =\ "a";} → \texttt{E0308} (çift tırnak = string)
\item
  \texttt{"42".parse::\textless{}i32\textgreater{}()} → \texttt{Result} döner (tür belirtmeden \texttt{parse()} yazarsan \texttt{E0284}: type annotations needed)
\end{itemize}

\textbf{GoT bağlantı haritası:}

\begin{verbatim}
scalar ──┬── integer ──┬── signed   (i8 ... i128, isize)
         │             ├── unsigned (u8 ... u128, usize)
         │             └── overflow ──► checked_ / wrapping_ / saturating_
         ├── float (f32, f64) ──► NaN
         ├── bool ──► if / && / ||
         └── char (4 byte) ──► as u32 / u8 as char
as ──► tür dönüştürme (kayıplı olabilir)
\end{verbatim}

\begin{center}\rule{0.5\linewidth}{0.5pt}\end{center}

\subsection{🟣 KATMAN 3 --- DERIN TEKNIK NOT}\label{katman-3-—-derin-teknik-not-3}

\textbf{Adım 1 --- Bellek Seviyesi}

\begin{itemize}
\tightlist
\item
  Önce soru: ``Bellekte ne oluyor?'' Cevap: bu tipler stack'te yaşar ve boyutları sabittir.
\end{itemize}

{\def\LTcaptype{none} % do not increment counter
\begin{longtable}[]{@{}ll@{}}
\toprule\noalign{}
Tip & Boyut \\
\midrule\noalign{}
\endhead
\bottomrule\noalign{}
\endlastfoot
\texttt{i8} / \texttt{u8} / \texttt{bool} & 1 byte \\
\texttt{i16} / \texttt{u16} & 2 byte \\
\texttt{i32} / \texttt{u32} / \texttt{f32} / \texttt{char} & 4 byte \\
\texttt{i64} / \texttt{u64} / \texttt{f64} & 8 byte \\
\texttt{i128} / \texttt{u128} & 16 byte \\
\texttt{isize} / \texttt{usize} & pointer genişliği (64-bit'te 8 byte) \\
\end{longtable}
}

\begin{itemize}
\tightlist
\item
  İşaretli sayılar two's complement (ikiye tümleyen) ile tutulur. Örneğin \texttt{-1\_i8} bellekte \texttt{0xFF}`tir. Bu yüzden \texttt{-1\_i32\ as\ u32} sonucu \texttt{4294967295} olur: bit'ler aynı, yorum farklı.
\item
  Float'lar IEEE 754 biçimindedir: \texttt{f64} için 1 bit işaret, 11 bit exponent, 52 bit mantissa. \texttt{0.1}`in ikilik tabanda tam karşılığı olmadığı için \texttt{0.1\ +\ 0.2\ !=\ 0.3} çıkar.
\item
  \texttt{char} bir Unicode scalar value'dur (0..=0x10FFFF, surrogate aralığı hariç) ve sabit 4 byte tutar. Bu yüzden \texttt{'a'} 4 byte, ama \texttt{"a"} string'i UTF-8 olduğu için 1 byte veri tutar.
\item
  \texttt{bool} 1 byte yer kaplar ve sadece \texttt{0} veya \texttt{1} değerini alabilir.
\item
  Kontrol için: \texttt{std::mem::size\_of::\textless{}char\textgreater{}()} → \texttt{4}
\end{itemize}

\textbf{Adım 2 --- Compiler Davranışı}

\begin{itemize}
\tightlist
\item
  Önce soru: ``Tür yazmazsam ne olur?'' Cevap: type inference, kısıt bulamazsa tam sayıya \texttt{i32}, ondalıklıya \texttt{f64} verir.
\item
  \texttt{let\ x:\ u8\ =\ 256;} ve \texttt{let\ y:\ u8\ =\ 255\ +\ 1;} gibi compile time'da bilinen taşmalar derleme hatasıdır (\texttt{overflowing\_literals} ve \texttt{arithmetic\_overflow} lint'leri varsayılan olarak deny).
\item
  Runtime taşması debug'da kontrol edilir. MIR'de toplama \texttt{AddWithOverflow} ve ardından bir \texttt{assert} olarak görünür. Release profilinde bu kontrol yoktur. \texttt{Cargo.toml}`da \texttt{overflow-checks\ =\ true} ile release'e de açabilirsin.
\item
  \texttt{as} dönüşümü hata vermez, bit seviyesinde keser veya yeniden yorumlar. Float'tan tam sayıya \texttt{as} ise saturating'dir: \texttt{300.7\_f64\ as\ u8} → \texttt{255}, \texttt{NaN\ as\ u8} → \texttt{0}.
\item
  MIR görmek için: \texttt{cargo\ rustc\ --\ --emit=mir}.
\end{itemize}

\textbf{Adım 3 --- Hata Kodları}

\begin{itemize}
\tightlist
\item
  \textbf{E0308:} tip uyuşmazlığı (\texttt{i32} + \texttt{i64}, \texttt{if\ 1}, \texttt{char} yerine string). Çözüm: aynı tipe getir (\texttt{b\ as\ i32} veya \texttt{i32::try\_from(b)}).
\item
  \textbf{E0277:} o tip çifti için işlem tanımlı değil (\texttt{\{integer\}\ /\ \{float\}}). Çözüm: iki tarafı da aynı türde yaz (\texttt{7.0\ /\ 2.0}).
\item
  \textbf{E0604:} sadece \texttt{u8}, \texttt{char}`a \texttt{as} ile çevrilebilir. Çözüm: \texttt{char::from\_u32(65)} (\texttt{Option\textless{}char\textgreater{}} döner).
\item
  \textbf{E0284:} \texttt{parse()} hangi tipe çevireceğini bilmiyor. Çözüm: \texttt{let\ n:\ u32\ =\ ...} veya \texttt{parse::\textless{}u32\textgreater{}()}.
\item
  \textbf{Kodsuz:} \texttt{literal\ out\ of\ range\ for\ u8}, \texttt{this\ arithmetic\ operation\ will\ overflow}
\item
  \textbf{Runtime panic:} \texttt{attempt\ to\ add\ with\ overflow}, \texttt{attempt\ to\ divide\ by\ zero}
\item
  Detay için: \texttt{rustc\ --explain\ E0308} (açıklama: beklenen tipin ve bulunan tipin eşleşmemesi).
\end{itemize}

\textbf{Adım 4 --- ``Neden Böyle?''}

\begin{itemize}
\tightlist
\item
  Çözülen problem: C'de \texttt{int} boyutu platforma bağlıdır ve tipler sessizce dönüşür (örneğin signed ile unsigned karşılaştırması). Rust'ta \texttt{i32}, \texttt{u64} gibi tipler sabit genişliklidir ve dönüşüm açıktır.
\item
  C'de signed overflow tanımsız davranıştır (undefined behavior). Rust'ta debug'da panic, release'te wrapping gibi tanımlı bir davranış vardır.
\item
  C'nin \texttt{char}`ı 1 byte'tır; Rust'ın \texttt{char}`ı Unicode scalar value'dur.
\item
  Trade-off: açık dönüşüm yazmak biraz daha uzun, ama hatalar derleme zamanında yakalanır.
\end{itemize}

\textbf{Adım 5 --- Rust 1.96.0 Farkı}

\begin{itemize}
\tightlist
\item
  Rust 1.96.0'daki değişiklikler (yeni Range* tipleri, \texttt{assert\_matches!} makrosu, Wasm linker güncellemeleri) scalar data type kurallarını değiştirmiyor.
\item
  Veri tiplerinin temel kuralları (boyutlar, overflow, \texttt{as}) değişmedi.
\item
  Yeni Range* tipleri scalar tiplerin temel kurallarını değiştirmiyor; \texttt{assert\_matches!} makrosu ve Wasm linker güncellemeleri de bu konuyu etkilemiyor.
\end{itemize}

\begin{center}\rule{0.5\linewidth}{0.5pt}\end{center}

\subsection{🔴 KATMAN 4 --- SIK UNUTULANLAR + DERINLEŞTIRME}\label{katman-4-—-sik-unutulanlar-derinleştirme-3}

\textbf{Sık Unutulan 1: Overflow debug ve release'te farklı davranır}

\begin{itemize}
\tightlist
\item
  Neden unutuluyor: kodu debug'da çalıştırırsın ve release'te farklı sonuç beklemezsin.
\item
  Yol 1 (analoji): kilometre sayacı, 999999'dan sonra 000000'a döner (wrapping). Debug'da ise sayaç alarm verir (panic).
\item
  Yol 2 (kod): \texttt{255u8.checked\_add(1)} → \texttt{None}
\item
  Yol 3 (hata): \texttt{thread\ 'main'\ panicked:\ attempt\ to\ add\ with\ overflow}
\end{itemize}

\textbf{Sık Unutulan 2: Tam sayı ve ondalıklıyı karıştırmak}

\begin{itemize}
\tightlist
\item
  Neden unutuluyor: matematikte \texttt{7\ /\ 2.0} doğal görünür.
\item
  Yol 1 (görsel): nokta var mı? \texttt{7\ /\ 2} = \texttt{3}, \texttt{7.0\ /\ 2.0} = \texttt{3.5}
\item
  Yol 2 (kod): \texttt{7\ as\ f64\ /\ 2.0} → \texttt{3.5}
\item
  Yol 3 (hata): \texttt{E0277:\ cannot\ divide\ \{integer\}\ by\ \{float\}}
\end{itemize}

\textbf{Sık Unutulan 3: Tek tırnak mı çift tırnak mı?}

\begin{itemize}
\tightlist
\item
  Neden unutuluyor: başka dillerde ikisi de string olabilir.
\item
  Yol 1 (analoji): tek tırnak = tek harf (tek kişilik koltuk), çift tırnak = metin.
\item
  Yol 2 (kod): \texttt{let\ c:\ char\ =\ 'a';} ve \texttt{let\ s:\ \&str\ =\ "a";}
\item
  Yol 3 (hata): \texttt{let\ c:\ char\ =\ "a";} → \texttt{E0308:\ mismatched\ types}
\end{itemize}

\textbf{Hazır Derinleştirme Prompt'u:}

\begin{verbatim}
Rust 1.96.0'da Data Types (int, float, bool, char) konusunu en detaylı şekilde anlat. Şunları kapsa:

1. BELLEK SEVİYESİ: Her tip stack'te kaç byte yer kaplar? Two's complement nasıl
   çalışır? IEEE 754 f32/f64 bit düzeni nedir? char ve bool'un geçerli değer
   aralıkları nelerdir?

2. COMPILER DAVRANIŞI: Tür çıkarımı varsayılan tipleri (i32, f64) nasıl seçiyor?
   Overflow kontrolü MIR'de nasıl görünür? overflowing_literals ve
   arithmetic_overflow lint'leri. as dönüşümü bit seviyesinde ne yapar?

3. HATA KODLARI: E0308, E0277, E0604, E0284. Her biri için: ne zaman, neden,
   nasıl çözülür. `rustc --explain EXXXX` çıktısını da ekle.

4. C/C++ KARŞILAŞTIRMASI: int, long, char, signed overflow davranışı. Rust neden
   sabit genişlikli tipler ve açık dönüşüm kullanıyor?

5. ANALOJİ: Günlük hayattan bir benzetme yap (örn: kutu boyutları, kilometre sayacı).

6. RUST 1.96.0: Bu konuda son sürümde değişen bir şey var mı?

7. PRATİK: 3 tane kod örneği ver — kolay, orta, zor.

Açıklama Türkçe, teknik terimler İngilizce, kod yorumları Türkçe olsun.
Her bölümü ayrı başlıkla ver. Kod bloklarını çalıştırılabilir yaz.
\end{verbatim}

\begin{center}\rule{0.5\linewidth}{0.5pt}\end{center}

\subsection{📝 GÖREV (Boş dosyaya yaz)}\label{görev-boş-dosyaya-yaz-2}

\textbf{🟢 Kolay}

\begin{enumerate}
\def\labelenumi{\arabic{enumi}.}
\tightlist
\item
  \texttt{i32}, \texttt{u8}, \texttt{f64}, \texttt{bool} ve \texttt{char} tiplerinde birer variable tanımla (tür yazarak) ve hepsini yazdır.
\item
  \texttt{7\ /\ 2}, \texttt{7\ \%\ 2} ve \texttt{7.0\ /\ 2.0} sonuçlarını yazdır.
\end{enumerate}

\textbf{🟡 Orta}
3. \texttt{255u8} için \texttt{checked\_add(1)}, \texttt{wrapping\_add(1)} ve \texttt{saturating\_add(1)} sonuçlarını yazdır.
4. \texttt{'a'} karakterini \texttt{u32}`ye, \texttt{97u8} değerini \texttt{char}`a çevirip yazdır.

\textbf{🔴 Zor}
5. \texttt{let\ a:\ i32\ =\ 10;\ let\ b:\ i64\ =\ 20;} ile toplam bul: \texttt{E0308}`i gör, sonra \texttt{as\ i64} ile düzelt. Ardından \texttt{3.9\_f64\ as\ i32} ve \texttt{-1\_i32\ as\ u32} sonuçlarını yazdır.

\textbf{Başarı kriteri:} Hepsini hatasız yazarsan → öğrendin ✅

\subsubsection{✅ Kendini Test Et (2-3 saat sonra, sıfırdan)}\label{kendini-test-et-2-3-saat-sonra-sıfırdan-3}

\begin{enumerate}
\def\labelenumi{\arabic{enumi}.}
\tightlist
\item
  Beş temel tipi birer örnekle yaz.
\item
  \texttt{u8} overflow'unu üç farklı yöntemle ele al.
\item
  \texttt{as} ile üç dönüşüm yap ve sonuçlarını tahmin et.
\end{enumerate}

\begin{center}\rule{0.5\linewidth}{0.5pt}\end{center}

\subsection{📚 KAYNAKÇA}\label{kaynakça-3}

\subsubsection{Resmi Kaynaklar}\label{resmi-kaynaklar-3}

\begin{itemize}
\tightlist
\item
  \href{https://doc.rust-lang.org/book/ch03-02-data-types.html}{Rust Book --- Data Types}
\item
  \href{https://doc.rust-lang.org/reference/types/numeric.html}{Rust Reference --- Numeric types}
\item
  \href{https://doc.rust-lang.org/reference/expressions/operator-expr.html\#type-cast-expressions}{Rust Reference --- Type cast expressions}
\item
  \href{https://doc.rust-lang.org/std/primitive.i32.html}{std --- i32 (primitive)}
\end{itemize}

\subsubsection{Sürüm Notları}\label{sürüm-notları-3}

\begin{itemize}
\tightlist
\item
  \href{https://blog.rust-lang.org/2026/05/28/Rust-1.96.0/}{Rust Blog --- Announcing 1.96.0}
\item
  \href{https://doc.rust-lang.org/stable/releases.html\#version-1960-2026-05-28}{Rust 1.96.0 Release Notes}
\end{itemize}

\subsubsection{Hata Kodları}\label{hata-kodları-3}

\begin{itemize}
\tightlist
\item
  \href{https://doc.rust-lang.org/error_codes/E0308.html}{E0308}
\item
  \href{https://doc.rust-lang.org/error_codes/E0277.html}{E0277}
\item
  \href{https://doc.rust-lang.org/error_codes/E0604.html}{E0604}
\item
  \href{https://doc.rust-lang.org/error_codes/E0284.html}{E0284}
\end{itemize}

\subsubsection{Ek Kaynaklar}\label{ek-kaynaklar-3}

\begin{itemize}
\tightlist
\item
  \href{https://doc.rust-lang.org/rust-by-example/primitives.html}{Rust By Example --- Primitives}
\item
  \href{https://jasonwalton.ca/rust-book-abridged/ch03-common-programming-concepts}{Rust Book Abridged --- Common Programming Concepts}
\end{itemize}

\emph{Not: \texttt{i32}/\texttt{f64} varsayılanları, overflow davranışı (debug panic, release wrapping) ve \texttt{char}`ın 4 byte olması aramada Rust Book üzerinden doğrulandı. \texttt{as} dönüşüm örnekleri, boyut tablosu ve hata kodu bağlantıları standart bilgidir; kodları kendi makinende çalıştırarak sonuçları kontrol et.}

\begin{center}\rule{0.5\linewidth}{0.5pt}\end{center}

\section{📘 KONU 5: Operators (+, -, *, /, \%, ==, \&\&, \textbar\textbar)}\label{konu-5-operators}

\textbf{Kategori:} Temel \textbar{} \textbf{Sürüm:} Rust 1.96.0

\begin{center}\rule{0.5\linewidth}{0.5pt}\end{center}

\subsection{🟢 KATMAN 1 --- ILK OKUMA}\label{katman-1-—-ilk-okuma-4}

\begin{Shaded}
\begin{Highlighting}[]
\KeywordTok{fn} \FunctionTok{main}\OperatorTok{()} \OperatorTok{\{}
    \KeywordTok{let}\NormalTok{ a }\OperatorTok{=} \DecValTok{17}\OperatorTok{;}
    \KeywordTok{let}\NormalTok{ b }\OperatorTok{=} \DecValTok{5}\OperatorTok{;}

\NormalTok{    println}\OperatorTok{!(}\StringTok{"\{\}"}\OperatorTok{,}\NormalTok{ a }\OperatorTok{+}\NormalTok{ b}\OperatorTok{);}            \CommentTok{// toplama}
\NormalTok{    println}\OperatorTok{!(}\StringTok{"\{\}"}\OperatorTok{,}\NormalTok{ a }\OperatorTok{-}\NormalTok{ b}\OperatorTok{);}            \CommentTok{// çıkarma}
\NormalTok{    println}\OperatorTok{!(}\StringTok{"\{\}"}\OperatorTok{,}\NormalTok{ a }\OperatorTok{*}\NormalTok{ b}\OperatorTok{);}            \CommentTok{// çarpma}
\NormalTok{    println}\OperatorTok{!(}\StringTok{"\{\}"}\OperatorTok{,}\NormalTok{ a }\OperatorTok{/}\NormalTok{ b}\OperatorTok{);}            \CommentTok{// tam sayı bölmesi}
\NormalTok{    println}\OperatorTok{!(}\StringTok{"\{\}"}\OperatorTok{,}\NormalTok{ a }\OperatorTok{\%}\NormalTok{ b}\OperatorTok{);}            \CommentTok{// kalan}
\NormalTok{    println}\OperatorTok{!(}\StringTok{"\{\}"}\OperatorTok{,}\NormalTok{ a }\OperatorTok{==}\NormalTok{ b}\OperatorTok{);}           \CommentTok{// eşit mi?}
\NormalTok{    println}\OperatorTok{!(}\StringTok{"\{\}"}\OperatorTok{,}\NormalTok{ a }\OperatorTok{!=}\NormalTok{ b}\OperatorTok{);}           \CommentTok{// eşit değil mi?}
\NormalTok{    println}\OperatorTok{!(}\StringTok{"\{\}"}\OperatorTok{,}\NormalTok{ a }\OperatorTok{\textgreater{}}\NormalTok{ b }\OperatorTok{\&\&}\NormalTok{ b }\OperatorTok{\textgreater{}} \DecValTok{0}\OperatorTok{);}   \CommentTok{// ve}
\NormalTok{    println}\OperatorTok{!(}\StringTok{"\{\}"}\OperatorTok{,}\NormalTok{ a }\OperatorTok{\textless{}}\NormalTok{ b }\OperatorTok{||}\NormalTok{ b }\OperatorTok{\textgreater{}} \DecValTok{0}\OperatorTok{);}   \CommentTok{// veya}
\NormalTok{    println}\OperatorTok{!(}\StringTok{"\{\}"}\OperatorTok{,} \OperatorTok{!(}\NormalTok{a }\OperatorTok{\textgreater{}}\NormalTok{ b}\OperatorTok{));}         \CommentTok{// değil}

    \KeywordTok{let} \KeywordTok{mut}\NormalTok{ c }\OperatorTok{=} \DecValTok{10}\OperatorTok{;}
\NormalTok{    c }\OperatorTok{+=} \DecValTok{5}\OperatorTok{;}                           \CommentTok{// c = c + 5}
\NormalTok{    c }\OperatorTok{*=} \DecValTok{2}\OperatorTok{;}                           \CommentTok{// c = c * 2}
\NormalTok{    println}\OperatorTok{!(}\StringTok{"\{c\}"}\OperatorTok{);}

\NormalTok{    println}\OperatorTok{!(}\StringTok{"\{\}"}\OperatorTok{,} \DecValTok{1} \OperatorTok{\textless{}\textless{}} \DecValTok{3}\OperatorTok{);}           \CommentTok{// bit kaydırma}
\NormalTok{    println}\OperatorTok{!(}\StringTok{"\{\}"}\OperatorTok{,} \BaseNTok{0b1100} \OperatorTok{\&} \BaseNTok{0b1010}\OperatorTok{);}  \CommentTok{// bit düzeyinde ve}
\OperatorTok{\}}
\CommentTok{// Çıktı:}
\CommentTok{// 22}
\CommentTok{// 12}
\CommentTok{// 85}
\CommentTok{// 3}
\CommentTok{// 2}
\CommentTok{// false}
\CommentTok{// true}
\CommentTok{// true}
\CommentTok{// true}
\CommentTok{// false}
\CommentTok{// 30}
\CommentTok{// 8}
\CommentTok{// 8}
\end{Highlighting}
\end{Shaded}

Operator (işleç), operand (işlenen) değerler üzerinde işlem yapan sembollerdir. Aritmetik işleçler \texttt{+\ -\ *\ /\ \%}, karşılaştırma işleçleri \texttt{==\ !=\ \textless{}\ \textgreater{}\ \textless{}=\ \textgreater{}=} ve mantıksal işleçler \texttt{\&\&\ \textbar{}\textbar{}\ !} olarak üç ana gruba ayrılır. Karşılaştırma ve mantıksal işleçler her zaman \texttt{bool} üretir. \texttt{+=} gibi compound assignment (bileşik atama) işleçleri değişkeni yerinde günceller ve variable'ın \texttt{mut} olmasını ister. \texttt{\&\&} ve \texttt{\textbar{}\textbar{}} short-circuit (kısa devre) çalışır: sol taraf sonucu belirliyorsa sağ tarafa bakmaz.

\textbf{1.96.0 notu:} İşleç kuralları 1.96.0'da aynı.

\begin{center}\rule{0.5\linewidth}{0.5pt}\end{center}

\subsection{🔵 KATMAN 2 --- RECALL KARTLARI}\label{katman-2-—-recall-kartlari-4}

\begin{itemize}
\item
  Aritmetik: \texttt{+\ -\ *\ /\ \%}

\begin{Shaded}
\begin{Highlighting}[]
\DecValTok{7} \OperatorTok{/} \DecValTok{2}      \CommentTok{// 3   (tam sayı bölmesi, sıfıra doğru kesilir)}
\OperatorTok{-}\DecValTok{7} \OperatorTok{/} \DecValTok{2}     \CommentTok{// -3}
\FloatTok{7.0} \OperatorTok{/} \FloatTok{2.0}  \CommentTok{// 3.5}
\end{Highlighting}
\end{Shaded}


  Hata (runtime): sıfıra bölme → panic \texttt{attempt\ to\ divide\ by\ zero}
\item
  Kalan (\texttt{\%}) işareti soldaki sayıyı izler:

\begin{Shaded}
\begin{Highlighting}[]
\OperatorTok{-}\DecValTok{7} \OperatorTok{\%} \DecValTok{3}               \CommentTok{// -1}
\OperatorTok{(-}\DecValTok{7\_i}\DecValTok{32}\OperatorTok{).}\FunctionTok{rem\_euclid}\OperatorTok{(}\DecValTok{3}\OperatorTok{)} \CommentTok{// 2}
\FloatTok{7.5} \OperatorTok{\%} \FloatTok{2.0}            \CommentTok{// 1.5 (float'ta da çalışır)}
\end{Highlighting}
\end{Shaded}
\item
  Karşılaştırma: iki taraf \textbf{aynı tipte} olmalı

\begin{Shaded}
\begin{Highlighting}[]
\DecValTok{1} \OperatorTok{==} \FloatTok{1.0}     \CommentTok{// E0277: can't compare `\{integer\}` with `\{float\}`}
\end{Highlighting}
\end{Shaded}


  İlgili: \texttt{PartialEq}, \texttt{PartialOrd}
\item
  Mantıksal: \texttt{\&\&} (ve), \texttt{\textbar{}\textbar{}} (veya), \texttt{!} (değil), hepsi \texttt{bool} ister

\begin{Shaded}
\begin{Highlighting}[]
\NormalTok{v}\OperatorTok{.}\FunctionTok{len}\OperatorTok{()} \OperatorTok{\textgreater{}} \DecValTok{0} \OperatorTok{\&\&}\NormalTok{ v}\OperatorTok{[}\DecValTok{0}\OperatorTok{]} \OperatorTok{==} \DecValTok{1}   \CommentTok{// sol false ise sağ hiç çalışmaz}
\end{Highlighting}
\end{Shaded}


  İlgili: \texttt{\&} ve \texttt{\textbar{}} bool üzerinde short-circuit yapmaz
\item
  Bit düzeyinde ve kaydırma:

\begin{Shaded}
\begin{Highlighting}[]
\BaseNTok{0b1100} \OperatorTok{\&} \BaseNTok{0b1010}   \CommentTok{// 8  (ve)}
\BaseNTok{0b1100} \OperatorTok{|} \BaseNTok{0b1010}   \CommentTok{// 14 (veya)}
\BaseNTok{0b1100} \OperatorTok{\^{}} \BaseNTok{0b1010}   \CommentTok{// 6  (xor)}
\OperatorTok{!}\DecValTok{0u}\DecValTok{8}              \CommentTok{// 255 (bit tersi)}
\DecValTok{1} \OperatorTok{\textless{}\textless{}} \DecValTok{3}            \CommentTok{// 8}
\DecValTok{16} \OperatorTok{\textgreater{}\textgreater{}} \DecValTok{2}           \CommentTok{// 4}
\end{Highlighting}
\end{Shaded}


  Hata (runtime): \texttt{1u8\ \textless{}\textless{}\ 8} → panic \texttt{attempt\ to\ shift\ left\ with\ overflow}
\item
  Compound assignment: \texttt{+=\ -=\ *=\ /=\ \%=\ \&=\ \textbar{}=\ \^{}=\ \textless{}\textless{}=\ \textgreater{}\textgreater{}=}
  Hata: \texttt{let\ x\ =\ 1;\ x\ +=\ 1;} → \texttt{E0384}: cannot assign twice to immutable variable
  Hata: \texttt{x++} → \texttt{error:\ Rust\ has\ no\ postfix\ increment\ operator} → \texttt{x\ +=\ 1} yaz
\item
  Öncelik (yüksekten düşüğe):

\begin{verbatim}
. () [] ?      (metot, çağrı, indeks, ?)
- ! * &        (tekli işleçler)
as
* / %
+ -
<< >>
&
^
|
== != < > <= >=   (zincirlenemez)
&&
||
.. ..=
= += -= ...
\end{verbatim}


  Şüphedeysen parantez kullan.
\item
  \texttt{as} çarpmadan \textbf{önce} gelir; \texttt{as} sonrası \texttt{\textless{}} sorun çıkarır:

\begin{Shaded}
\begin{Highlighting}[]
\NormalTok{a }\KeywordTok{as} \DataTypeTok{u64} \OperatorTok{*}\NormalTok{ b          }\CommentTok{// (a as u64) * b}
\OperatorTok{(}\NormalTok{a }\KeywordTok{as} \DataTypeTok{u64}\OperatorTok{)} \OperatorTok{\textless{}}\NormalTok{ b        }\CommentTok{// parantez şart}
\end{Highlighting}
\end{Shaded}


  Hata: \texttt{a\ as\ u64\ \textless{}\ b} → \texttt{error:\ \textless{}\ is\ interpreted\ as\ a\ start\ of\ generic\ arguments\ for\ u64,\ not\ a\ comparison}
\item
  Zincirleme karşılaştırma yasak:

\begin{Shaded}
\begin{Highlighting}[]
\DecValTok{1} \OperatorTok{\textless{}}\NormalTok{ x }\OperatorTok{\textless{}} \DecValTok{10}             \CommentTok{// error: comparison operators cannot be chained}
\DecValTok{1} \OperatorTok{\textless{}}\NormalTok{ x }\OperatorTok{\&\&}\NormalTok{ x }\OperatorTok{\textless{}} \DecValTok{10}        \CommentTok{// doğrusu}
\end{Highlighting}
\end{Shaded}
\item
  \texttt{=} ifade olarak \texttt{()} döner, koşul \texttt{bool} ister:
  Hata: \texttt{if\ x\ =\ 5\ \{\ \}} → \texttt{E0308}: mismatched types (expected \texttt{bool}, found \texttt{()})
\item
  Kendi tipin için işleç:
  Hata: \texttt{p1\ +\ p2} (\texttt{Point}) → \texttt{E0369}: binary operation \texttt{+} cannot be applied to type \texttt{Point}
  İlgili: \texttt{impl\ Add\ for\ Point}
\end{itemize}

\textbf{GoT bağlantı haritası:}

\begin{verbatim}
operator ──┬── arithmetic (+ - * / %) ──► overflow / sıfıra bölme panic
           ├── comparison (== != < > <= >=) ──► bool ──► if / while
           ├── logical (&& || !) ──► short-circuit
           ├── bitwise (& | ^ << >>)
           ├── compound (+= ...) ──► mut gerekir
           └── as ──► tür dönüştürme (öncelik yüksek)
trait'ler: Add / PartialEq / PartialOrd ──► E0369
\end{verbatim}

\begin{center}\rule{0.5\linewidth}{0.5pt}\end{center}

\subsection{🟣 KATMAN 3 --- DERIN TEKNIK NOT}\label{katman-3-—-derin-teknik-not-4}

\textbf{Adım 1 --- Bellek Seviyesi}

\begin{itemize}
\tightlist
\item
  Önce soru: ``Bellekte ne oluyor?'' Cevap: primitive tiplerde işleçler doğrudan CPU komutuna dönüşür (\texttt{add}, \texttt{sub}, \texttt{imul}, \texttt{idiv}, \texttt{cmp}). İşlenenler çoğunlukla register'larda yaşar, stack'e çoğu zaman gerek kalmaz.
\item
  Debug build'de toplama sonrasına bir overflow (taşma) kontrolü eklenir (örneğin \texttt{jo}, yani taşma varsa atla). Release'te bu kontrol yoktur.
\item
  \texttt{i32::MIN\ /\ -1} donanımda tuzak (trap) oluşturur; Rust bunu önceden kontrol edip panic'e çevirir.
\item
  \texttt{\&\&} ve \texttt{\textbar{}\textbar{}} koşullu atlama (branch) olarak derlenir; sol taraf sonucu belirliyorsa sağ tarafın kodu atlanır.
\item
  İşaretsiz sayıda \texttt{x\ /\ 8} gibi ifadeler optimizasyonla bit kaydırmaya dönüşebilir.
\end{itemize}

\textbf{Adım 2 --- Compiler Davranışı}

\begin{itemize}
\tightlist
\item
  Önce soru: ``Compiler \texttt{a\ +\ b}`yi nasıl görüyor?'' Cevap: primitive tiplerde yerleşik işleç; kendi tiplerinde ise trait (özellik) çağrısı: \texttt{a\ +\ b} → \texttt{Add::add(a,\ b)}, \texttt{a\ +=\ b} → \texttt{AddAssign}, \texttt{a\ ==\ b} → \texttt{PartialEq::eq(\&a,\ \&b)}, \texttt{a\ \textless{}\ b} → \texttt{PartialOrd}.
\item
  \texttt{==} referans alır, bu yüzden \texttt{s1\ ==\ s2} String'leri \textbf{taşımaz}. \texttt{s1\ +\ \&s2} ise \texttt{s1}`i taşır.
\item
  Öncelik parser aşamasında çözülür. Karşılaştırma işleçleri associative (birleşme yönü) olmadığı için zincirlenemez.
\item
  Compile time'da bilinen ifadeler (\texttt{2\ +\ 3}) constant folding ile hesaplanır.
\item
  MIR'de aritmetik \texttt{Add}, \texttt{Rem}, \texttt{Lt} gibi işlemler olarak, kontrollü olanlar \texttt{AddWithOverflow} + \texttt{assert} olarak görünür. Görmek için: \texttt{cargo\ rustc\ --\ --emit=mir}.
\end{itemize}

\textbf{Adım 3 --- Hata Kodları}

\begin{itemize}
\tightlist
\item
  \textbf{E0369:} işleç o tip için tanımlı değil. Çözüm: ilgili trait'i implement et (\texttt{impl\ Add\ for\ Point}) veya alanlara ayrı ayrı işlem yap. Detay: \texttt{rustc\ --explain\ E0369}.
\item
  \textbf{E0277:} tip çifti için işlem yok (\texttt{\{integer\}} ile \texttt{\{float\}} karşılaştırması). Çözüm: tipleri eşitle.
\item
  \textbf{E0308:} \texttt{i32\ +\ i64} veya \texttt{if\ x\ =\ 5}. Çözüm: aynı tipe çevir / \texttt{==} yaz.
\item
  \textbf{E0384:} \texttt{mut} olmayan variable'a \texttt{+=}. Çözüm: \texttt{let\ mut}.
\item
  \textbf{E0600:} \texttt{!} işleci o tipe uygulanamaz (örneğin \texttt{!"text"}). Çözüm: \texttt{bool} veya tam sayıya uygula.
\item
  \textbf{Kodsuz:} \texttt{comparison\ operators\ cannot\ be\ chained}, \texttt{Rust\ has\ no\ postfix\ increment\ operator}, \texttt{\textless{}\ is\ interpreted\ as\ a\ start\ of\ generic\ arguments}
\item
  \textbf{Runtime panic:} \texttt{attempt\ to\ divide\ by\ zero}, \texttt{attempt\ to\ calculate\ the\ remainder\ with\ a\ divisor\ of\ zero}, \texttt{attempt\ to\ shift\ left\ with\ overflow}
\end{itemize}

\textbf{Adım 4 --- ``Neden Böyle?''}

\begin{itemize}
\tightlist
\item
  C/C++`ta \texttt{if\ (x\ =\ 5)} geçerli bir koşuldur ve klasik bir hata kaynağıdır. Rust'ta koşul \texttt{bool} olmak zorunda olduğu için derlenmez.
\item
  C'de \texttt{x++} ifadesinin değerlendirme sırası karışık davranışlara yol açabilir. Rust bu operatörü hiç sunmaz, \texttt{x\ +=\ 1} kullanırsın.
\item
  C'de \texttt{a\ \&\ b\ ==\ c}, \texttt{a\ \&\ (b\ ==\ c)} olarak ayrışır (sürpriz). Rust'ta \texttt{\&} karşılaştırmadan önce gelir.
\item
  Rust'ta tipler arası otomatik dönüşüm yok: \texttt{1\ ==\ 1.0} derlenmez.
\item
  Trade-off: biraz daha çok yazarsın, ama anlam belirsizliği azalır.
\end{itemize}

\textbf{Adım 5 --- Rust 1.96.0 Farkı}

\begin{itemize}
\tightlist
\item
  Rust 1.96.0'daki değişiklikler (yeni Range* tipleri, \texttt{assert\_matches!} makrosu, Wasm linker güncellemeleri) operatör önceliğini ve temel overflow davranışını değiştirmiyor.
\item
  Range* tiplerinin stabilize edilmesi range operatörleriyle ilgili API'leri etkiler; ancak diğer operatör kuralları, öncelik tablosu ve overflow davranışı değişmedi.
\end{itemize}

\begin{center}\rule{0.5\linewidth}{0.5pt}\end{center}

\subsection{🔴 KATMAN 4 --- SIK UNUTULANLAR + DERINLEŞTIRME}\label{katman-4-—-sik-unutulanlar-derinleştirme-4}

\textbf{Sık Unutulan 1: \texttt{=} ile \texttt{==} karışması}

\begin{itemize}
\tightlist
\item
  Neden unutuluyor: başka dillerde \texttt{if\ (x\ =\ 5)} bazen sessizce çalışır.
\item
  Yol 1 (analoji): \texttt{=} ``ata'' (emir), \texttt{==} ``eşit mi?'' (soru).
\item
  Yol 2 (kod): \texttt{if\ x\ ==\ 5\ \{\ \}}
\item
  Yol 3 (hata): \texttt{if\ x\ =\ 5\ \{\ \}} → \texttt{E0308:\ mismatched\ types,\ expected\ bool,\ found\ ()}
\end{itemize}

\textbf{Sık Unutulan 2: \texttt{x++} yok}

\begin{itemize}
\tightlist
\item
  Neden unutuluyor: C, Java, JavaScript alışkanlığı.
\item
  Yol 1 (görsel): \texttt{++} yok, ama \texttt{+=\ 1} var.
\item
  Yol 2 (kod): \texttt{let\ mut\ x\ =\ 0;\ x\ +=\ 1;}
\item
  Yol 3 (hata): \texttt{x++} → \texttt{error:\ Rust\ has\ no\ postfix\ increment\ operator}
\end{itemize}

\textbf{Sık Unutulan 3: \texttt{\&\&} mi \texttt{\&} mi?}

\begin{itemize}
\tightlist
\item
  Neden unutuluyor: \texttt{bool} üzerinde ikisi de çalışır.
\item
  Yol 1 (analoji): \texttt{\&\&} = önce kapıyı kontrol et, kapalıysa içeri bakma. \texttt{\&} = her durumda içeri bak.
\item
  Yol 2 (kod): \texttt{v.len()\ \textgreater{}\ 0\ \&\&\ v{[}0{]}\ ==\ 1} güvenli; \texttt{\&} ile yazarsan \texttt{v{[}0{]}} boş vector'de panic yapar.
\item
  Yol 3 (hata): \texttt{index\ out\ of\ bounds:\ the\ len\ is\ 0\ but\ the\ index\ is\ 0}
\end{itemize}

\textbf{Hazır Derinleştirme Prompt'u:}

\begin{verbatim}
Rust 1.96.0'da Operators (+, -, *, /, %, ==, &&, ||) konusunu en detaylı şekilde anlat. Şunları kapsa:

1. BELLEK SEVİYESİ: İşleçler CPU seviyesinde nasıl çalışır? Overflow kontrolü
   assembly'de nasıl görünür? Short-circuit branch olarak nasıl derlenir?

2. COMPILER DAVRANIŞI: a + b, Add::add'e nasıl dönüşür? == neden referans alır?
   Öncelik tablosu ve as'ın önceliği. MIR çıktısında işleçler nasıl görünür?
   (--emit=mir)

3. HATA KODLARI: E0369, E0277, E0308, E0384, E0600. Her biri için: ne zaman, neden,
   nasıl çözülür. `rustc --explain EXXXX` çıktısını da ekle.

4. C/C++ KARŞILAŞTIRMASI: x++, if (x = 5), a & b == c. Rust neden farklı?

5. ANALOJİ: Günlük hayattan bir benzetme yap (örn: kapı kontrolü, hesap makinesi).

6. RUST 1.96.0: Bu konuda son sürümde değişen bir şey var mı?

7. PRATİK: 3 tane kod örneği ver — kolay, orta, zor.

Açıklama Türkçe, teknik terimler İngilizce, kod yorumları Türkçe olsun.
Her bölümü ayrı başlıkla ver. Kod bloklarını çalıştırılabilir yaz.
\end{verbatim}

\begin{center}\rule{0.5\linewidth}{0.5pt}\end{center}

\subsection{📝 GÖREV (Boş dosyaya yaz)}\label{görev-boş-dosyaya-yaz-3}

\textbf{🟢 Kolay}

\begin{enumerate}
\def\labelenumi{\arabic{enumi}.}
\tightlist
\item
  İki \texttt{i32} variable tanımla; \texttt{+\ -\ *\ /\ \%} sonuçlarını yazdır.
\item
  \texttt{==}, \texttt{!=}, \texttt{\textless{}}, \texttt{\textgreater{}=} karşılaştırmalarının sonuçlarını yazdır.
\end{enumerate}

\textbf{🟡 Orta}
3. \texttt{let\ mut\ score\ =\ 10;} ile \texttt{+=}, \texttt{-=}, \texttt{*=} kullanarak değeri değiştir ve her adımda yazdır.
4. \texttt{age\ \textgreater{}=\ 18\ \&\&\ has\_id} ve \texttt{is\_weekend\ \textbar{}\textbar{}\ is\_holiday} ifadelerini \texttt{bool} variable'larıyla yaz ve yazdır.

\textbf{🔴 Zor}
5. \texttt{-7\ /\ 2}, \texttt{-7\ \%\ 3} ve \texttt{(-7\_i32).rem\_euclid(3)} sonuçlarını yazdır. Sonra \texttt{let\ v:\ Vec\textless{}i32\textgreater{}\ =\ vec!{[}{]};} için \texttt{v.len()\ \textgreater{}\ 0\ \&\&\ v{[}0{]}\ ==\ 1} ifadesini panic olmadan yazdır. Son olarak \texttt{a\ as\ u64\ *\ b} ile \texttt{(a\ as\ u64)\ \textless{}\ b} kullan.

\textbf{Başarı kriteri:} Hepsini hatasız yazarsan → öğrendin ✅

\subsubsection{✅ Kendini Test Et (2-3 saat sonra, sıfırdan)}\label{kendini-test-et-2-3-saat-sonra-sıfırdan-4}

\begin{enumerate}
\def\labelenumi{\arabic{enumi}.}
\tightlist
\item
  Beş aritmetik işleci bir örnekle yaz.
\item
  \texttt{\&\&}, \texttt{\textbar{}\textbar{}}, \texttt{!} ile bir koşul yaz.
\item
  \texttt{x\ +=\ 1} yaz ve \texttt{x++} yazınca ne olduğunu söyle.
\end{enumerate}

\begin{center}\rule{0.5\linewidth}{0.5pt}\end{center}

\subsection{📚 KAYNAKÇA}\label{kaynakça-4}

\subsubsection{Resmi Kaynaklar}\label{resmi-kaynaklar-4}

\begin{itemize}
\tightlist
\item
  \href{https://doc.rust-lang.org/book/appendix-02-operators.html}{Rust Book --- Appendix B: Operators and Symbols}
\item
  \href{https://doc.rust-lang.org/reference/expressions/operator-expr.html}{Rust Reference --- Operator expressions}
\item
  \href{https://doc.rust-lang.org/reference/expressions.html\#expression-precedence}{Rust Reference --- Expression precedence}
\item
  \href{https://doc.rust-lang.org/std/ops/index.html}{std::ops}
\end{itemize}

\subsubsection{Sürüm Notları}\label{sürüm-notları-4}

\begin{itemize}
\tightlist
\item
  \href{https://blog.rust-lang.org/2026/05/28/Rust-1.96.0/}{Rust Blog --- Announcing 1.96.0}
\item
  \href{https://doc.rust-lang.org/stable/releases.html\#version-1960-2026-05-28}{Rust 1.96.0 Release Notes}
\end{itemize}

\subsubsection{Hata Kodları}\label{hata-kodları-4}

\begin{itemize}
\tightlist
\item
  \href{https://doc.rust-lang.org/error_codes/E0369.html}{E0369}
\item
  \href{https://doc.rust-lang.org/error_codes/E0277.html}{E0277}
\item
  \href{https://doc.rust-lang.org/error_codes/E0308.html}{E0308}
\item
  \href{https://doc.rust-lang.org/error_codes/E0384.html}{E0384}
\item
  \href{https://doc.rust-lang.org/error_codes/E0600.html}{E0600}
\end{itemize}

\subsubsection{Ek Kaynaklar}\label{ek-kaynaklar-4}

\begin{itemize}
\tightlist
\item
  \href{https://doc.rust-lang.org/rust-by-example/primitives.html}{Rust By Example --- Primitives (operators)}
\item
  \href{https://m.appbook.qq.com/read/1036699272/45}{Operator expressions (kısa devre açıklaması)}
\end{itemize}

\emph{Not: Bu konuda arama sonuçları zayıftı. \texttt{\&\&} ve \texttt{\textbar{}\textbar{}} işleçlerinin kısa devre çalıştığı ve \texttt{\&}/\texttt{\textbar{}} yerine tercih edilmesi gerektiği bilgisi sadece ikincil bir kaynakta doğrulandı. Öncelik tablosu, hata kodları ve hata mesajları bilgime dayanıyor ve resmi Reference sayfasıyla tam karşılaştırılmadı. Öncelik tablosunu Reference'taki ``Expression precedence'' bölümünden, hata mesajlarını da kendi makinende çalıştırarak kontrol et.}

\begin{center}\rule{0.5\linewidth}{0.5pt}\end{center}

\section{📘 KONU 6: String vs \&str}\label{konu-6-string-vs-str}

\textbf{Kategori:} Temel \textbar{} \textbf{Sürüm:} Rust 1.96.0

\begin{center}\rule{0.5\linewidth}{0.5pt}\end{center}

\subsection{🟢 KATMAN 1 --- ILK OKUMA}\label{katman-1-—-ilk-okuma-5}

\begin{Shaded}
\begin{Highlighting}[]
\KeywordTok{fn} \FunctionTok{main}\OperatorTok{()} \OperatorTok{\{}
    \KeywordTok{let}\NormalTok{ literal}\OperatorTok{:} \OperatorTok{\&}\DataTypeTok{str} \OperatorTok{=} \StringTok{"hello"}\OperatorTok{;}                    \CommentTok{// sabit metin (borrowed)}
    \KeywordTok{let} \KeywordTok{mut}\NormalTok{ owned}\OperatorTok{:} \DataTypeTok{String} \OperatorTok{=} \DataTypeTok{String}\OperatorTok{::}\FunctionTok{from}\OperatorTok{(}\StringTok{"hello"}\OperatorTok{);}  \CommentTok{// heap'te büyüyebilen metin}
\NormalTok{    owned}\OperatorTok{.}\FunctionTok{push\_str}\OperatorTok{(}\StringTok{", world"}\OperatorTok{);}                      \CommentTok{// \&str ekler}
\NormalTok{    owned}\OperatorTok{.}\FunctionTok{push}\OperatorTok{(}\CharTok{'!'}\OperatorTok{);}                                \CommentTok{// tek char ekler}

    \KeywordTok{let}\NormalTok{ slice}\OperatorTok{:} \OperatorTok{\&}\DataTypeTok{str} \OperatorTok{=} \OperatorTok{\&}\NormalTok{owned}\OperatorTok{[}\DecValTok{0}\OperatorTok{..}\DecValTok{5}\OperatorTok{];}                 \CommentTok{// String'in bir parçası}

\NormalTok{    println}\OperatorTok{!(}\StringTok{"\{literal\}"}\OperatorTok{);}
\NormalTok{    println}\OperatorTok{!(}\StringTok{"\{owned\}"}\OperatorTok{);}
\NormalTok{    println}\OperatorTok{!(}\StringTok{"\{slice\}"}\OperatorTok{);}
\NormalTok{    println}\OperatorTok{!(}\StringTok{"\{\}"}\OperatorTok{,}\NormalTok{ owned}\OperatorTok{.}\FunctionTok{len}\OperatorTok{());}                    \CommentTok{// byte sayısı}
\NormalTok{    println}\OperatorTok{!(}\StringTok{"\{\}"}\OperatorTok{,} \StringTok{"héllo"}\OperatorTok{.}\FunctionTok{len}\OperatorTok{());}                  \CommentTok{// é iki byte}
\NormalTok{    println}\OperatorTok{!(}\StringTok{"\{\}"}\OperatorTok{,} \StringTok{"héllo"}\OperatorTok{.}\FunctionTok{chars}\OperatorTok{().}\FunctionTok{count}\OperatorTok{());}        \CommentTok{// karakter sayısı}

    \FunctionTok{greet}\OperatorTok{(\&}\NormalTok{owned}\OperatorTok{);}                                  \CommentTok{// \&String → \&str}
    \FunctionTok{greet}\OperatorTok{(}\StringTok{"literal"}\OperatorTok{);}                               \CommentTok{// \&str direkt}
\OperatorTok{\}}

\KeywordTok{fn} \FunctionTok{greet}\OperatorTok{(}\NormalTok{name}\OperatorTok{:} \OperatorTok{\&}\DataTypeTok{str}\OperatorTok{)} \OperatorTok{\{}
\NormalTok{    println}\OperatorTok{!(}\StringTok{"Hi, \{name\}!"}\OperatorTok{);}
\OperatorTok{\}}
\CommentTok{// Çıktı:}
\CommentTok{// hello}
\CommentTok{// hello, world!}
\CommentTok{// hello}
\CommentTok{// 13}
\CommentTok{// 6}
\CommentTok{// 5}
\CommentTok{// Hi, hello, world!!}
\CommentTok{// Hi, literal!}
\end{Highlighting}
\end{Shaded}

Rust'ta iki ana metin tipi vardır: \texttt{String} owned (sahipli) bir tiptir, heap'te yaşar ve büyüyebilir. \texttt{\&str} ise string slice (metin dilimi) denen, başka bir yerdeki UTF-8 metne borrowed (ödünç alınmış) bir bakıştır; \texttt{"hello"} gibi literal'ler bu tiptedir. Bir fonksiyon sadece okuyacaksa parametresini \texttt{\&str} yaz, çünkü hem \texttt{\&String} hem literal kabul eder. \texttt{len()} karakter değil byte sayısını verir; \texttt{"héllo"} için \texttt{6} çıkar çünkü \texttt{é} iki byte'tır. Metnin \texttt{n}. karakterine \texttt{s{[}n{]}} ile erişemezsin; bunun nedenini Katman 3'te göreceksin.

\textbf{1.96.0 notu:} \texttt{String} ve \texttt{\&str} davranışı 1.96.0'da aynı.

\begin{center}\rule{0.5\linewidth}{0.5pt}\end{center}

\subsection{🔵 KATMAN 2 --- RECALL KARTLARI}\label{katman-2-—-recall-kartlari-5}

\begin{itemize}
\item
  Oluşturma:

\begin{Shaded}
\begin{Highlighting}[]
\KeywordTok{let}\NormalTok{ a}\OperatorTok{:} \OperatorTok{\&}\DataTypeTok{str} \OperatorTok{=} \StringTok{"hi"}\OperatorTok{;}                 \CommentTok{// literal}
\KeywordTok{let}\NormalTok{ b }\OperatorTok{=} \DataTypeTok{String}\OperatorTok{::}\FunctionTok{from}\OperatorTok{(}\StringTok{"hi"}\OperatorTok{);}         \CommentTok{// \&str → String}
\KeywordTok{let}\NormalTok{ c }\OperatorTok{=} \StringTok{"hi"}\OperatorTok{.}\FunctionTok{to\_string}\OperatorTok{();}           \CommentTok{// aynı iş}
\KeywordTok{let}\NormalTok{ d }\OperatorTok{=} \StringTok{"hi"}\OperatorTok{.}\FunctionTok{to\_owned}\OperatorTok{();}            \CommentTok{// aynı iş}
\KeywordTok{let}\NormalTok{ e }\OperatorTok{=} \DataTypeTok{String}\OperatorTok{::}\FunctionTok{new}\OperatorTok{();}              \CommentTok{// boş}
\KeywordTok{let}\NormalTok{ f }\OperatorTok{=} \DataTypeTok{String}\OperatorTok{::}\FunctionTok{with\_capacity}\OperatorTok{(}\DecValTok{32}\OperatorTok{);}  \CommentTok{// önceden yer ayır}
\end{Highlighting}
\end{Shaded}


  Hata: \texttt{let\ s:\ String\ =\ "hi";} → \texttt{E0308}: expected \texttt{String}, found \texttt{\&str}
\item
  Ekleme:

\begin{Shaded}
\begin{Highlighting}[]
\KeywordTok{let} \KeywordTok{mut}\NormalTok{ s }\OperatorTok{=} \DataTypeTok{String}\OperatorTok{::}\FunctionTok{from}\OperatorTok{(}\StringTok{"a"}\OperatorTok{);}
\NormalTok{s}\OperatorTok{.}\FunctionTok{push\_str}\OperatorTok{(}\StringTok{"bc"}\OperatorTok{);}   \CommentTok{// \&str ekler}
\NormalTok{s}\OperatorTok{.}\FunctionTok{push}\OperatorTok{(}\CharTok{'d'}\OperatorTok{);}        \CommentTok{// char ekler}
\NormalTok{s }\OperatorTok{+=} \StringTok{"ef"}\OperatorTok{;}          \CommentTok{// += da \&str ister}
\end{Highlighting}
\end{Shaded}


  Hata: \texttt{s} \texttt{mut} değilse → \texttt{E0596}: cannot borrow \texttt{s} as mutable, as it is not declared as mutable
\item
  Birleştirme:

\begin{Shaded}
\begin{Highlighting}[]
\KeywordTok{let}\NormalTok{ s3 }\OperatorTok{=}\NormalTok{ s1 }\OperatorTok{+} \OperatorTok{\&}\NormalTok{s2}\OperatorTok{;}                  \CommentTok{// s1 TAŞINIR, s2 kalır}
\KeywordTok{let}\NormalTok{ s4 }\OperatorTok{=}\NormalTok{ format}\OperatorTok{!(}\StringTok{"\{s2\}-\{s3\}"}\OperatorTok{);}      \CommentTok{// hiçbiri taşınmaz}
\end{Highlighting}
\end{Shaded}


  Hata: \texttt{s1} sonra kullanılırsa → \texttt{E0382}: borrow of moved value
  Hata: \texttt{"a"\ +\ "b"} → \texttt{E0369}: cannot add \texttt{\&str} to \texttt{\&str}
\item
  Fonksiyon parametresi:

\begin{Shaded}
\begin{Highlighting}[]
\KeywordTok{fn} \FunctionTok{f}\OperatorTok{(}\NormalTok{s}\OperatorTok{:} \OperatorTok{\&}\DataTypeTok{str}\OperatorTok{)} \OperatorTok{\{\}}
\FunctionTok{f}\OperatorTok{(}\StringTok{"lit"}\OperatorTok{);}        \CommentTok{// OK}
\FunctionTok{f}\OperatorTok{(\&}\NormalTok{owned}\OperatorTok{);}       \CommentTok{// OK: \&String → \&str (deref coercion)}
\FunctionTok{f}\OperatorTok{(}\NormalTok{owned}\OperatorTok{);}        \CommentTok{// E0308: expected \&str, found String}
\end{Highlighting}
\end{Shaded}
\item
  Uzunluk:

\begin{Shaded}
\begin{Highlighting}[]
\StringTok{"héllo"}\OperatorTok{.}\FunctionTok{len}\OperatorTok{()}              \CommentTok{// 6  (byte)}
\StringTok{"héllo"}\OperatorTok{.}\FunctionTok{chars}\OperatorTok{().}\FunctionTok{count}\OperatorTok{()}    \CommentTok{// 5  (karakter)}
\end{Highlighting}
\end{Shaded}
\item
  Indexleme yok:

\begin{Shaded}
\begin{Highlighting}[]
\KeywordTok{let}\NormalTok{ s }\OperatorTok{=} \DataTypeTok{String}\OperatorTok{::}\FunctionTok{from}\OperatorTok{(}\StringTok{"hello"}\OperatorTok{);}
\KeywordTok{let}\NormalTok{ c }\OperatorTok{=}\NormalTok{ s}\OperatorTok{[}\DecValTok{0}\OperatorTok{];}     \CommentTok{// E0277: the type `str` cannot be indexed by `\{integer\}`}
\end{Highlighting}
\end{Shaded}
\item
  Dilimleme (byte aralığı):

\begin{Shaded}
\begin{Highlighting}[]
\OperatorTok{\&}\NormalTok{s}\OperatorTok{[}\DecValTok{0}\OperatorTok{..}\DecValTok{2}\OperatorTok{]}          \CommentTok{// "he"}
\end{Highlighting}
\end{Shaded}


  Hata (runtime): \texttt{\&"é"{[}0..1{]}} → panic \texttt{byte\ index\ 1\ is\ not\ a\ char\ boundary}
\item
  Dolaşma:

\begin{Shaded}
\begin{Highlighting}[]
\ControlFlowTok{for}\NormalTok{ c }\KeywordTok{in}\NormalTok{ s}\OperatorTok{.}\FunctionTok{chars}\OperatorTok{()} \OperatorTok{\{\}}               \CommentTok{// karakterler}
\ControlFlowTok{for}\NormalTok{ b }\KeywordTok{in}\NormalTok{ s}\OperatorTok{.}\FunctionTok{bytes}\OperatorTok{()} \OperatorTok{\{\}}               \CommentTok{// byte'lar}
\ControlFlowTok{for} \OperatorTok{(}\NormalTok{i}\OperatorTok{,}\NormalTok{ c}\OperatorTok{)} \KeywordTok{in}\NormalTok{ s}\OperatorTok{.}\FunctionTok{char\_indices}\OperatorTok{()} \OperatorTok{\{\}}   \CommentTok{// byte konumu + karakter}
\end{Highlighting}
\end{Shaded}
\item
  Dönüşümler:

\begin{Shaded}
\begin{Highlighting}[]
\NormalTok{owned}\OperatorTok{.}\FunctionTok{as\_str}\OperatorTok{()}               \CommentTok{// String → \&str}
\OperatorTok{\&}\NormalTok{owned}\OperatorTok{[..]}                   \CommentTok{// aynı}
\StringTok{"42"}\OperatorTok{.}\NormalTok{parse}\OperatorTok{::\textless{}}\DataTypeTok{i32}\OperatorTok{\textgreater{}()}          \CommentTok{// Result\textless{}i32, \_\textgreater{}}
\DataTypeTok{String}\OperatorTok{::}\FunctionTok{from\_utf8}\OperatorTok{(}\NormalTok{vec}\OperatorTok{)}       \CommentTok{// Result\textless{}String, \_\textgreater{}}
\end{Highlighting}
\end{Shaded}
\end{itemize}

\textbf{GoT bağlantı haritası:}

\begin{verbatim}
str (UTF-8 byte'lar) ──► &str (ptr, len) ──┬── literal ("..." → binary içinde)
                                           └── slice (&s[a..b])
String (ptr, cap, len) ──► heap ──┬── push_str / push / +=
                                  ├── + / format!
                                  └── &String ──deref coercion──► &str
len() = byte   |   chars() = karakter   ──► s[0] yok (E0277)
\end{verbatim}

\begin{center}\rule{0.5\linewidth}{0.5pt}\end{center}

\subsection{🟣 KATMAN 3 --- DERIN TEKNIK NOT}\label{katman-3-—-derin-teknik-not-5}

\textbf{Adım 1 --- Bellek Seviyesi}

\begin{itemize}
\tightlist
\item
  Önce soru: ``Veri nerede?'' Cevap:

  \begin{itemize}
  \tightlist
  \item
    \texttt{\&str} stack'te bir fat pointer (geniş pointer) tutar: (ptr, len), 64-bit'te 16 byte. Literal'in metni programın binary dosyasında, salt okunur bölümde durur ve program boyunca yaşar (\texttt{\&'static\ str}).
  \item
    \texttt{String} stack'te (ptr, cap, len) tutar, 64-bit'te 24 byte. Metnin kendisi heap'te bir buffer'dadır.
  \end{itemize}
\item
  \texttt{String::from("hello")} heap'te yer ayırır ve 5 byte'ı kopyalar. Metin büyüdükçe capacity (kapasite) dolunca buffer daha büyük yere taşınır (genelde katlanarak büyür). Bu yüzden \texttt{with\_capacity} ile baştan yer ayırmak hızlıdır.
\item
  \texttt{owned} scope sonunda drop edilir ve heap buffer serbest bırakılır.
\item
  \texttt{\&owned{[}0..5{]}} yeni bir kopya oluşturmaz, aynı buffer'a bakan yeni (ptr, len) üretir.
\item
  Kontrol için: \texttt{std::mem::size\_of::\textless{}String\textgreater{}()} → \texttt{24}, \texttt{std::mem::size\_of::\textless{}\&str\textgreater{}()} → \texttt{16} (64-bit sistem).
\end{itemize}

\textbf{Adım 2 --- Compiler Davranışı}

\begin{itemize}
\tightlist
\item
  Önce soru: ''\texttt{greet(\&owned)} nasıl derleniyor?'' Cevap: \texttt{\&String} beklenen \texttt{\&str}`ye deref coercion (otomatik başvuru dönüşümü) ile çevrilir; \texttt{\&owned} kabaca \texttt{\&owned{[}..{]}} olur.
\item
  \texttt{+} işleci \texttt{String} için \texttt{add(self,\ s:\ \&str)\ -\textgreater{}\ String} olarak tanımlıdır. Bu yüzden soldaki \texttt{String} taşınır (move), sağdaki \texttt{\&String} ise \texttt{\&str}`ye coercion ile çevrilir.
\item
  Metnin UTF-8 geçerli olması garanti edilir. Dilimleme sırasında sınırın bir karakter başlangıcına denk gelip gelmediği \textbf{runtime'da} kontrol edilir, eşleşmezse panic olur.
\item
  İndexleme bilerek yasaktır: UTF-8'de karakterler 1-4 byte olduğundan \texttt{s{[}n{]}} ne byte'ı ne karakteri güvenle verir ve O(1) olamaz.
\item
  MIR'de \texttt{String::from} ve \texttt{push\_str} fonksiyon çağrıları olarak görünür. Görmek için: \texttt{cargo\ rustc\ --\ --emit=mir}.
\end{itemize}

\textbf{Adım 3 --- Hata Kodları}

\begin{itemize}
\tightlist
\item
  \textbf{E0308:} \texttt{String} ile \texttt{\&str} karıştı. Çözüm: \texttt{String::from(...)}/\texttt{.to\_string()} veya \texttt{\&s}/\texttt{.as\_str()}.
\item
  \textbf{E0277:} \texttt{s{[}0{]}} yazdın. Çözüm: \texttt{s.chars().nth(0)} veya byte aralığıyla \texttt{\&s{[}0..1{]}} (dikkatli).
\item
  \textbf{E0369:} iki \texttt{\&str}`yi \texttt{+} ile topladın. Çözüm: \texttt{format!("\{a\}\{b\}")} veya soldaki \texttt{String} yap.
\item
  \textbf{E0382:} \texttt{s1\ +\ \&s2}`den sonra \texttt{s1}`i kullandın. Çözüm: \texttt{s1.clone()\ +\ \&s2} veya \texttt{format!}.
\item
  \textbf{E0596:} \texttt{mut} olmayan \texttt{String}`e \texttt{push\_str}. Çözüm: \texttt{let\ mut}.
\item
  \textbf{Runtime panic:} \texttt{byte\ index\ 1\ is\ not\ a\ char\ boundary}. Çözüm: \texttt{char\_indices()} veya \texttt{s.get(0..1)} (\texttt{Option\textless{}\&str\textgreater{}} döner).
\item
  Detay için: \texttt{rustc\ --explain\ E0382}.
\end{itemize}

\textbf{Adım 4 --- ``Neden Böyle?''}

\begin{itemize}
\tightlist
\item
  C'de metin \texttt{char*} ve NUL ile biter: uzunluk bilinmez, taşma (buffer overflow) kolaydır. Rust'ta uzunluk her zaman (ptr, len) içinde taşınır.
\item
  C++`taki \texttt{std::string} ve \texttt{std::string\_view} ikilisi \texttt{String} ve \texttt{\&str} ile benzer işi görür. Fark: Rust'ta borrow checker, dilimin sahibinden uzun yaşamasını derleme zamanında engeller.
\item
  UTF-8 zorunluluğu her yerde aynı kodlamayı garanti eder; karşılığında karakter bazlı O(1) indexleme yoktur.
\item
  Trade-off: iki tip öğrenmek zorundasın, ama kim sahip, kim ödünç aldı sorusu tipte açıkça yazar.
\end{itemize}

\textbf{Adım 5 --- Rust 1.96.0 Farkı}

\begin{itemize}
\tightlist
\item
  Rust 1.96.0'daki değişiklikler (yeni Range* tipleri, \texttt{assert\_matches!} makrosu, Wasm linker güncellemeleri) \texttt{String}/\texttt{\&str}, UTF-8 veya deref coercion kurallarını değiştirmiyor.
\item
  Bu üç 1.96.0 değişikliği string slice veya \texttt{String} kullanımını etkilemiyor.
\end{itemize}

\begin{center}\rule{0.5\linewidth}{0.5pt}\end{center}

\subsection{🔴 KATMAN 4 --- SIK UNUTULANLAR + DERINLEŞTIRME}\label{katman-4-—-sik-unutulanlar-derinleştirme-5}

\textbf{Sık Unutulan 1: Ne zaman \texttt{String}, ne zaman \texttt{\&str}?}

\begin{itemize}
\tightlist
\item
  Neden unutuluyor: ikisi de ``metin'' gibi görünüyor.
\item
  Yol 1 (analoji): \texttt{String} = sahip olduğun defter, \texttt{\&str} = defterin bir sayfasına bakmak.
\item
  Yol 2 (kod): parametre \texttt{fn\ f(s:\ \&str)}, saklayan alan \texttt{struct\ User\ \{\ name:\ String\ \}}.
\item
  Yol 3 (hata): \texttt{f(owned)} → \texttt{E0308:\ expected\ \&str,\ found\ String}
\end{itemize}

\textbf{Sık Unutulan 2: \texttt{s{[}0{]}} yok}

\begin{itemize}
\tightlist
\item
  Neden unutuluyor: Python ve JavaScript'te çalışır.
\item
  Yol 1 (görsel): \texttt{é} iki byte, \texttt{😀} dört byte; ``3. karakter'' hangi byte?
\item
  Yol 2 (kod): \texttt{s.chars().nth(0)} → \texttt{Option\textless{}char\textgreater{}}
\item
  Yol 3 (hata): \texttt{E0277:\ the\ type\ str\ cannot\ be\ indexed\ by\ \{integer\}}
\end{itemize}

\textbf{Sık Unutulan 3: \texttt{+} soldakini taşır}

\begin{itemize}
\tightlist
\item
  Neden unutuluyor: sayılarda \texttt{+} hiçbir şeyi taşımaz.
\item
  Yol 1 (analoji): \texttt{a\ +\ \&b} = a'nın defterine b'yi ekleyip defteri yeni isme devretmek.
\item
  Yol 2 (kod): taşımadan birleştirmek için \texttt{format!("\{a\}\{b\}")}
\item
  Yol 3 (hata): \texttt{E0382:\ borrow\ of\ moved\ value:\ a}
\end{itemize}

\textbf{Hazır Derinleştirme Prompt'u:}

\begin{verbatim}
Rust 1.96.0'da String vs &str konusunu en detaylı şekilde anlat. Şunları kapsa:

1. BELLEK SEVİYESİ: Stack'te ne var, heap'te ne var? Kaç byte? String'in
   (ptr, cap, len) yapısı ve &str'nin fat pointer yapısı. String büyürken buffer
   nasıl yeniden ayrılıyor? Literal'ler nerede saklanıyor?

2. COMPILER DAVRANIŞI: Deref coercion nasıl çalışıyor? + işleci neden soldaki
   String'i taşıyor? Dilimlemede char boundary kontrolü nasıl yapılıyor?
   MIR çıktısında String işlemleri nasıl görünür? (--emit=mir)

3. HATA KODLARI: E0308, E0277, E0369, E0382, E0596. Her biri için: ne zaman,
   neden, nasıl çözülür. `rustc --explain EXXXX` çıktısını da ekle.

4. C/C++ KARŞILAŞTIRMASI: char*, std::string, std::string_view ile karşılaştır.
   Rust neden UTF-8 zorunlu kılıyor ve neden indexleme yok?

5. ANALOJİ: Günlük hayattan bir benzetme yap (örn: defter ve sayfa, kitap ve
   fotokopi).

6. RUST 1.96.0: Bu konuda son sürümde değişen bir şey var mı?

7. PRATİK: 3 tane kod örneği ver — kolay, orta, zor.

Açıklama Türkçe, teknik terimler İngilizce, kod yorumları Türkçe olsun.
Her bölümü ayrı başlıkla ver. Kod bloklarını çalıştırılabilir yaz.
\end{verbatim}

\begin{center}\rule{0.5\linewidth}{0.5pt}\end{center}

\subsection{📝 GÖREV (Boş dosyaya yaz)}\label{görev-boş-dosyaya-yaz-4}

\textbf{🟢 Kolay}

\begin{enumerate}
\def\labelenumi{\arabic{enumi}.}
\tightlist
\item
  Bir \texttt{\&str} literal ve bir \texttt{String} oluştur, ikisini de yazdır.
\item
  \texttt{let\ mut\ s\ =\ String::from("Rust");} yap; \texttt{push\_str("\ is\ fun")} ve \texttt{push('!')} ekle, yazdır.
\end{enumerate}

\textbf{🟡 Orta}
3. \texttt{fn\ shout(s:\ \&str)\ -\textgreater{}\ String} yaz (\texttt{s.to\_uppercase()} döndürsün). Hem \texttt{"hello"} literal'i hem \texttt{\&String} ile çağır.
4. \texttt{"héllo"} için \texttt{len()} ve \texttt{chars().count()} değerlerini yazdır.

\textbf{🔴 Zor}
5. \texttt{let\ a\ =\ String::from("Hello");\ let\ b\ =\ String::from("\ Rust");\ let\ c\ =\ a\ +\ \&b;} yaz ve yazdır. Sonra \texttt{a}`yı yazdırmayı dene (\texttt{E0382}`yi gör), ardından \texttt{format!} ile taşımadan birleştir. Son olarak \texttt{\&c{[}0..5{]}} ile dilimle ve \texttt{"é"} içeren bir metinde sınırı bozan dilimlemeyle panic'i gör.

\textbf{Başarı kriteri:} Hepsini hatasız yazarsan → öğrendin ✅

\subsubsection{✅ Kendini Test Et (2-3 saat sonra, sıfırdan)}\label{kendini-test-et-2-3-saat-sonra-sıfırdan-5}

\begin{enumerate}
\def\labelenumi{\arabic{enumi}.}
\tightlist
\item
  \texttt{String} oluştur, büyüt ve \texttt{\&str} parametreli fonksiyona ver.
\item
  \texttt{len()} ile \texttt{chars().count()} farkını göster.
\item
  \texttt{+} ile \texttt{format!} farkını (move) göster.
\end{enumerate}

\begin{center}\rule{0.5\linewidth}{0.5pt}\end{center}

\subsection{📚 KAYNAKÇA}\label{kaynakça-5}

\subsubsection{Resmi Kaynaklar}\label{resmi-kaynaklar-5}

\begin{itemize}
\tightlist
\item
  \href{https://doc.rust-lang.org/book/ch08-02-strings.html}{Rust Book --- Storing UTF-8 Encoded Text with Strings}
\item
  \href{https://doc.rust-lang.org/book/ch04-03-slices.html}{Rust Book --- The Slice Type}
\item
  \href{https://doc.rust-lang.org/std/string/struct.String.html}{std --- String}
\item
  \href{https://doc.rust-lang.org/std/primitive.str.html}{std --- str}
\end{itemize}

\subsubsection{Sürüm Notları}\label{sürüm-notları-5}

\begin{itemize}
\tightlist
\item
  \href{https://blog.rust-lang.org/2026/05/28/Rust-1.96.0/}{Rust Blog --- Announcing 1.96.0}
\item
  \href{https://doc.rust-lang.org/stable/releases.html\#version-1960-2026-05-28}{Rust 1.96.0 Release Notes}
\end{itemize}

\subsubsection{Hata Kodları}\label{hata-kodları-5}

\begin{itemize}
\tightlist
\item
  \href{https://doc.rust-lang.org/error_codes/E0308.html}{E0308}
\item
  \href{https://doc.rust-lang.org/error_codes/E0277.html}{E0277}
\item
  \href{https://doc.rust-lang.org/error_codes/E0369.html}{E0369}
\item
  \href{https://doc.rust-lang.org/error_codes/E0382.html}{E0382}
\item
  \href{https://doc.rust-lang.org/error_codes/E0596.html}{E0596}
\end{itemize}

\subsubsection{Ek Kaynaklar}\label{ek-kaynaklar-5}

\begin{itemize}
\tightlist
\item
  \href{https://doc.rust-lang.org/rust-by-example/std/str.html}{Rust By Example --- Strings}
\item
  \href{https://moodle.cs.pdx.edu/mod/page/view.php?id=212}{PSU CS notları --- Strings and I/O}
\end{itemize}

\emph{Not: \texttt{String} ve \texttt{\&str} tanımları, \texttt{add(self,\ s:\ \&str)} imzası, deref coercion ve \texttt{E0277} indexleme hatası aramada Rust Book üzerinden doğrulandı. Boyutlar (24 ve 16 byte), büyüme stratejisi ve \texttt{char\ boundary} panic mesajı bilgime dayanıyor; \texttt{size\_of} ve panic örneğini kendi makinende çalıştırarak kontrol et.}

\begin{center}\rule{0.5\linewidth}{0.5pt}\end{center}

\section{📘 KONU 7: if / else if / else}\label{konu-7-if-else-if-else}

\textbf{Kategori:} Kontrol \textbar{} \textbf{Sürüm:} Rust 1.96.0

\begin{center}\rule{0.5\linewidth}{0.5pt}\end{center}

\subsection{🟢 KATMAN 1 --- ILK OKUMA}\label{katman-1-—-ilk-okuma-6}

\begin{Shaded}
\begin{Highlighting}[]
\KeywordTok{fn} \FunctionTok{main}\OperatorTok{()} \OperatorTok{\{}
    \KeywordTok{let}\NormalTok{ number }\OperatorTok{=} \DecValTok{6}\OperatorTok{;}

    \ControlFlowTok{if}\NormalTok{ number }\OperatorTok{\%} \DecValTok{4} \OperatorTok{==} \DecValTok{0} \OperatorTok{\{}
\NormalTok{        println}\OperatorTok{!(}\StringTok{"divisible by 4"}\OperatorTok{);}
    \OperatorTok{\}} \ControlFlowTok{else} \ControlFlowTok{if}\NormalTok{ number }\OperatorTok{\%} \DecValTok{3} \OperatorTok{==} \DecValTok{0} \OperatorTok{\{}
\NormalTok{        println}\OperatorTok{!(}\StringTok{"divisible by 3"}\OperatorTok{);}
    \OperatorTok{\}} \ControlFlowTok{else} \ControlFlowTok{if}\NormalTok{ number }\OperatorTok{\%} \DecValTok{2} \OperatorTok{==} \DecValTok{0} \OperatorTok{\{}
\NormalTok{        println}\OperatorTok{!(}\StringTok{"divisible by 2"}\OperatorTok{);}
    \OperatorTok{\}} \ControlFlowTok{else} \OperatorTok{\{}
\NormalTok{        println}\OperatorTok{!(}\StringTok{"not divisible by 4, 3, or 2"}\OperatorTok{);}
    \OperatorTok{\}}

    \KeywordTok{let}\NormalTok{ is\_big }\OperatorTok{=}\NormalTok{ number }\OperatorTok{\textgreater{}} \DecValTok{5}\OperatorTok{;}
    \KeywordTok{let}\NormalTok{ label }\OperatorTok{=} \ControlFlowTok{if}\NormalTok{ is\_big }\OperatorTok{\{} \StringTok{"big"} \OperatorTok{\}} \ControlFlowTok{else} \OperatorTok{\{} \StringTok{"small"} \OperatorTok{\};}  \CommentTok{// if bir ifade}
\NormalTok{    println}\OperatorTok{!(}\StringTok{"\{label\}"}\OperatorTok{);}

    \KeywordTok{let}\NormalTok{ temp }\OperatorTok{=} \DecValTok{25}\OperatorTok{;}
    \ControlFlowTok{if}\NormalTok{ temp }\OperatorTok{\textgreater{}} \DecValTok{20} \OperatorTok{\&\&}\NormalTok{ temp }\OperatorTok{\textless{}} \DecValTok{30} \OperatorTok{\{}                        \CommentTok{// iki koşul birlikte}
\NormalTok{        println}\OperatorTok{!(}\StringTok{"warm"}\OperatorTok{);}
    \OperatorTok{\}}
\OperatorTok{\}}
\CommentTok{// Çıktı:}
\CommentTok{// divisible by 3}
\CommentTok{// big}
\CommentTok{// warm}
\end{Highlighting}
\end{Shaded}

\texttt{if} bir condition (koşul) alır ve bu koşul mutlaka bool (mantıksal değer) olmalıdır; sayıyı otomatik \texttt{true} ya da \texttt{false} yapmaz. Zincirdeki arm'lardan (kol) yalnızca ilk \texttt{true} olan çalışır, kalanlara bakılmaz. \texttt{if} bir expression (ifade) olduğu için değer üretir ve \texttt{let\ label\ =\ if\ ...\ \{\ ...\ \}\ else\ \{\ ...\ \};} biçiminde kullanılabilir. Bu durumda her iki kolun da aynı tipte değer üretmesi gerekir. Koşulun etrafına parantez yazmak zorunda değilsin, ama süslü parantez her zaman zorunludur.

\textbf{1.96.0 notu:} \texttt{if} kuralları 1.96.0'da aynı.

\begin{center}\rule{0.5\linewidth}{0.5pt}\end{center}

\subsection{🔵 KATMAN 2 --- RECALL KARTLARI}\label{katman-2-—-recall-kartlari-6}

\begin{itemize}
\item
  Temel yapı:

\begin{Shaded}
\begin{Highlighting}[]
\ControlFlowTok{if}\NormalTok{ cond }\OperatorTok{\{} \OperatorTok{...} \OperatorTok{\}} \ControlFlowTok{else} \ControlFlowTok{if}\NormalTok{ cond2 }\OperatorTok{\{} \OperatorTok{...} \OperatorTok{\}} \ControlFlowTok{else} \OperatorTok{\{} \OperatorTok{...} \OperatorTok{\}}
\end{Highlighting}
\end{Shaded}


  İlgili: \texttt{else} zorunlu değil
\item
  Sadece ilk \texttt{true} kol çalışır, sıra önemlidir:

\begin{Shaded}
\begin{Highlighting}[]
\KeywordTok{let}\NormalTok{ n }\OperatorTok{=} \DecValTok{12}\OperatorTok{;}
\ControlFlowTok{if}\NormalTok{ n }\OperatorTok{\%} \DecValTok{2} \OperatorTok{==} \DecValTok{0} \OperatorTok{\{} \OperatorTok{\}}      \CommentTok{// burada durur}
\ControlFlowTok{else} \ControlFlowTok{if}\NormalTok{ n }\OperatorTok{\%} \DecValTok{4} \OperatorTok{==} \DecValTok{0} \OperatorTok{\{} \OperatorTok{\}} \CommentTok{// asla çalışmaz}
\end{Highlighting}
\end{Shaded}
\item
  Koşul \texttt{bool} olmalı:

\begin{Shaded}
\begin{Highlighting}[]
\KeywordTok{let}\NormalTok{ n }\OperatorTok{=} \DecValTok{3}\OperatorTok{;}
\ControlFlowTok{if}\NormalTok{ n }\OperatorTok{\{} \OperatorTok{\}}         \CommentTok{// E0308: expected `bool`, found integer}
\ControlFlowTok{if}\NormalTok{ n }\OperatorTok{!=} \DecValTok{0} \OperatorTok{\{} \OperatorTok{\}}    \CommentTok{// doğrusu}
\end{Highlighting}
\end{Shaded}
\item
  Süslü parantez zorunlu, normal parantez gereksiz:

\begin{Shaded}
\begin{Highlighting}[]
\ControlFlowTok{if} \OperatorTok{(}\NormalTok{n }\OperatorTok{\textgreater{}} \DecValTok{5}\OperatorTok{)} \OperatorTok{\{} \OperatorTok{\}}   \CommentTok{// warning: unnecessary parentheses}
\ControlFlowTok{if}\NormalTok{ n }\OperatorTok{\textgreater{}} \DecValTok{5}\NormalTok{ println}\OperatorTok{!(}\StringTok{"x"}\OperatorTok{);}   \CommentTok{// error: expected `\{`, found `println`}
\end{Highlighting}
\end{Shaded}
\item
  \texttt{if} ifade olarak değer döndürür (kolların sonunda \texttt{;} yok):

\begin{Shaded}
\begin{Highlighting}[]
\KeywordTok{let}\NormalTok{ x }\OperatorTok{=} \ControlFlowTok{if}\NormalTok{ cond }\OperatorTok{\{} \DecValTok{5} \OperatorTok{\}} \ControlFlowTok{else} \OperatorTok{\{} \DecValTok{6} \OperatorTok{\};}
\end{Highlighting}
\end{Shaded}


  İlgili: ternary operatörü yok
\item
  Kollar aynı tipte olmalı:

\begin{Shaded}
\begin{Highlighting}[]
\KeywordTok{let}\NormalTok{ x }\OperatorTok{=} \ControlFlowTok{if}\NormalTok{ cond }\OperatorTok{\{} \DecValTok{5} \OperatorTok{\}} \ControlFlowTok{else} \OperatorTok{\{} \StringTok{"six"} \OperatorTok{\};}
\CommentTok{// E0308: `if` and `else` have incompatible types}
\end{Highlighting}
\end{Shaded}
\item
  Değer alıyorsan \texttt{else} zorunlu:

\begin{Shaded}
\begin{Highlighting}[]
\KeywordTok{let}\NormalTok{ x }\OperatorTok{=} \ControlFlowTok{if}\NormalTok{ cond }\OperatorTok{\{} \DecValTok{5} \OperatorTok{\};}
\CommentTok{// E0317: `if` may be missing an `else` clause}
\end{Highlighting}
\end{Shaded}
\item
  Kolun sonuna \texttt{;} koyma tuzağı:

\begin{Shaded}
\begin{Highlighting}[]
\KeywordTok{let}\NormalTok{ x }\OperatorTok{=} \ControlFlowTok{if}\NormalTok{ cond }\OperatorTok{\{} \DecValTok{5}\OperatorTok{;} \OperatorTok{\}} \ControlFlowTok{else} \OperatorTok{\{} \DecValTok{6} \OperatorTok{\};}
\CommentTok{// E0308: mismatched types (ilk kol () döndürdü)}
\end{Highlighting}
\end{Shaded}
\item
  Birleşik koşul: \texttt{\&\&}, \texttt{\textbar{}\textbar{}}, \texttt{!}

\begin{Shaded}
\begin{Highlighting}[]
\ControlFlowTok{if}\NormalTok{ a }\OperatorTok{\textgreater{}} \DecValTok{0} \OperatorTok{\&\&} \OperatorTok{(}\NormalTok{b }\OperatorTok{\textgreater{}} \DecValTok{0} \OperatorTok{||}\NormalTok{ c }\OperatorTok{\textgreater{}} \DecValTok{0}\OperatorTok{)} \OperatorTok{\{} \OperatorTok{\}}
\end{Highlighting}
\end{Shaded}


  İlgili: kısa devre (Konu 5)
\item
  Çok \texttt{else\ if} varsa \texttt{match} kullan (Konu 8); desen eşleştirme için \texttt{if\ let} (Konu 18, 20)
\end{itemize}

\textbf{GoT bağlantı haritası:}

\begin{verbatim}
if ──┬── koşul: bool ──► && || ! ──► E0308 (bool değilse)
     ├── else if ──► ilk true olan çalışır
     ├── else ──► değer alıyorsan zorunlu (E0317)
     └── ifade ──► let x = if ... ──► kollar aynı tip (E0308)
çok kol ──► match (Konu 8) ; desen ──► if let (Konu 18, 20)
\end{verbatim}

\begin{center}\rule{0.5\linewidth}{0.5pt}\end{center}

\subsection{🟣 KATMAN 3 --- DERIN TEKNIK NOT}\label{katman-3-—-derin-teknik-not-6}

\textbf{Adım 1 --- Bellek Seviyesi}

\begin{itemize}
\tightlist
\item
  Önce soru: ``Bellekte ne oluyor?'' Cevap: \texttt{if} için ayrı bir veri yapısı yoktur. Koşul bir karşılaştırma (\texttt{cmp}) komutuna ve ardından koşullu atlamaya (\texttt{jne}, \texttt{jle} gibi) dönüşür.
\item
  \texttt{let\ x\ =\ if\ c\ \{\ 5\ \}\ else\ \{\ 6\ \};} ifadesinde sonuç tek bir yerde (register ya da stack slot) tutulur. Her iki kolun tipi aynı olmak zorunda olmasının sebebi budur: \texttt{x}`in boyutu compile time'da bilinmelidir.
\item
  Çalışmayan kolun kodu yürütülmez; bu yüzden yan etkileri de (örneğin \texttt{println!}) görünmez.
\item
  Basit durumlarda derleyici dal yerine \texttt{cmov} gibi koşullu seçim komutu üretebilir.
\end{itemize}

\textbf{Adım 2 --- Compiler Davranışı}

\begin{itemize}
\tightlist
\item
  Önce soru: ``Compiler ne kontrol ediyor?'' Cevap: koşulun tipinin \texttt{bool} olması ve kolların tip birleştirmesi (unification).
\item
  Değer kullanılmadığında \texttt{if} kolları \texttt{()} döndürür. Bu yüzden \texttt{else}`siz \texttt{if}`in gövdesi \texttt{()} olmak zorundadır.
\item
  \texttt{return}, \texttt{break}, \texttt{panic!} gibi hiç dönmeyen ifadelerin tipi \texttt{!}`dir ve herhangi bir tipe uyar:

\begin{Shaded}
\begin{Highlighting}[]
\KeywordTok{let}\NormalTok{ x }\OperatorTok{=} \ControlFlowTok{if}\NormalTok{ c }\OperatorTok{\{} \DecValTok{5} \OperatorTok{\}} \ControlFlowTok{else} \OperatorTok{\{} \ControlFlowTok{return}\OperatorTok{;} \OperatorTok{\};}
\end{Highlighting}
\end{Shaded}
\item
  Sabit koşullarda (\texttt{if\ true}) ölü kod elenir.
\item
  MIR'de \texttt{if} bir \texttt{switchInt} olarak görünür. Görmek için: \texttt{cargo\ rustc\ --\ --emit=mir}.
\end{itemize}

\textbf{Adım 3 --- Hata Kodları}

\begin{itemize}
\tightlist
\item
  \textbf{E0308 (bool bekleniyor):} koşul \texttt{bool} değil. Çözüm: \texttt{n\ !=\ 0} gibi karşılaştırma yaz.
\item
  \textbf{E0308 (\texttt{if} and \texttt{else} have incompatible types):} kollar farklı tipte. Çözüm: ikisini aynı tipe getir.
\item
  \textbf{E0308 (mismatched types, \texttt{;} yüzünden):} kolun sonundaki \texttt{;} değeri \texttt{()} yaptı. Çözüm: \texttt{;}`yi kaldır.
\item
  \textbf{E0317:} değer üreten \texttt{if}`te \texttt{else} yok. Çözüm: \texttt{else} ekle.
\item
  \textbf{Kodsuz hata:} \texttt{expected\ \{,\ found\ ...} (süslü parantez eksik)
\item
  \textbf{Uyarı:} \texttt{unnecessary\ parentheses\ around\ if\ condition}
\item
  Detay için: \texttt{rustc\ --explain\ E0317}.
\end{itemize}

\textbf{Adım 4 --- ``Neden Böyle?''}

\begin{itemize}
\tightlist
\item
  C/C++`ta \texttt{if\ (x)} tam sayıyla çalışır ve \texttt{if\ (x\ =\ 5)} geçerlidir. Rust bunu \texttt{bool} zorunluluğuyla engeller.
\item
  C'de süslü parantezsiz \texttt{if} yazılabilir; yanlış girinti kritik güvenlik hatalarına yol açmıştır. Rust'ta süslü parantez zorunludur.
\item
  C'deki \texttt{cond\ ?\ a\ :\ b} yerine Rust'ta \texttt{if} zaten bir ifadedir; ayrı bir ternary operatörüne gerek yoktur.
\item
  Trade-off: biraz daha uzun yazarsın, ama anlam netleşir.
\end{itemize}

\textbf{Adım 5 --- Rust 1.96.0 Farkı}

\begin{itemize}
\tightlist
\item
  Rust 1.96.0'daki değişiklikler (yeni Range* tipleri, \texttt{assert\_matches!} makrosu, Wasm linker güncellemeleri) temel \texttt{if} / \texttt{else\ if} / \texttt{else} davranışını değiştirmiyor.
\item
  Yakın dönemde \texttt{if} ile ilgili asıl yenilik let chains'tir (\texttt{if\ let\ Some(x)\ =\ a\ \&\&\ x\ \textgreater{}\ 5}). Bu özelliğin Rust 1.88'de edition 2024 için stabilize edildiğini hatırlıyorum, ama bu oturumda doğrulamadım. Henüz öğrenmediğin bir konu olduğu için şimdilik atla.
\item
  Bu değişiklikler \texttt{if} koşullarının değer üretme ve \texttt{bool} değerlendirme kurallarını etkilemiyor.
\end{itemize}

\begin{center}\rule{0.5\linewidth}{0.5pt}\end{center}

\subsection{🔴 KATMAN 4 --- SIK UNUTULANLAR + DERINLEŞTIRME}\label{katman-4-—-sik-unutulanlar-derinleştirme-6}

\textbf{Sık Unutulan 1: Koşul \texttt{bool} olmak zorunda}

\begin{itemize}
\tightlist
\item
  Neden unutuluyor: C, Python, JavaScript'te sayı veya metin koşul olabilir.
\item
  Yol 1 (analoji): Rust'ta kapıcı sadece ``evet'' ya da ``hayır'' cevabını kabul eder, ``3'' demek yetmez.
\item
  Yol 2 (kod): \texttt{if\ n\ !=\ 0\ \{\ \}}
\item
  Yol 3 (hata): \texttt{if\ n\ \{\ \}} → \texttt{E0308:\ expected\ bool,\ found\ integer}
\end{itemize}

\textbf{Sık Unutulan 2: \texttt{if} ifade olarak kullanılırken \texttt{;} ve tip}

\begin{itemize}
\tightlist
\item
  Neden unutuluyor: kolun son satırına alışkanlıkla \texttt{;} koyarsın.
\item
  Yol 1 (görsel): son satırda noktalı virgül = ``değer üretme'' demek.
\item
  Yol 2 (kod): \texttt{let\ x\ =\ if\ c\ \{\ 5\ \}\ else\ \{\ 6\ \};} ve kolların içinde \texttt{5} yazıyor, \texttt{5;} değil.
\item
  Yol 3 (hata): \texttt{E0308} ve \texttt{E0317}
\end{itemize}

\textbf{Sık Unutulan 3: Sadece ilk \texttt{true} kol çalışır}

\begin{itemize}
\tightlist
\item
  Neden unutuluyor: tüm koşulların kontrol edildiğini sanırsın.
\item
  Yol 1 (analoji): bankodaki ilk boş gişeye gidersin, diğerlerine bakmazsın.
\item
  Yol 2 (kod): \texttt{n\ =\ 12} için \texttt{n\ \%\ 2\ ==\ 0} önce yazılırsa \texttt{n\ \%\ 4\ ==\ 0} kolu asla çalışmaz; özel durumu başa yaz.
\item
  Yol 3 (hata): derleyici hata vermez; sonuç mantıksal olarak yanlış çıkar.
\end{itemize}

\textbf{Hazır Derinleştirme Prompt'u:}

\begin{verbatim}
Rust 1.96.0'da if / else if / else konusunu en detaylı şekilde anlat. Şunları kapsa:

1. BELLEK SEVİYESİ: if assembly'de nasıl görünür? let x = if ... { } else { }
   ifadesinde sonuç nerede tutulur? Neden kolların tipi aynı olmalı?

2. COMPILER DAVRANIŞI: Koşulun bool kontrolü ve kolların tip birleştirmesi nasıl
   yapılıyor? Never type (!) kolu nasıl uyuyor? MIR'de switchInt nasıl görünür?
   (--emit=mir)

3. HATA KODLARI: E0308 (üç farklı durum) ve E0317. Her biri için: ne zaman, neden,
   nasıl çözülür. `rustc --explain EXXXX` çıktısını da ekle.

4. C/C++ KARŞILAŞTIRMASI: if (x), ternary operatörü, süslü parantezsiz if.
   Rust neden farklı?

5. ANALOJİ: Günlük hayattan bir benzetme yap (örn: yol ayrımı, gişe kuyruğu).

6. RUST 1.96.0: Bu konuda son sürümde değişen bir şey var mı? (let chains dahil)

7. PRATİK: 3 tane kod örneği ver — kolay, orta, zor.

Açıklama Türkçe, teknik terimler İngilizce, kod yorumları Türkçe olsun.
Her bölümü ayrı başlıkla ver. Kod bloklarını çalıştırılabilir yaz.
\end{verbatim}

\begin{center}\rule{0.5\linewidth}{0.5pt}\end{center}

\subsection{📝 GÖREV (Boş dosyaya yaz)}\label{görev-boş-dosyaya-yaz-5}

\textbf{🟢 Kolay}

\begin{enumerate}
\def\labelenumi{\arabic{enumi}.}
\tightlist
\item
  \texttt{let\ number\ =\ 7;} için \texttt{number\ \textgreater{}\ 5} ise ``big'', değilse ``small'' yazdır (\texttt{if} / \texttt{else}).
\item
  \texttt{score} değişkenine göre not yazdır: \texttt{\textgreater{}=\ 90} → ``A'', \texttt{\textgreater{}=\ 80} → ``B'', diğerleri → ``C'' (\texttt{else\ if} zinciri).
\end{enumerate}

\textbf{🟡 Orta}
3. \texttt{let\ label\ =\ if\ number\ \textgreater{}\ 5\ \{\ "big"\ \}\ else\ \{\ "small"\ \};} yaz ve \texttt{label}`ı yazdır.
4. Bir sayının çift mi tek mi olduğunu \texttt{\%} ile bulup yazdır.

\textbf{🔴 Zor}
5. \texttt{let\ n\ =\ 15;} için FizzBuzz mantığı yaz (önce \texttt{n\ \%\ 15}, sonra \texttt{n\ \%\ 3}, \texttt{n\ \%\ 5}). Sonra bilerek \texttt{E0317} (\texttt{else}`siz değer alan \texttt{if}) ve \texttt{E0308} (\texttt{if\ \{\ 5\ \}\ else\ \{\ "six"\ \}}) hatalarını üretip mesajları oku.

\textbf{Başarı kriteri:} Hepsini hatasız yazarsan → öğrendin ✅

\subsubsection{✅ Kendini Test Et (2-3 saat sonra, sıfırdan)}\label{kendini-test-et-2-3-saat-sonra-sıfırdan-6}

\begin{enumerate}
\def\labelenumi{\arabic{enumi}.}
\tightlist
\item
  \texttt{if} / \texttt{else\ if} / \texttt{else} zinciri yaz.
\item
  \texttt{let\ x\ =\ if\ ...\ \{\ \}\ else\ \{\ \};} yaz.
\item
  \texttt{if\ 1\ \{\ \}} neden hata verir, söyle.
\end{enumerate}

\begin{center}\rule{0.5\linewidth}{0.5pt}\end{center}

\subsection{📚 KAYNAKÇA}\label{kaynakça-6}

\subsubsection{Resmi Kaynaklar}\label{resmi-kaynaklar-6}

\begin{itemize}
\tightlist
\item
  \href{https://doc.rust-lang.org/book/ch03-05-control-flow.html}{Rust Book --- Control Flow}
\item
  \href{https://doc.rust-lang.org/reference/expressions/if-expr.html}{Rust Reference --- if expressions}
\end{itemize}

\subsubsection{Sürüm Notları}\label{sürüm-notları-6}

\begin{itemize}
\tightlist
\item
  \href{https://blog.rust-lang.org/2026/05/28/Rust-1.96.0/}{Rust Blog --- Announcing 1.96.0}
\item
  \href{https://doc.rust-lang.org/stable/releases.html\#version-1960-2026-05-28}{Rust 1.96.0 Release Notes}
\end{itemize}

\subsubsection{Hata Kodları}\label{hata-kodları-6}

\begin{itemize}
\tightlist
\item
  \href{https://doc.rust-lang.org/error_codes/E0308.html}{E0308}
\item
  \href{https://doc.rust-lang.org/error_codes/E0317.html}{E0317}
\end{itemize}

\subsubsection{Ek Kaynaklar}\label{ek-kaynaklar-6}

\begin{itemize}
\tightlist
\item
  \href{https://doc.rust-lang.org/rust-by-example/flow_control/if_else.html}{Rust By Example --- if/else}
\item
  \href{https://google.github.io/comprehensive-rust/pl/pattern-matching/let-control-flow.html}{Comprehensive Rust --- if let, let else, while let}
\end{itemize}

\emph{Not: Koşulun \texttt{bool} olması, \texttt{else\ if} zincirinde ilk \texttt{true} kolun çalışması, \texttt{if}`in ifade olarak kullanılması ve \texttt{E0308} (``if and else have incompatible types'') aramada Rust Book üzerinden doğrulandı. \texttt{E0317}, \texttt{;} tuzağı, MIR \texttt{switchInt} ve let chains'in 1.88'de stabilize edildiği bilgisi bilgime dayanıyor; hata mesajlarını kendi makinende çalıştırarak kontrol et.}

\begin{center}\rule{0.5\linewidth}{0.5pt}\end{center}

\section{📘 KONU 8: match}\label{konu-8-match}

\textbf{Kategori:} Kontrol \textbar{} \textbf{Sürüm:} Rust 1.96.0

\begin{center}\rule{0.5\linewidth}{0.5pt}\end{center}

\subsection{🟢 KATMAN 1 --- ILK OKUMA}\label{katman-1-—-ilk-okuma-7}

\begin{Shaded}
\begin{Highlighting}[]
\KeywordTok{fn} \FunctionTok{main}\OperatorTok{()} \OperatorTok{\{}
    \KeywordTok{let}\NormalTok{ number }\OperatorTok{=} \DecValTok{3}\OperatorTok{;}

    \ControlFlowTok{match}\NormalTok{ number }\OperatorTok{\{}
        \DecValTok{1} \OperatorTok{=\textgreater{}}\NormalTok{ println}\OperatorTok{!(}\StringTok{"one"}\OperatorTok{),}
        \DecValTok{2} \OperatorTok{|} \DecValTok{3} \OperatorTok{=\textgreater{}}\NormalTok{ println}\OperatorTok{!(}\StringTok{"two or three"}\OperatorTok{),}   \CommentTok{// | = veya}
        \DecValTok{4}\OperatorTok{..=}\DecValTok{9} \OperatorTok{=\textgreater{}}\NormalTok{ println}\OperatorTok{!(}\StringTok{"four to nine"}\OperatorTok{),}   \CommentTok{// aralık (9 dahil)}
\NormalTok{        \_ }\OperatorTok{=\textgreater{}}\NormalTok{ println}\OperatorTok{!(}\StringTok{"something else"}\OperatorTok{),}     \CommentTok{// geri kalan her şey}
    \OperatorTok{\}}

    \KeywordTok{let}\NormalTok{ size }\OperatorTok{=} \ControlFlowTok{match}\NormalTok{ number }\OperatorTok{\{}                \CommentTok{// match bir ifade}
        \DecValTok{0} \OperatorTok{=\textgreater{}} \StringTok{"zero"}\OperatorTok{,}
\NormalTok{        n }\ControlFlowTok{if}\NormalTok{ n }\OperatorTok{\textless{}} \DecValTok{0} \OperatorTok{=\textgreater{}} \StringTok{"negative"}\OperatorTok{,}            \CommentTok{// guard: ek koşul}
        \DecValTok{1}\OperatorTok{..=}\DecValTok{5} \OperatorTok{=\textgreater{}} \StringTok{"small"}\OperatorTok{,}
\NormalTok{        \_ }\OperatorTok{=\textgreater{}} \StringTok{"large"}\OperatorTok{,}
    \OperatorTok{\};}
\NormalTok{    println}\OperatorTok{!(}\StringTok{"\{size\}"}\OperatorTok{);}

    \KeywordTok{let}\NormalTok{ pair }\OperatorTok{=} \OperatorTok{(}\DecValTok{0}\OperatorTok{,} \OperatorTok{-}\DecValTok{2}\OperatorTok{);}
    \ControlFlowTok{match}\NormalTok{ pair }\OperatorTok{\{}
        \OperatorTok{(}\DecValTok{0}\OperatorTok{,}\NormalTok{ y}\OperatorTok{)} \OperatorTok{=\textgreater{}}\NormalTok{ println}\OperatorTok{!(}\StringTok{"x is zero, y = \{y\}"}\OperatorTok{),}  \CommentTok{// tuple parçalama}
        \OperatorTok{(}\NormalTok{x}\OperatorTok{,} \DecValTok{0}\OperatorTok{)} \OperatorTok{=\textgreater{}}\NormalTok{ println}\OperatorTok{!(}\StringTok{"y is zero, x = \{x\}"}\OperatorTok{),}
\NormalTok{        \_ }\OperatorTok{=\textgreater{}}\NormalTok{ println}\OperatorTok{!(}\StringTok{"no zeros"}\OperatorTok{),}
    \OperatorTok{\}}
\OperatorTok{\}}
\CommentTok{// Çıktı:}
\CommentTok{// two or three}
\CommentTok{// small}
\CommentTok{// x is zero, y = -2}
\end{Highlighting}
\end{Shaded}

\texttt{match}, bir değeri sırayla pattern'lerle (desen) karşılaştırır ve eşleşen arm'ın (kol) kodunu çalıştırır. İlk eşleşen kol seçilir ve kalanlara bakılmaz. \texttt{match} exhaustive (eksiksiz) olmak zorundadır: değerin alabileceği her ihtimali karşılamalısın. Geri kalan her şeyi \texttt{\_} wildcard'ı (joker) yakalar. \texttt{match} de \texttt{if} gibi bir expression (ifade) olduğu için \texttt{let\ size\ =\ match\ ...\ \{\ ...\ \};} biçiminde değer üretebilir ve tüm kolların tipi aynı olmalıdır.

\textbf{1.96.0 notu:} \texttt{match} davranışı 1.96.0'da aynı.

\begin{center}\rule{0.5\linewidth}{0.5pt}\end{center}

\subsection{🔵 KATMAN 2 --- RECALL KARTLARI}\label{katman-2-—-recall-kartlari-7}

\begin{itemize}
\item
  Temel yapı, ilk eşleşen kol çalışır:

\begin{Shaded}
\begin{Highlighting}[]
\ControlFlowTok{match}\NormalTok{ value }\OperatorTok{\{}
\NormalTok{    pattern1 }\OperatorTok{=\textgreater{}}\NormalTok{ expr1}\OperatorTok{,}
\NormalTok{    pattern2 }\OperatorTok{=\textgreater{}} \OperatorTok{\{}\NormalTok{ expr\_a}\OperatorTok{;}\NormalTok{ expr\_b }\OperatorTok{\}}   \CommentTok{// birden fazla satır → süslü parantez}
\NormalTok{    \_ }\OperatorTok{=\textgreater{}}\NormalTok{ expr3}\OperatorTok{,}
\OperatorTok{\}}
\end{Highlighting}
\end{Shaded}


  İlgili: kollar \texttt{,} ile ayrılır
\item
  Eksiksizlik, \texttt{\_} joker:

\begin{Shaded}
\begin{Highlighting}[]
\KeywordTok{let}\NormalTok{ n}\OperatorTok{:} \DataTypeTok{u8} \OperatorTok{=} \DecValTok{5}\OperatorTok{;}
\ControlFlowTok{match}\NormalTok{ n }\OperatorTok{\{} \DecValTok{0} \OperatorTok{=\textgreater{}} \OperatorTok{\{\},} \DecValTok{1} \OperatorTok{=\textgreater{}} \OperatorTok{\{\}} \OperatorTok{\}}
\CommentTok{// E0004: non-exhaustive patterns: `2\_u8..=u8::MAX` not covered}
\end{Highlighting}
\end{Shaded}


  Çözüm: sona \texttt{\_\ =\textgreater{}\ \{\}} ekle
\item
  Birden fazla desen: \texttt{1\ \textbar{}\ 2\ \textbar{}\ 3\ =\textgreater{}\ ...}
\item
  Aralık: \texttt{1..=5} (5 dahil), \texttt{'a'..='z'}
  İlgili: \texttt{a..b} (b hariç) desenlerde Rust 1.80'den beri kullanılabilir
\item
  Değeri bir isme bağlama (binding):

\begin{Shaded}
\begin{Highlighting}[]
\ControlFlowTok{match}\NormalTok{ n }\OperatorTok{\{}
\NormalTok{    x }\OperatorTok{@} \DecValTok{1}\OperatorTok{..=}\DecValTok{5} \OperatorTok{=\textgreater{}}\NormalTok{ println}\OperatorTok{!(}\StringTok{"small \{x\}"}\OperatorTok{),}   \CommentTok{// @ ile hem aralık hem isim}
\NormalTok{    other }\OperatorTok{=\textgreater{}}\NormalTok{ println}\OperatorTok{!(}\StringTok{"other \{other\}"}\OperatorTok{),}   \CommentTok{// her şeyi yakalar}
\OperatorTok{\}}
\end{Highlighting}
\end{Shaded}
\item
  Guard (\texttt{if} koşulu):

\begin{Shaded}
\begin{Highlighting}[]
\ControlFlowTok{match}\NormalTok{ n }\OperatorTok{\{}
\NormalTok{    x }\ControlFlowTok{if}\NormalTok{ x }\OperatorTok{\textless{}} \DecValTok{0} \OperatorTok{=\textgreater{}} \StringTok{"negative"}\OperatorTok{,}
    \DecValTok{0} \OperatorTok{=\textgreater{}} \StringTok{"zero"}\OperatorTok{,}
\NormalTok{    \_ }\OperatorTok{=\textgreater{}} \StringTok{"positive"}\OperatorTok{,}
\OperatorTok{\}}
\end{Highlighting}
\end{Shaded}


  Hata: guard'lar eksiksizlik hesabına \textbf{katılmaz}:

\begin{Shaded}
\begin{Highlighting}[]
\ControlFlowTok{match}\NormalTok{ n }\OperatorTok{\{}\NormalTok{ x }\ControlFlowTok{if}\NormalTok{ x }\OperatorTok{\textless{}} \DecValTok{0} \OperatorTok{=\textgreater{}} \DecValTok{1}\OperatorTok{,}\NormalTok{ x }\ControlFlowTok{if}\NormalTok{ x }\OperatorTok{\textgreater{}=} \DecValTok{0} \OperatorTok{=\textgreater{}} \DecValTok{2} \OperatorTok{\}}
\CommentTok{// E0004: non-exhaustive patterns: `\_` not covered}
\end{Highlighting}
\end{Shaded}
\item
  \texttt{match} ifade olarak, tüm kollar aynı tipte:

\begin{Shaded}
\begin{Highlighting}[]
\KeywordTok{let}\NormalTok{ s }\OperatorTok{=} \ControlFlowTok{match}\NormalTok{ n }\OperatorTok{\{} \DecValTok{10} \OperatorTok{=\textgreater{}} \DecValTok{8}\OperatorTok{,}\NormalTok{ \_ }\OperatorTok{=\textgreater{}} \StringTok{"Not ten"} \OperatorTok{\};}
\CommentTok{// E0308: `match` arms have incompatible types}
\end{Highlighting}
\end{Shaded}
\item
  Tuple ve bool eşleştirme:

\begin{Shaded}
\begin{Highlighting}[]
\ControlFlowTok{match} \OperatorTok{(}\NormalTok{a}\OperatorTok{,}\NormalTok{ b}\OperatorTok{)} \OperatorTok{\{} \OperatorTok{(}\ConstantTok{true}\OperatorTok{,} \ConstantTok{true}\OperatorTok{)} \OperatorTok{=\textgreater{}} \OperatorTok{\{\},} \OperatorTok{(}\ConstantTok{true}\OperatorTok{,} \ConstantTok{false}\OperatorTok{)} \OperatorTok{=\textgreater{}} \OperatorTok{\{\},} \OperatorTok{(}\ConstantTok{false}\OperatorTok{,}\NormalTok{ \_}\OperatorTok{)} \OperatorTok{=\textgreater{}} \OperatorTok{\{\}} \OperatorTok{\}}
\end{Highlighting}
\end{Shaded}
\item
  Joker en başa konursa:

\begin{Shaded}
\begin{Highlighting}[]
\ControlFlowTok{match}\NormalTok{ n }\OperatorTok{\{}\NormalTok{ \_ }\OperatorTok{=\textgreater{}} \OperatorTok{\{\},} \DecValTok{1} \OperatorTok{=\textgreater{}} \OperatorTok{\{\}} \OperatorTok{\}}
\CommentTok{// warning: unreachable pattern}
\end{Highlighting}
\end{Shaded}
\item
  Kollarda farklı isimler:

\begin{Shaded}
\begin{Highlighting}[]
\ControlFlowTok{match}\NormalTok{ n }\OperatorTok{\{} \DecValTok{1} \OperatorTok{|}\NormalTok{ m }\OperatorTok{=\textgreater{}} \OperatorTok{\{\}} \OperatorTok{\}}
\CommentTok{// E0408: variable `m` is not bound in all patterns}
\end{Highlighting}
\end{Shaded}
\end{itemize}

\textbf{GoT bağlantı haritası:}

\begin{verbatim}
match ──┬── desenler: literal | aralık | tuple | _ | isim | @
        ├── guard (if) ──► eksiksizliğe sayılmaz
        ├── eksiksiz olmalı ──► E0004 ──► _ ile çöz
        └── ifade ──► kollar aynı tip (E0308)
çok kol ──► if / else if zincirinden daha net
enum ve Option ile birlikte ──► Konu 17, 18, 20
\end{verbatim}

\begin{center}\rule{0.5\linewidth}{0.5pt}\end{center}

\subsection{🟣 KATMAN 3 --- DERIN TEKNIK NOT}\label{katman-3-—-derin-teknik-not-7}

\textbf{Adım 1 --- Bellek Seviyesi}

\begin{itemize}
\tightlist
\item
  Önce soru: ``Bellekte ne oluyor?'' Cevap: \texttt{match} için ayrı bir veri yapısı yoktur. Değerler karşılaştırma ve atlama komutlarına dönüşür.
\item
  Yoğun tam sayı kolları (\texttt{0}, \texttt{1}, \texttt{2}, \ldots) çoğunlukla bir jump table (atlama tablosu) olarak derlenir. Aralıklar karşılaştırma zincirine dönüşür.
\item
  Kazanan kolun sonucu, \texttt{if} ifadesindeki gibi tek bir yerde (register ya da stack slot) tutulur. Bu yüzden tüm kolların tipi aynı olmalıdır.
\item
  Binding'ler değeri kopyalar veya taşır: \texttt{Copy} olan tiplerde (sayılar) kopya olur, \texttt{String} gibi tiplerde değer bağlamaya taşınır. Taşımak istemezsen \texttt{match\ \&s\ \{\ ...\ \}} yaz.
\end{itemize}

\textbf{Adım 2 --- Compiler Davranışı}

\begin{itemize}
\tightlist
\item
  Önce soru: ``Eksiksizliği kim kontrol ediyor?'' Cevap: derleyici, desenleri bir matris gibi analiz eden bir algoritma çalıştırır ve kapsanmayan değer aralıklarını bulur. Bu yüzden hata mesajı eksik olan aralığı da gösterir (\texttt{2\_u8..=u8::MAX}).
\item
  Guard'lar derleyici için opak ifadelerdir; içeriklerine bakılmaz. Bu yüzden \texttt{x\ if\ x\ \textgreater{}=\ 0} gibi bir kol eksiksizliğe katkı yapmaz.
\item
  Kendinden önceki desenlerce zaten kapsanan bir kol için \texttt{unreachable\ pattern} uyarısı verir.
\item
  Kollar tip olarak birleştirilir. \texttt{return}, \texttt{panic!}, \texttt{break} gibi dönmeyen kolların tipi \texttt{!}`dir ve her tipe uyar.
\item
  MIR'de \texttt{match}, bir decision tree (karar ağacı) olarak \texttt{switchInt} bloklarına çevrilir. Görmek için: \texttt{cargo\ rustc\ --\ --emit=mir}.
\end{itemize}

\textbf{Adım 3 --- Hata Kodları}

\begin{itemize}
\tightlist
\item
  \textbf{E0004:} \texttt{match} tüm ihtimalleri kapsamıyor. Çözüm: eksik değerleri yaz veya \texttt{\_\ =\textgreater{}\ ...} ekle. Detay: \texttt{rustc\ --explain\ E0004}.
\item
  \textbf{E0308 (\texttt{match} arms have incompatible types):} kollar farklı tip üretiyor. Çözüm: tipleri eşitle.
\item
  \textbf{E0408:} \texttt{\textbar{}} ile bağlanan değişken her desende yok. Çözüm: aynı isimleri her desende kullan.
\item
  \textbf{Uyarı:} \texttt{unreachable\ pattern}, \texttt{unused\ variable}
\item
  Not: bir enum varyantı eksikse aynı \texttt{E0004} çıkar (Konu 17'de göreceksin).
\end{itemize}

\textbf{Adım 4 --- ``Neden Böyle?''}

\begin{itemize}
\tightlist
\item
  C'nin \texttt{switch}`inde fallthrough (alta düşme) vardır: \texttt{break} unutulursa bir sonraki kol da çalışır. Rust'ta kollar birbirinden bağımsızdır, \texttt{break} yoktur.
\item
  C'de eksik case sessizce geçilir. Rust'ta eksiksizlik zorunlu olduğu için yeni bir değer eklendiğinde derleyici eksik yerleri gösterir.
\item
  \texttt{match} bir ifade olduğundan her kol değer üretir; geçici değişken ya da ayrı atamalara gerek kalmaz.
\item
  Trade-off: \texttt{\_} her şeyi yakaladığı için, yeni bir değer eklendiğinde derleyicinin sana yardım etmesini engeller. Mümkünse değerleri tek tek yaz.
\end{itemize}

\textbf{Adım 5 --- Rust 1.96.0 Farkı}

\begin{itemize}
\tightlist
\item
  Rust 1.96.0'daki değişiklikler içinde stabilize edilen \texttt{assert\_matches!} makrosu desen eşleşmesini kontrol etmek için yeni bir araç sağlıyor; \texttt{match} ve pattern kuralları değişmedi.
\item
  Aralık desenlerinde \texttt{a..b} biçimi 1.80'den beri stabil; yeni bir şey değil.
\item
  Yeni Range* tipleri ve Wasm linker güncellemeleri de \texttt{match} sözdizimi ve pattern kurallarını değiştirmiyor.
\end{itemize}

\begin{center}\rule{0.5\linewidth}{0.5pt}\end{center}

\subsection{🔴 KATMAN 4 --- SIK UNUTULANLAR + DERINLEŞTIRME}\label{katman-4-—-sik-unutulanlar-derinleştirme-7}

\textbf{Sık Unutulan 1: \texttt{\_} kolunu unutmak}

\begin{itemize}
\tightlist
\item
  Neden unutuluyor: \texttt{if} zincirinde \texttt{else} olmadan da yazabilirdin.
\item
  Yol 1 (analoji): \texttt{\_} = ``başka ne gelirse'' kutusu, onsuz sınav sorusu cevapsız kalır.
\item
  Yol 2 (kod): \texttt{match\ n\ \{\ 1\ =\textgreater{}\ ...,\ \_\ =\textgreater{}\ ...\ \}}
\item
  Yol 3 (hata): \texttt{E0004:\ non-exhaustive\ patterns:\ ...\ not\ covered}
\end{itemize}

\textbf{Sık Unutulan 2: Guard eksiksizliğe sayılmaz}

\begin{itemize}
\tightlist
\item
  Neden unutuluyor: \texttt{x\ if\ x\ \textless{}\ 0} ve \texttt{x\ if\ x\ \textgreater{}=\ 0} mantıksal olarak her şeyi kapsıyor gibi görünür.
\item
  Yol 1 (analoji): derleyici kapıdaki ``özel şart'' kağıtlarını okumaz, sadece desene bakar.
\item
  Yol 2 (kod): son kolu guard'sız yaz: \texttt{\_\ =\textgreater{}\ ...} veya \texttt{x\ =\textgreater{}\ ...}
\item
  Yol 3 (hata): \texttt{E0004:\ non-exhaustive\ patterns:\ \_\ not\ covered}
\end{itemize}

\textbf{Sık Unutulan 3: Kollar aynı tipte olmalı, kol sonları \texttt{,}}

\begin{itemize}
\tightlist
\item
  Neden unutuluyor: \texttt{if} ile aynı kural olsa da \texttt{match}`te kol sayısı fazla olunca gözden kaçar.
\item
  Yol 1 (görsel): her \texttt{=\textgreater{}} sonrası aynı tip ve sonunda virgül.
\item
  Yol 2 (kod): \texttt{let\ s\ =\ match\ n\ \{\ 0\ =\textgreater{}\ "zero",\ \_\ =\textgreater{}\ "other"\ \};}
\item
  Yol 3 (hata): \texttt{E0308:\ match\ arms\ have\ incompatible\ types}
\end{itemize}

\textbf{Hazır Derinleştirme Prompt'u:}

\begin{verbatim}
Rust 1.96.0'da match konusunu en detaylı şekilde anlat. Şunları kapsa:

1. BELLEK SEVİYESİ: match assembly'de nasıl görünür (jump table, karşılaştırma zinciri)?
   Binding'ler değeri kopyalar mı taşır mı? match &s ile match s farkı nedir?

2. COMPILER DAVRANIŞI: Eksiksizlik kontrolü nasıl çalışıyor? Guard'lar neden
   hesaba katılmıyor? unreachable pattern uyarısı. MIR'de decision tree ve
   switchInt nasıl görünür? (--emit=mir)

3. HATA KODLARI: E0004, E0308 (match arms), E0408. Her biri için: ne zaman, neden,
   nasıl çözülür. `rustc --explain EXXXX` çıktısını da ekle.

4. C/C++ KARŞILAŞTIRMASI: switch, break ve fallthrough. Rust neden eksiksizlik
   zorunlu kılıyor?

5. ANALOJİ: Günlük hayattan bir benzetme yap (örn: posta ayırma, yol ayrımı levhası).

6. RUST 1.96.0: Bu konuda son sürümde değişen bir şey var mı?

7. PRATİK: 3 tane kod örneği ver — kolay, orta, zor.

Açıklama Türkçe, teknik terimler İngilizce, kod yorumları Türkçe olsun.
Her bölümü ayrı başlıkla ver. Kod bloklarını çalıştırılabilir yaz.
\end{verbatim}

\begin{center}\rule{0.5\linewidth}{0.5pt}\end{center}

\subsection{📝 GÖREV (Boş dosyaya yaz)}\label{görev-boş-dosyaya-yaz-6}

\textbf{🟢 Kolay}

\begin{enumerate}
\def\labelenumi{\arabic{enumi}.}
\tightlist
\item
  \texttt{let\ day\ =\ 3;} için \texttt{1}, \texttt{2}, \texttt{3} kollarını ve \texttt{\_} kolunu yazıp bir mesaj yazdır.
\item
  \texttt{1\ \textbar{}\ 2} ve \texttt{4..=9} desenlerini kullanan bir \texttt{match} yaz.
\end{enumerate}

\textbf{🟡 Orta}
3. \texttt{let\ size\ =\ match\ n\ \{\ 0\ =\textgreater{}\ "zero",\ 1..=5\ =\textgreater{}\ "small",\ \_\ =\textgreater{}\ "large"\ \};} yaz ve \texttt{size}`ı yazdır.
4. Guard ile negatif, sıfır ve pozitif ayrımı yap: \texttt{x\ if\ x\ \textless{}\ 0}, \texttt{0}, \texttt{\_}.

\textbf{🔴 Zor}
5. \texttt{let\ pair\ =\ (0,\ -2);} için \texttt{(0,\ y)}, \texttt{(x,\ 0)}, \texttt{\_} kollarını yaz. Sonra bilerek \texttt{E0004} (eksik \texttt{\_}) ve \texttt{E0308} (\texttt{match} arms have incompatible types) hatalarını üretip mesajları oku. Son olarak \texttt{x\ @\ 1..=5} desenini kullan.

\textbf{Başarı kriteri:} Hepsini hatasız yazarsan → öğrendin ✅

\subsubsection{✅ Kendini Test Et (2-3 saat sonra, sıfırdan)}\label{kendini-test-et-2-3-saat-sonra-sıfırdan-7}

\begin{enumerate}
\def\labelenumi{\arabic{enumi}.}
\tightlist
\item
  \texttt{\_} kolu olan bir \texttt{match} yaz.
\item
  \texttt{match}`i bir \texttt{let} ile değer döndürecek şekilde kullan.
\item
  \texttt{E0004}`ün nedenini ve çözümünü söyle.
\end{enumerate}

\begin{center}\rule{0.5\linewidth}{0.5pt}\end{center}

\subsection{📚 KAYNAKÇA}\label{kaynakça-7}

\subsubsection{Resmi Kaynaklar}\label{resmi-kaynaklar-7}

\begin{itemize}
\tightlist
\item
  \href{https://doc.rust-lang.org/book/ch06-02-match.html}{Rust Book --- The match Control Flow Construct}
\item
  \href{https://doc.rust-lang.org/book/ch19-03-pattern-syntax.html}{Rust Book --- Pattern Syntax}
\item
  \href{https://doc.rust-lang.org/reference/expressions/match-expr.html}{Rust Reference --- match expressions}
\item
  \href{https://doc.rust-lang.org/reference/patterns.html}{Rust Reference --- Patterns}
\end{itemize}

\subsubsection{Sürüm Notları}\label{sürüm-notları-7}

\begin{itemize}
\tightlist
\item
  \href{https://blog.rust-lang.org/2026/05/28/Rust-1.96.0/}{Rust Blog --- Announcing 1.96.0}
\item
  \href{https://doc.rust-lang.org/stable/releases.html\#version-1960-2026-05-28}{Rust 1.96.0 Release Notes}
\end{itemize}

\subsubsection{Hata Kodları}\label{hata-kodları-7}

\begin{itemize}
\tightlist
\item
  \href{https://doc.rust-lang.org/error_codes/E0004.html}{E0004}
\item
  \href{https://doc.rust-lang.org/error_codes/E0308.html}{E0308}
\item
  \href{https://doc.rust-lang.org/error_codes/E0408.html}{E0408}
\end{itemize}

\subsubsection{Ek Kaynaklar}\label{ek-kaynaklar-7}

\begin{itemize}
\tightlist
\item
  \href{https://doc.rust-lang.org/rust-by-example/flow_control/match.html}{Rust By Example --- match}
\item
  \href{https://dev.to/mark_saward/rust-match-guards-and-exhaustive-conditions-ad2}{DEV Community --- Rust match guards and exhaustive conditions}
\end{itemize}

\emph{Not: \texttt{match}`in eksiksiz olması, \texttt{\_} joker kolu, ilk eşleşen kolda durması, kolların aynı tipte olması (\texttt{E0308}), \texttt{E0004} mesajı ve guard'ların eksiksizliğe katılmaması aramada doğrulandı. Jump table, \texttt{E0408}, MIR \texttt{switchInt} ve \texttt{a..b} desenlerinin 1.80'den beri stabil olması bilgime dayanıyor; mesajları kendi makinende çalıştırarak kontrol et.}

\begin{center}\rule{0.5\linewidth}{0.5pt}\end{center}

\section{📘 KONU 9: loop / while / for}\label{konu-9-loop-while-for}

\textbf{Kategori:} Kontrol \textbar{} \textbf{Sürüm:} Rust 1.96.0

\begin{center}\rule{0.5\linewidth}{0.5pt}\end{center}

\subsection{🟢 KATMAN 1 --- ILK OKUMA}\label{katman-1-—-ilk-okuma-8}

\begin{Shaded}
\begin{Highlighting}[]
\KeywordTok{fn} \FunctionTok{main}\OperatorTok{()} \OperatorTok{\{}
    \KeywordTok{let} \KeywordTok{mut}\NormalTok{ count }\OperatorTok{=} \DecValTok{0}\OperatorTok{;}
    \KeywordTok{let}\NormalTok{ result }\OperatorTok{=} \ControlFlowTok{loop} \OperatorTok{\{}            \CommentTok{// loop: sonsuz, break değer döndürebilir}
\NormalTok{        count }\OperatorTok{+=} \DecValTok{1}\OperatorTok{;}
        \ControlFlowTok{if}\NormalTok{ count }\OperatorTok{==} \DecValTok{5} \OperatorTok{\{}
            \ControlFlowTok{break}\NormalTok{ count }\OperatorTok{*} \DecValTok{2}\OperatorTok{;}       \CommentTok{// break ile değer çıkar}
        \OperatorTok{\}}
    \OperatorTok{\};}
\NormalTok{    println}\OperatorTok{!(}\StringTok{"result = \{result\}"}\OperatorTok{);}

    \KeywordTok{let} \KeywordTok{mut}\NormalTok{ n }\OperatorTok{=} \DecValTok{3}\OperatorTok{;}
    \ControlFlowTok{while}\NormalTok{ n }\OperatorTok{!=} \DecValTok{0} \OperatorTok{\{}                 \CommentTok{// while: koşul doğruyken döner}
\NormalTok{        println}\OperatorTok{!(}\StringTok{"\{n\}!"}\OperatorTok{);}
\NormalTok{        n }\OperatorTok{-=} \DecValTok{1}\OperatorTok{;}
    \OperatorTok{\}}
\NormalTok{    println}\OperatorTok{!(}\StringTok{"LIFTOFF!"}\OperatorTok{);}

    \KeywordTok{let}\NormalTok{ arr }\OperatorTok{=} \OperatorTok{[}\DecValTok{10}\OperatorTok{,} \DecValTok{20}\OperatorTok{,} \DecValTok{30}\OperatorTok{];}
    \ControlFlowTok{for}\NormalTok{ x }\KeywordTok{in}\NormalTok{ arr }\OperatorTok{\{}                 \CommentTok{// for: iterator üzerinde döner}
\NormalTok{        println}\OperatorTok{!(}\StringTok{"value = \{x\}"}\OperatorTok{);}
    \OperatorTok{\}}

    \ControlFlowTok{for}\NormalTok{ i }\KeywordTok{in} \DecValTok{0}\OperatorTok{..}\DecValTok{3} \OperatorTok{\{}                \CommentTok{// Range: 0, 1, 2 (3 hariç)}
\NormalTok{        println}\OperatorTok{!(}\StringTok{"i = \{i\}"}\OperatorTok{);}
    \OperatorTok{\}}
\OperatorTok{\}}
\CommentTok{// Çıktı:}
\CommentTok{// result = 10}
\CommentTok{// 3!}
\CommentTok{// 2!}
\CommentTok{// 1!}
\CommentTok{// LIFTOFF!}
\CommentTok{// value = 10}
\CommentTok{// value = 20}
\CommentTok{// value = 30}
\CommentTok{// i = 0}
\CommentTok{// i = 1}
\CommentTok{// i = 2}
\end{Highlighting}
\end{Shaded}

Rust'ta üç loop (döngü) yapısı vardır. \texttt{loop} sonsuz döner, çıkmak için \texttt{break} gerekir; \texttt{break} bir değer taşıyabilir ve bu durumda \texttt{loop} bir expression (ifade) olur. \texttt{while} her turdan önce bir condition (koşul) kontrol eder, koşul \texttt{false} olunca durur. \texttt{for} bir iterator (yineleyici) üzerinde gezer; dizi, Vec veya \texttt{0..3} gibi bir Range (aralık) üzerinde çalışır ve en güvenli, en çok kullanılan döngüdür çünkü indeks hatası yapma riski yoktur.

\textbf{1.96.0 notu:} \texttt{loop}/\texttt{while}/\texttt{for} davranışı 1.96.0'da aynı.

\begin{center}\rule{0.5\linewidth}{0.5pt}\end{center}

\subsection{🔵 KATMAN 2 --- RECALL KARTLARI}\label{katman-2-—-recall-kartlari-8}

\begin{itemize}
\item
  \texttt{loop\ \{\ ...\ break;\ \}} → sonsuz döngü, \texttt{break} olmadan asla durmaz
  İlgili: \texttt{while\ true\ \{\ \}} ile aynı iş ama \texttt{loop} niyeti daha nettir
\item
  \texttt{break\ değer;} → sadece \texttt{loop}`ta çalışır:

\begin{Shaded}
\begin{Highlighting}[]
\KeywordTok{let}\NormalTok{ r }\OperatorTok{=} \ControlFlowTok{loop} \OperatorTok{\{} \ControlFlowTok{break} \DecValTok{7}\OperatorTok{;} \OperatorTok{\};}   \CommentTok{// r = 7}
\end{Highlighting}
\end{Shaded}


  Hata: \texttt{while} veya \texttt{for}`da \texttt{break\ değer;} → \texttt{E0571}: \texttt{break} with value from a \texttt{while} loop
\item
  \texttt{continue} → o turu atlar, bir sonraki tura geçer

\begin{Shaded}
\begin{Highlighting}[]
\ControlFlowTok{for}\NormalTok{ i }\KeywordTok{in} \DecValTok{0}\OperatorTok{..}\DecValTok{5} \OperatorTok{\{}
    \ControlFlowTok{if}\NormalTok{ i }\OperatorTok{\%} \DecValTok{2} \OperatorTok{==} \DecValTok{0} \OperatorTok{\{} \ControlFlowTok{continue}\OperatorTok{;} \OperatorTok{\}}
\NormalTok{    println}\OperatorTok{!(}\StringTok{"\{i\}"}\OperatorTok{);}   \CommentTok{// 1, 3}
\OperatorTok{\}}
\end{Highlighting}
\end{Shaded}
\item
  \texttt{while} koşulu \texttt{bool} olmalı:

\begin{Shaded}
\begin{Highlighting}[]
\ControlFlowTok{while}\NormalTok{ n }\OperatorTok{\{} \OperatorTok{\}}   \CommentTok{// E0308: expected bool, found integer}
\end{Highlighting}
\end{Shaded}
\item
  Label (etiket) ile iç içe döngüden çıkma:

\begin{Shaded}
\begin{Highlighting}[]
\NormalTok{'outer}\OperatorTok{:} \ControlFlowTok{for}\NormalTok{ i }\KeywordTok{in} \DecValTok{0}\OperatorTok{..}\DecValTok{3} \OperatorTok{\{}
    \ControlFlowTok{for}\NormalTok{ j }\KeywordTok{in} \DecValTok{0}\OperatorTok{..}\DecValTok{3} \OperatorTok{\{}
        \ControlFlowTok{if}\NormalTok{ j }\OperatorTok{==} \DecValTok{1} \OperatorTok{\{} \ControlFlowTok{continue}\NormalTok{ 'outer}\OperatorTok{;} \OperatorTok{\}}
        \ControlFlowTok{if}\NormalTok{ i }\OperatorTok{==} \DecValTok{2} \OperatorTok{\{} \ControlFlowTok{break}\NormalTok{ 'outer}\OperatorTok{;} \OperatorTok{\}}
\NormalTok{        println}\OperatorTok{!(}\StringTok{"\{i\},\{j\}"}\OperatorTok{);}
    \OperatorTok{\}}
\OperatorTok{\}}
\end{Highlighting}
\end{Shaded}


  İlgili: etiketler \texttt{'isim:} ile başlar, \texttt{break}/\texttt{continue} içindeki değişkenlerden ayrı bir isim alanıdır
\item
  \texttt{for} ile Range:

\begin{Shaded}
\begin{Highlighting}[]
\ControlFlowTok{for}\NormalTok{ i }\KeywordTok{in} \DecValTok{0}\OperatorTok{..}\DecValTok{5} \OperatorTok{\{} \OperatorTok{\}}     \CommentTok{// 0..4 (5 hariç)}
\ControlFlowTok{for}\NormalTok{ i }\KeywordTok{in} \DecValTok{0}\OperatorTok{..=}\DecValTok{5} \OperatorTok{\{} \OperatorTok{\}}    \CommentTok{// 0..5 (5 dahil)}
\ControlFlowTok{for}\NormalTok{ i }\KeywordTok{in} \OperatorTok{(}\DecValTok{0}\OperatorTok{..}\DecValTok{5}\OperatorTok{).}\FunctionTok{rev}\OperatorTok{()} \OperatorTok{\{} \OperatorTok{\}}   \CommentTok{// ters sırada: 4,3,2,1,0}
\end{Highlighting}
\end{Shaded}
\item
  \texttt{for} iterator'ı tüketir, \texttt{in} sonrası değer taşınır (Copy değilse):

\begin{Shaded}
\begin{Highlighting}[]
\KeywordTok{let}\NormalTok{ v }\OperatorTok{=}\NormalTok{ vec}\OperatorTok{![}\DecValTok{1}\OperatorTok{,} \DecValTok{2}\OperatorTok{,} \DecValTok{3}\OperatorTok{];}
\ControlFlowTok{for}\NormalTok{ x }\KeywordTok{in}\NormalTok{ v }\OperatorTok{\{} \OperatorTok{\}}        \CommentTok{// v taşındı}
\ControlFlowTok{for}\NormalTok{ x }\KeywordTok{in} \OperatorTok{\&}\NormalTok{v }\OperatorTok{\{} \OperatorTok{\}}       \CommentTok{// v kalır, x: \&i32}
\ControlFlowTok{for}\NormalTok{ x }\KeywordTok{in}\NormalTok{ v}\OperatorTok{.}\FunctionTok{iter}\OperatorTok{()} \OperatorTok{\{} \OperatorTok{\}} \CommentTok{// \&v ile aynı}
\end{Highlighting}
\end{Shaded}


  Hata: \texttt{for} sonrası \texttt{v} kullanılırsa → \texttt{E0382}: borrow of moved value
\item
  \texttt{\_} ile değeri kullanmamak: \texttt{for\ \_\ in\ 0..3\ \{\ \}}
\item
  \texttt{while\ let} (desen eşleşirken döner --- Konu 18, 20'de detaylandırılacak):

\begin{Shaded}
\begin{Highlighting}[]
\KeywordTok{let} \KeywordTok{mut}\NormalTok{ stack }\OperatorTok{=}\NormalTok{ vec}\OperatorTok{![}\DecValTok{1}\OperatorTok{,} \DecValTok{2}\OperatorTok{,} \DecValTok{3}\OperatorTok{];}
\ControlFlowTok{while} \KeywordTok{let} \ConstantTok{Some}\OperatorTok{(}\NormalTok{top}\OperatorTok{)} \OperatorTok{=}\NormalTok{ stack}\OperatorTok{.}\FunctionTok{pop}\OperatorTok{()} \OperatorTok{\{}\NormalTok{ println}\OperatorTok{!(}\StringTok{"\{top\}"}\OperatorTok{);} \OperatorTok{\}}
\end{Highlighting}
\end{Shaded}
\end{itemize}

\textbf{GoT bağlantı haritası:}

\begin{verbatim}
loop ──┬── break ──► değer döndürebilir (break val)
       └── sonsuz, koşulsuz
while ──┬── koşul: bool
        └── break değer YOK (E0571)
for ──┬── iterator üzerinde ──► Range (0..n, 0..=n)
      ├── in v / in &v / in v.iter() ──► move / borrow
      └── en güvenli, indeks hatası yok
ortak: break / continue / label ('isim:)
\end{verbatim}

\begin{center}\rule{0.5\linewidth}{0.5pt}\end{center}

\subsection{🟣 KATMAN 3 --- DERIN TEKNIK NOT}\label{katman-3-—-derin-teknik-not-8}

\textbf{Adım 1 --- Bellek Seviyesi}

\begin{itemize}
\tightlist
\item
  Önce soru: ``Bellekte ne oluyor?'' Cevap: üç döngü de assembly'de bir geri-atlama (jump) ve koşullu çıkışa dönüşür; ek bir veri yapısı yoktur.
\item
  \texttt{for\ x\ in\ collection} aslında \texttt{IntoIterator::into\_iter(collection)} çağırır ve dönen iterator üzerinde tekrar tekrar \texttt{next()} çağırır; \texttt{Some(x)} geldikçe gövde çalışır, \texttt{None} gelince döngü biter. Bu yüzden \texttt{for} aslında bir \texttt{loop} + \texttt{match} olarak genişler (syntactic sugar --- sözdizimsel şeker).
\item
  \texttt{for\ x\ in\ v} yazarsan \texttt{v}`nin sahipliği \texttt{into\_iter()}`a geçer (taşınır); \texttt{for\ x\ in\ \&v} yazarsan sadece referans alınır, \texttt{v} scope sonunda hâlâ geçerlidir.
\item
  \texttt{loop} içindeki \texttt{break\ val} ifadesinde \texttt{val}`in bellekteki yeri, \texttt{if}`teki gibi tek bir sonuç slotudur.
\end{itemize}

\textbf{Adım 2 --- Compiler Davranışı}

\begin{itemize}
\tightlist
\item
  Önce soru: ''\texttt{for} derleyiciye nasıl görünür?'' Cevap: desugar (çözülme) sonrası şuna benzer:

\begin{Shaded}
\begin{Highlighting}[]
\KeywordTok{let} \KeywordTok{mut}\NormalTok{ iter }\OperatorTok{=}\NormalTok{ IntoIterator}\OperatorTok{::}\FunctionTok{into\_iter}\OperatorTok{(}\NormalTok{collection}\OperatorTok{);}
\ControlFlowTok{loop} \OperatorTok{\{}
    \ControlFlowTok{match}\NormalTok{ iter}\OperatorTok{.}\FunctionTok{next}\OperatorTok{()} \OperatorTok{\{}
        \ConstantTok{Some}\OperatorTok{(}\NormalTok{x}\OperatorTok{)} \OperatorTok{=\textgreater{}} \OperatorTok{\{} \CommentTok{/* gövde */} \OperatorTok{\}}
        \ConstantTok{None} \OperatorTok{=\textgreater{}} \ControlFlowTok{break}\OperatorTok{,}
    \OperatorTok{\}}
\OperatorTok{\}}
\end{Highlighting}
\end{Shaded}
\item
  \texttt{loop}`un tipi, içindeki tüm \texttt{break\ val} ifadelerinin tipinden çıkarılır; hiç \texttt{break\ val} yoksa (veya hiç \texttt{break} yoksa) \texttt{loop}`un tipi \texttt{!} (never type) olur ve bu da onu her yerde kullanılabilir yapar.
\item
  \texttt{while} ve \texttt{for}`un tipi her zaman \texttt{()}`tir, çünkü gövdenin hiç çalışmama ihtimali vardır; bu yüzden \texttt{break\ val} kabul etmezler (\texttt{E0571}).
\item
  Etiketler (\texttt{'outer}) lifetime (yaşam süresi) etiketleriyle aynı sözdizimini kullanır ama farklı bir isim alanındadır; birbirine karışmaz.
\item
  MIR'de her üç döngü de benzer bir \texttt{loop} bloğu + \texttt{SwitchInt} + geri sıçrama olarak görünür. Görmek için: \texttt{cargo\ rustc\ --\ --emit=mir}.
\end{itemize}

\textbf{Adım 3 --- Hata Kodları}

\begin{itemize}
\tightlist
\item
  \textbf{E0571:} \texttt{while}/\texttt{for} içinde \texttt{break\ değer;}. Çözüm: \texttt{loop} kullan veya değer için ayrı bir \texttt{mut} değişken tut.
\item
  \textbf{E0308:} \texttt{while} koşulu \texttt{bool} değil. Çözüm: karşılaştırma yaz (\texttt{n\ !=\ 0}).
\item
  \textbf{E0382:} \texttt{for\ x\ in\ v} sonrası \texttt{v}`yi tekrar kullandın. Çözüm: \texttt{for\ x\ in\ \&v} kullan.
\item
  \textbf{E0499 / E0502 (ileride göreceksin):} döngü içinde aynı anda birden fazla mutable borrow almaya çalışmak.
\item
  \textbf{Derleyici uyarısı:} \texttt{unused\_labels}, kullanılmayan \texttt{'outer} etiketi.
\item
  Detay için: \texttt{rustc\ --explain\ E0571}.
\end{itemize}

\textbf{Adım 4 --- ``Neden Böyle?''}

\begin{itemize}
\tightlist
\item
  C'de \texttt{for\ (int\ i\ =\ 0;\ i\ \textless{}\ n;\ i++)} üç ayrı ifade elle yazılır; sınır hatası (\texttt{\textless{}=} yerine \texttt{\textless{}} unutmak) klasik bir bug kaynağıdır. Rust'ın \texttt{for\ x\ in\ 0..n} yapısı bu sınır hesabını Range'e devreder.
\item
  C'de \texttt{do\ \{\ \}\ while} vardır; Rust'ta bunun yerine \texttt{loop\ \{\ ...;\ if\ !cond\ \{\ break;\ \}\ \}} yazılır --- ayrı bir sözdizimine gerek yoktur.
\item
  \texttt{break}`in değer taşıyabilmesi, döngü sonucunu geçici bir \texttt{mut} değişkene atama ihtiyacını ortadan kaldırır.
\item
  Label'lı \texttt{break}/\texttt{continue}, C'deki \texttt{goto}`ya ihtiyaç duymadan iç içe döngülerden çıkmayı sağlar.
\end{itemize}

\textbf{Adım 5 --- Rust 1.96.0 Farkı}

\begin{itemize}
\tightlist
\item
  Rust 1.96.0'daki Range* tiplerinin stabilize edilmesi range tabanlı iteration API'lerini etkiliyor; ancak \texttt{loop}/\texttt{while}/\texttt{for} sözdizimi ve label kuralları değişmedi.
\item
  \texttt{while\ let} zincirleri (\texttt{while\ let\ Some(x)\ =\ a\ \&\&\ x\ \textgreater{}\ 5} gibi) edition 2024 ile ilgili bir yenilik alanıdır; bu konunun temel \texttt{while}/\texttt{for} kurallarını etkilemez ve henüz öğrenmediğin bir konu olduğu için şimdilik atla.
\item
  \texttt{assert\_matches!} makrosu ve Wasm linker güncellemeleri de temel loop/while/for davranışını etkilemiyor.
\end{itemize}

\begin{center}\rule{0.5\linewidth}{0.5pt}\end{center}

\subsection{🔴 KATMAN 4 --- SIK UNUTULANLAR + DERINLEŞTIRME}\label{katman-4-—-sik-unutulanlar-derinleştirme-8}

\textbf{Sık Unutulan 1: \texttt{break\ val} sadece \texttt{loop}`ta çalışır}

\begin{itemize}
\tightlist
\item
  Neden unutuluyor: üçü de ``döngü'' olduğu için aynı kurallara sahip sanılır.
\item
  Yol 1 (analoji): sadece \texttt{loop} bir ``sonuç kutusu'' sunar; \texttt{while}/\texttt{for} sadece ``yap ve çık'' der.
\item
  Yol 2 (kod): \texttt{let\ r\ =\ loop\ \{\ break\ 5;\ \};}
\item
  Yol 3 (hata): \texttt{while\ true\ \{\ break\ 5;\ \}} → \texttt{E0571:\ break\ with\ value\ from\ a\ while\ loop}
\end{itemize}

\textbf{Sık Unutulan 2: \texttt{for\ v} ile \texttt{for\ \&v} farkı}

\begin{itemize}
\tightlist
\item
  Neden unutuluyor: döngüden sonra \texttt{v}`yi tekrar kullanmak isteyebilirsin.
\item
  Yol 1 (analoji): \texttt{for\ x\ in\ v} = kutuyu devretmek; \texttt{for\ x\ in\ \&v} = kutuya bakmak, geri vermek.
\item
  Yol 2 (kod): \texttt{for\ x\ in\ \&v\ \{\ println!("\{x\}");\ \}\ println!("\{:?\}",\ v);}
\item
  Yol 3 (hata): \texttt{for\ x\ in\ v\ \{\ \}\ println!("\{:?\}",\ v);} → \texttt{E0382:\ borrow\ of\ moved\ value:\ v}
\end{itemize}

\textbf{Sık Unutulan 3: Label'ı unutup yanlış döngüden çıkmak}

\begin{itemize}
\tightlist
\item
  Neden unutuluyor: \texttt{break} varsayılan olarak sadece en yakın (innermost) döngüyü keser.
\item
  Yol 1 (görsel): etiketsiz \texttt{break} sadece ``bir alt kat'' çıkar, bütün binadan çıkmaz.
\item
  Yol 2 (kod): \texttt{'outer:\ for\ i\ in\ 0..3\ \{\ for\ j\ in\ 0..3\ \{\ if\ i\ ==\ 1\ \{\ break\ 'outer;\ \}\ \}\ \}}
\item
  Yol 3 (hata): derleyici hata vermez; sadece istediğin döngüden çıkmamış olursun (mantık hatası).
\end{itemize}

\textbf{Hazır Derinleştirme Prompt'u:}

\begin{verbatim}
Rust 1.96.0'da loop / while / for konusunu en detaylı şekilde anlat. Şunları kapsa:

1. BELLEK SEVİYESİ: Üç döngü de assembly'de nasıl görünür? break val'in sonucu
   nerede tutulur? for x in v ile for x in &v arasında bellekte ne fark var?

2. COMPILER DAVRANIŞI: for döngüsü IntoIterator/Iterator::next() olarak nasıl
   desugar oluyor? loop'un tipi break ifadelerinden nasıl çıkarılıyor (never type
   dahil)? Etiketlerin isim alanı nasıl çalışıyor? MIR'de loop blokları nasıl
   görünür? (--emit=mir)

3. HATA KODLARI: E0571, E0308, E0382. Her biri için: ne zaman, neden, nasıl
   çözülür. `rustc --explain EXXXX` çıktısını da ekle.

4. C/C++ KARŞILAŞTIRMASI: for(;;), do-while, goto. Rust'ın loop+break val ve
   label'lı break/continue yapısı bunların yerini nasıl alıyor?

5. ANALOJİ: Günlük hayattan bir benzetme yap (örn: koşu pisti, asansör katları).

6. RUST 1.96.0: Bu konuda son sürümde değişen bir şey var mı?

7. PRATİK: 3 tane kod örneği ver — kolay, orta, zor.

Açıklama Türkçe, teknik terimler İngilizce, kod yorumları Türkçe olsun.
Her bölümü ayrı başlıkla ver. Kod bloklarını çalıştırılabilir yaz.
\end{verbatim}

\begin{center}\rule{0.5\linewidth}{0.5pt}\end{center}

\subsection{📝 GÖREV (Boş dosyaya yaz)}\label{görev-boş-dosyaya-yaz-7}

\textbf{🟢 Kolay}

\begin{enumerate}
\def\labelenumi{\arabic{enumi}.}
\tightlist
\item
  \texttt{loop} ile \texttt{count}`u \texttt{0}`dan \texttt{5}`e kadar artır, \texttt{5} olunca \texttt{break} ile çık ve yazdır.
\item
  \texttt{while} ile \texttt{3}`ten \texttt{1}`e geri sayım yap ve her adımda yazdır.
\end{enumerate}

\textbf{🟡 Orta}
3. \texttt{let\ arr\ =\ {[}1,\ 2,\ 3,\ 4,\ 5{]};} için \texttt{for} ile dolaş, sadece çift sayıları yazdır (\texttt{continue} kullan).
4. \texttt{let\ result\ =\ loop\ \{\ ...\ break\ sayı\ *\ 2;\ \};} yaz; \texttt{sayı}`yı döngü içinde artırıp bir koşulda \texttt{break} ile iki katını döndür.

\textbf{🔴 Zor}
5. İç içe iki \texttt{for} döngüsü yaz; \texttt{'outer:} etiketi koy. \texttt{i\ ==\ j} olduğunda \texttt{continue\ 'outer}, \texttt{i\ ==\ 3} olduğunda \texttt{break\ 'outer} yap. Sonra \texttt{let\ v\ =\ vec!{[}1,2,3{]};\ for\ x\ in\ v\ \{\ \}} sonrası \texttt{v}`yi yazdırmayı dene, \texttt{E0382}`yi gör ve \texttt{\&v} ile düzelt.

\textbf{Başarı kriteri:} Hepsini hatasız yazarsan → öğrendin ✅

\subsubsection{✅ Kendini Test Et (2-3 saat sonra, sıfırdan)}\label{kendini-test-et-2-3-saat-sonra-sıfırdan-8}

\begin{enumerate}
\def\labelenumi{\arabic{enumi}.}
\tightlist
\item
  \texttt{loop} + \texttt{break\ değer} ile bir sonuç üret.
\item
  \texttt{for} ile bir Range üzerinde dolaş.
\item
  Label'lı \texttt{break}/\texttt{continue} örneği yaz.
\end{enumerate}

\begin{center}\rule{0.5\linewidth}{0.5pt}\end{center}

\subsection{📚 KAYNAKÇA}\label{kaynakça-8}

\subsubsection{Resmi Kaynaklar}\label{resmi-kaynaklar-8}

\begin{itemize}
\tightlist
\item
  \href{https://doc.rust-lang.org/book/ch03-05-control-flow.html\#repetition-with-loops}{Rust Book --- Repetition with Loops}
\item
  \href{https://doc.rust-lang.org/reference/expressions/loop-expr.html}{Rust Reference --- Loops and other breakable expressions}
\item
  \href{https://doc.rust-lang.org/std/iter/trait.Iterator.html}{std::iter::Iterator}
\end{itemize}

\subsubsection{Sürüm Notları}\label{sürüm-notları-8}

\begin{itemize}
\tightlist
\item
  \href{https://blog.rust-lang.org/2026/05/28/Rust-1.96.0/}{Rust Blog --- Announcing 1.96.0}
\item
  \href{https://doc.rust-lang.org/stable/releases.html\#version-1960-2026-05-28}{Rust 1.96.0 Release Notes}
\end{itemize}

\subsubsection{Hata Kodları}\label{hata-kodları-8}

\begin{itemize}
\tightlist
\item
  \href{https://doc.rust-lang.org/error_codes/E0571.html}{E0571}
\item
  \href{https://doc.rust-lang.org/error_codes/E0308.html}{E0308}
\item
  \href{https://doc.rust-lang.org/error_codes/E0382.html}{E0382}
\end{itemize}

\subsubsection{Ek Kaynaklar}\label{ek-kaynaklar-8}

\begin{itemize}
\tightlist
\item
  \href{https://doc.rust-lang.org/rust-by-example/flow_control/loop/nested.html}{Rust By Example --- loop, nesting and labels}
\item
  \href{https://www.luisllamas.es/en/rust-loops-while-for/}{Luis Llamas --- Rust loops: loop, while, for}
\end{itemize}

\emph{Not: \texttt{break\ val}`in sadece \texttt{loop}`ta çalıştığı, label'lı \texttt{break}/\texttt{continue} davranışı ve \texttt{for}`un iterator üzerinde çalıştığı aramada Reference ve ek kaynaklar üzerinden doğrulandı. \texttt{for}`un desugar edilmiş hâli, MIR görünümü ve hata kodları (E0571, E0382) bilgime dayanıyor; mesajları kendi makinende çalıştırarak kontrol et.}

\begin{center}\rule{0.5\linewidth}{0.5pt}\end{center}

\section{📘 KONU 10: Functions (fn, params, return)}\label{konu-10-functions-fn-params-return}

\textbf{Kategori:} Fonksiyon \textbar{} \textbf{Sürüm:} Rust 1.96.0

\begin{center}\rule{0.5\linewidth}{0.5pt}\end{center}

\subsection{🟢 KATMAN 1 --- ILK OKUMA}\label{katman-1-—-ilk-okuma-9}

\begin{Shaded}
\begin{Highlighting}[]
\KeywordTok{fn} \FunctionTok{main}\OperatorTok{()} \OperatorTok{\{}
    \FunctionTok{greet}\OperatorTok{(}\StringTok{"Ferris"}\OperatorTok{);}                    \CommentTok{// çağrı: fonksiyon main'den SONRA tanımlı olsa da çalışır}
    \KeywordTok{let}\NormalTok{ sum }\OperatorTok{=} \FunctionTok{add}\OperatorTok{(}\DecValTok{3}\OperatorTok{,} \DecValTok{4}\OperatorTok{);}
\NormalTok{    println}\OperatorTok{!(}\StringTok{"sum = \{sum\}"}\OperatorTok{);}
\NormalTok{    println}\OperatorTok{!(}\StringTok{"doubled = \{\}"}\OperatorTok{,} \FunctionTok{double}\OperatorTok{(}\NormalTok{sum}\OperatorTok{));}
\OperatorTok{\}}

\KeywordTok{fn} \FunctionTok{greet}\OperatorTok{(}\NormalTok{name}\OperatorTok{:} \OperatorTok{\&}\DataTypeTok{str}\OperatorTok{)} \OperatorTok{\{}                  \CommentTok{// parametre tipi zorunlu}
\NormalTok{    println}\OperatorTok{!(}\StringTok{"Hello, \{name\}!"}\OperatorTok{);}
\OperatorTok{\}}

\KeywordTok{fn} \FunctionTok{add}\OperatorTok{(}\NormalTok{a}\OperatorTok{:} \DataTypeTok{i32}\OperatorTok{,}\NormalTok{ b}\OperatorTok{:} \DataTypeTok{i32}\OperatorTok{)} \OperatorTok{-\textgreater{}} \DataTypeTok{i32} \OperatorTok{\{}         \CommentTok{// -\textgreater{} dönüş tipi}
\NormalTok{    a }\OperatorTok{+}\NormalTok{ b                               }\CommentTok{// son ifade, noktalı virgülsüz → dönüş değeri}
\OperatorTok{\}}

\KeywordTok{fn} \FunctionTok{double}\OperatorTok{(}\NormalTok{x}\OperatorTok{:} \DataTypeTok{i32}\OperatorTok{)} \OperatorTok{-\textgreater{}} \DataTypeTok{i32} \OperatorTok{\{}
    \ControlFlowTok{return}\NormalTok{ x }\OperatorTok{*} \DecValTok{2}\OperatorTok{;}                       \CommentTok{// return da kullanılabilir (erken çıkış için tipik)}
\OperatorTok{\}}
\CommentTok{// Çıktı:}
\CommentTok{// Hello, Ferris!}
\CommentTok{// sum = 7}
\CommentTok{// doubled = 14}
\end{Highlighting}
\end{Shaded}

Rust'ta fonksiyonlar \texttt{fn} ile tanımlanır ve her parametrenin tipi açıkça yazılmak zorundadır; Rust bunu çıkaramaz. Fonksiyonlar dosyada herhangi bir sırada tanımlanabilir, \texttt{main}`den önce ya da sonra olması önemli değildir --- önemli olan aynı scope'ta (kapsam) görünür olmasıdır. Bir fonksiyonun gövdesi bir block (blok) expression'dır (ifade); son satır noktalı virgülsüz bırakılırsa o satırın değeri otomatik dönüş değeri olur. \texttt{return} anahtar kelimesi genelde fonksiyonun ortasından erken çıkmak için kullanılır.

\textbf{1.96.0 notu:} Fonksiyon tanımlama kuralları 1.96.0'da aynı.

\begin{center}\rule{0.5\linewidth}{0.5pt}\end{center}

\subsection{🔵 KATMAN 2 --- RECALL KARTLARI}\label{katman-2-—-recall-kartlari-9}

\begin{itemize}
\tightlist
\item
  Temel imza: \texttt{fn\ isim(param:\ Tip,\ param2:\ Tip2)\ -\textgreater{}\ DönüşTipi\ \{\ ...\ \}}
  İlgili: dönüş tipi yoksa \texttt{()} (unit) varsayılır
\item
  Parametre tipi zorunlu:

\begin{Shaded}
\begin{Highlighting}[]
\KeywordTok{fn} \FunctionTok{f}\OperatorTok{(}\NormalTok{x}\OperatorTok{)} \OperatorTok{\{} \OperatorTok{\}}   \CommentTok{// error: expected one of `:`, found `)`}
\KeywordTok{fn} \FunctionTok{f}\OperatorTok{(}\NormalTok{x}\OperatorTok{:} \DataTypeTok{i32}\OperatorTok{)} \OperatorTok{\{} \OperatorTok{\}}   \CommentTok{// doğrusu}
\end{Highlighting}
\end{Shaded}
\item
  Son ifade = dönüş değeri (noktalı virgülsüz):

\begin{Shaded}
\begin{Highlighting}[]
\KeywordTok{fn} \FunctionTok{add}\OperatorTok{(}\NormalTok{a}\OperatorTok{:} \DataTypeTok{i32}\OperatorTok{,}\NormalTok{ b}\OperatorTok{:} \DataTypeTok{i32}\OperatorTok{)} \OperatorTok{-\textgreater{}} \DataTypeTok{i32} \OperatorTok{\{}
\NormalTok{    a }\OperatorTok{+}\NormalTok{ b       }\CommentTok{// dönüş değeri}
\OperatorTok{\}}
\end{Highlighting}
\end{Shaded}


  Hata: sona \texttt{;} koyarsan:

\begin{Shaded}
\begin{Highlighting}[]
\KeywordTok{fn} \FunctionTok{add}\OperatorTok{(}\NormalTok{a}\OperatorTok{:} \DataTypeTok{i32}\OperatorTok{,}\NormalTok{ b}\OperatorTok{:} \DataTypeTok{i32}\OperatorTok{)} \OperatorTok{-\textgreater{}} \DataTypeTok{i32} \OperatorTok{\{}
\NormalTok{    a }\OperatorTok{+}\NormalTok{ b}\OperatorTok{;}      \CommentTok{// artık bir statement, değer üretmiyor}
\OperatorTok{\}}
\CommentTok{// E0308: mismatched types, expected i32, found ()}
\end{Highlighting}
\end{Shaded}
\item
  Erken çıkış: \texttt{return\ değer;}

\begin{Shaded}
\begin{Highlighting}[]
\KeywordTok{fn} \FunctionTok{abs}\OperatorTok{(}\NormalTok{x}\OperatorTok{:} \DataTypeTok{i32}\OperatorTok{)} \OperatorTok{-\textgreater{}} \DataTypeTok{i32} \OperatorTok{\{}
    \ControlFlowTok{if}\NormalTok{ x }\OperatorTok{\textless{}} \DecValTok{0} \OperatorTok{\{} \ControlFlowTok{return} \OperatorTok{-}\NormalTok{x}\OperatorTok{;} \OperatorTok{\}}
\NormalTok{    x}
\OperatorTok{\}}
\end{Highlighting}
\end{Shaded}
\item
  Dönüş tipi \texttt{()} (hiçbir şey döndürmeyen fonksiyon):

\begin{Shaded}
\begin{Highlighting}[]
\KeywordTok{fn} \FunctionTok{log}\OperatorTok{(}\NormalTok{msg}\OperatorTok{:} \OperatorTok{\&}\DataTypeTok{str}\OperatorTok{)} \OperatorTok{\{}\NormalTok{ println}\OperatorTok{!(}\StringTok{"\{msg\}"}\OperatorTok{);} \OperatorTok{\}}   \CommentTok{// -\textgreater{} () yazmaya gerek yok}
\end{Highlighting}
\end{Shaded}


  Hata: gövdede değer varsa ama \texttt{-\textgreater{}} yoksa:

\begin{Shaded}
\begin{Highlighting}[]
\KeywordTok{fn} \FunctionTok{f}\OperatorTok{()} \OperatorTok{\{} \DecValTok{5} \OperatorTok{\}}
\CommentTok{// E0308: mismatched types, expected `()`, found integer}
\end{Highlighting}
\end{Shaded}
\item
  Fonksiyon sırası önemli değil (hoisting benzeri davranış):

\begin{Shaded}
\begin{Highlighting}[]
\KeywordTok{fn} \FunctionTok{main}\OperatorTok{()} \OperatorTok{\{} \FunctionTok{helper}\OperatorTok{();} \OperatorTok{\}}
\KeywordTok{fn} \FunctionTok{helper}\OperatorTok{()} \OperatorTok{\{}\NormalTok{ println}\OperatorTok{!(}\StringTok{"called"}\OperatorTok{);} \OperatorTok{\}}   \CommentTok{// main'den sonra tanımlı, sorun yok}
\end{Highlighting}
\end{Shaded}
\item
  Parametre ile argument (argüman) farkı: tanımdaki \texttt{x:\ i32} parametre, çağrıdaki \texttt{5} argümandır
\item
  Statement (ifade dizisi, değer üretmez) vs expression (değer üretir):

\begin{Shaded}
\begin{Highlighting}[]
\KeywordTok{let}\NormalTok{ y }\OperatorTok{=} \OperatorTok{(}\KeywordTok{let}\NormalTok{ x }\OperatorTok{=} \DecValTok{5}\OperatorTok{);}   \CommentTok{// error: `let` ifadedeği, statement'tır}
\KeywordTok{let}\NormalTok{ z }\OperatorTok{=} \OperatorTok{\{} \KeywordTok{let}\NormalTok{ a }\OperatorTok{=} \DecValTok{3}\OperatorTok{;}\NormalTok{ a }\OperatorTok{+} \DecValTok{1} \OperatorTok{\};}   \CommentTok{// z = 4, blok bir expression}
\end{Highlighting}
\end{Shaded}
\item
  Fonksiyon pointer'ı (ileri konularda detaylanacak):

\begin{Shaded}
\begin{Highlighting}[]
\KeywordTok{let}\NormalTok{ f}\OperatorTok{:} \KeywordTok{fn}\OperatorTok{(}\DataTypeTok{i32}\OperatorTok{,} \DataTypeTok{i32}\OperatorTok{)} \OperatorTok{-\textgreater{}} \DataTypeTok{i32} \OperatorTok{=}\NormalTok{ add}\OperatorTok{;}
\end{Highlighting}
\end{Shaded}
\end{itemize}

\textbf{GoT bağlantı haritası:}

\begin{verbatim}
fn isim(param: Tip) -> DönüşTip ──┬── gövde = block expression
                                   ├── son ifade (;'siz) ──► return değeri
                                   ├── return değer; ──► erken çıkış
                                   └── -> yoksa ──► () (unit)
statement (;'li) vs expression (;'siz) ──► E0308 kaynağı
tanım sırası önemsiz ──► aynı scope yeterli
\end{verbatim}

\begin{center}\rule{0.5\linewidth}{0.5pt}\end{center}

\subsection{🟣 KATMAN 3 --- DERIN TEKNIK NOT}\label{katman-3-—-derin-teknik-not-9}

\textbf{Adım 1 --- Bellek Seviyesi}

\begin{itemize}
\tightlist
\item
  Önce soru: ``Fonksiyon çağrılınca bellekte ne oluyor?'' Cevap: her çağrı, çağıran fonksiyonun üstüne yeni bir stack frame (yığın çerçevesi) iter. Parametreler bu frame'e kopyalanır (Copy tiplerde) veya taşınır (owned tiplerde, örneğin \texttt{String}).
\item
  Dönüş değeri, çağıranın beklediği yere (genelde bir register, örn. x86-64'te \texttt{rax}) yazılır. Büyük struct'lar dönülürken derleyici gizli bir pointer (return slot) da kullanabilir.
\item
  Fonksiyon gövdesindeki yerel değişkenler fonksiyon dönünce (scope bitince) drop edilir; frame stack'ten kalkar.
\item
  \texttt{return} ile erken çıkışta da aynı drop kuralları işler; o ana kadar oluşturulmuş yerel değişkenler sırayla temizlenir.
\end{itemize}

\textbf{Adım 2 --- Compiler Davranışı}

\begin{itemize}
\tightlist
\item
  Önce soru: ``Derleyici neyi kontrol ediyor?'' Cevap: her parametrenin ve dönüş değerinin tipi; gövdenin son ifadesinin tipi dönüş tipiyle eşleşmeli.
\item
  \texttt{;} bir expression'ı statement'a çevirir ve değeri \texttt{()}`e düşürür. Bu yüzden son satırdaki yanlışlıkla eklenen \texttt{;} tip hatasına yol açar.
\item
  Dosyadaki fonksiyonlar derleme öncesi bir item (öğe) listesi olarak toplanır; bu yüzden çağrı sırası, tanım sırasından bağımsızdır (ileri referans mümkündür).
\item
  Basit fonksiyonlar genelde inline (gövdeye gömme) edilebilir; derleyici, çağrı maliyetini release build'de optimize edebilir. Bunu garanti eden bir kural yoktur, bir optimizasyondur.
\item
  MIR'de her fonksiyon ayrı bir birim olarak görünür, çağrılar \texttt{Call} terminator'ı ile temsil edilir. Görmek için: \texttt{cargo\ rustc\ --\ --emit=mir}.
\end{itemize}

\textbf{Adım 3 --- Hata Kodları}

\begin{itemize}
\tightlist
\item
  \textbf{E0308 (son ifadede \texttt{;}):} gövde \texttt{()} üretti ama dönüş tipi başka. Çözüm: son \texttt{;}`yi kaldır veya \texttt{return} kullan.
\item
  \textbf{E0308 (dönüş tipi yokken değer var):} \texttt{-\textgreater{}\ Tip} eksik ama gövde değer üretiyor. Çözüm: \texttt{-\textgreater{}\ Tip} ekle veya son ifadeye \texttt{;} koy.
\item
  \textbf{E0061:} yanlış sayıda argüman. \texttt{this\ function\ takes\ 2\ arguments\ but\ 1\ argument\ was\ supplied}. Çözüm: argüman sayısını eşitle.
\item
  \textbf{E0308 (parametre tipi uyuşmazlığı):} \texttt{add("a",\ "b")} gibi yanlış tipte argüman. Çözüm: doğru tipte argüman ver.
\item
  \textbf{Kodsuz:} \texttt{missing\ type\ for\ function\ argument} (parametre tipi yazılmadı)
\item
  Detay için: \texttt{rustc\ --explain\ E0308} veya \texttt{rustc\ --explain\ E0061}.
\end{itemize}

\textbf{Adım 4 --- ``Neden Böyle?''}

\begin{itemize}
\tightlist
\item
  C'de parametre tipi yazmak zorunludur ama dönüş tipi yazmazsan eskiden \texttt{int} varsayılırdı (artık yasak). Rust'ta hem parametre hem dönüş tipi (varsa) her zaman açık yazılır; belirsizlik bırakılmaz.
\item
  C'de \texttt{return} zorunludur; Rust'ta son ifadenin otomatik dönmesi, ternary benzeri kısa ifadeler yazmayı kolaylaştırır (\texttt{if}/\texttt{match} gibi expression'larla birleşir).
\item
  Statement/expression ayrımı, Rust'ın hemen her yapıyı (if, match, block) bir değere çevirebilmesinin temelidir.
\item
  Trade-off: \texttt{;} unutmak ya da fazladan koymak yeni başlayanlar için sık karşılaşılan bir tuzaktır, ama kuralı öğrenince çok az hataya yol açar.
\end{itemize}

\textbf{Adım 5 --- Rust 1.96.0 Farkı}

\begin{itemize}
\tightlist
\item
  Rust 1.96.0'daki değişiklikler (yeni Range* tipleri, \texttt{assert\_matches!} makrosu, Wasm linker güncellemeleri) fonksiyon tanımlama ve parametre/dönüş tipi kurallarını değiştirmiyor.
\item
  Bu değişiklikler fonksiyonların son-ifade-dönüş kuralını da etkilemiyor.
\end{itemize}

\begin{center}\rule{0.5\linewidth}{0.5pt}\end{center}

\subsection{🔴 KATMAN 4 --- SIK UNUTULANLAR + DERINLEŞTIRME}\label{katman-4-—-sik-unutulanlar-derinleştirme-9}

\textbf{Sık Unutulan 1: Son satıra yanlışlıkla \texttt{;} koymak}

\begin{itemize}
\tightlist
\item
  Neden unutuluyor: her satırı \texttt{;} ile bitirme alışkanlığı güçlüdür.
\item
  Yol 1 (analoji): \texttt{;} bir kapıyı kapatır, değeri içeride hapseder; dönüş değeri dışarı çıkamaz.
\item
  Yol 2 (kod): \texttt{fn\ add(a:\ i32,\ b:\ i32)\ -\textgreater{}\ i32\ \{\ a\ +\ b\ \}} (son satırda \texttt{;} yok)
\item
  Yol 3 (hata): \texttt{a\ +\ b;} yazarsan → \texttt{E0308:\ mismatched\ types,\ expected\ i32,\ found\ ()}
\end{itemize}

\textbf{Sık Unutulan 2: Parametre tipini yazmayı unutmak}

\begin{itemize}
\tightlist
\item
  Neden unutuluyor: \texttt{let} içinde tip çoğu zaman yazılmaz, aynı alışkanlık fonksiyona taşınır.
\item
  Yol 1 (görsel): \texttt{let}`te tür çıkarımı var, \texttt{fn} parametresinde yok --- fonksiyon imzası bir sözleşmedir.
\item
  Yol 2 (kod): \texttt{fn\ f(x:\ i32)\ \{\ \}}
\item
  Yol 3 (hata): \texttt{fn\ f(x)\ \{\ \}} → \texttt{error:\ expected\ one\ of\ ...\ found\ ')'}
\end{itemize}

\textbf{Sık Unutulan 3: Argüman sayısı/tipi uyuşmazlığı}

\begin{itemize}
\tightlist
\item
  Neden unutuluyor: Python gibi dillerde varsayılan parametre olabilir, Rust'ta yok.
\item
  Yol 1 (analoji): fonksiyon imzası bir form gibidir; eksik ya da yanlış tipte alan bırakamazsın.
\item
  Yol 2 (kod): \texttt{add(3,\ 4)} --- tam iki \texttt{i32} argüman.
\item
  Yol 3 (hata): \texttt{add(3)} → \texttt{E0061:\ this\ function\ takes\ 2\ arguments\ but\ 1\ argument\ was\ supplied}
\end{itemize}

\textbf{Hazır Derinleştirme Prompt'u:}

\begin{verbatim}
Rust 1.96.0'da Functions (fn, params, return) konusunu en detaylı şekilde anlat. Şunları kapsa:

1. BELLEK SEVİYESİ: Fonksiyon çağrısında stack frame nasıl kurulur? Parametreler
   nasıl kopyalanır/taşınır? Dönüş değeri nerede tutulur (register vs return slot)?
   Yerel değişkenler ne zaman drop edilir?

2. COMPILER DAVRANIŞI: Statement/expression ayrımı nasıl çalışıyor? ; bir ifadeyi
   nasıl () yapıyor? Fonksiyonların tanım sırası neden önemsiz? MIR'de bir
   fonksiyon çağrısı (Call terminator) nasıl görünür? (--emit=mir)

3. HATA KODLARI: E0308 (iki farklı durum) ve E0061. Her biri için: ne zaman,
   neden, nasıl çözülür. `rustc --explain EXXXX` çıktısını da ekle.

4. C/C++ KARŞILAŞTIRMASI: Zorunlu return, parametre/dönüş tipi yazımı. Rust'ın
   son-ifade-dönüş kuralı C'den nasıl farklı?

5. ANALOJİ: Günlük hayattan bir benzetme yap (örn: fabrika hattı, form doldurma).

6. RUST 1.96.0: Bu konuda son sürümde değişen bir şey var mı?

7. PRATİK: 3 tane kod örneği ver — kolay, orta, zor.

Açıklama Türkçe, teknik terimler İngilizce, kod yorumları Türkçe olsun.
Her bölümü ayrı başlıkla ver. Kod bloklarını çalıştırılabilir yaz.
\end{verbatim}

\begin{center}\rule{0.5\linewidth}{0.5pt}\end{center}

\subsection{📝 GÖREV (Boş dosyaya yaz)}\label{görev-boş-dosyaya-yaz-8}

\textbf{🟢 Kolay}

\begin{enumerate}
\def\labelenumi{\arabic{enumi}.}
\tightlist
\item
  Parametresiz, dönüş değersiz bir \texttt{greet()} fonksiyonu yaz ve \texttt{main}`den çağır.
\item
  İki \texttt{i32} parametre alıp çarpımını döndüren \texttt{multiply(a,\ b)} fonksiyonu yaz.
\end{enumerate}

\textbf{🟡 Orta}
3. \texttt{is\_even(n:\ i32)\ -\textgreater{}\ bool} yaz (\texttt{n\ \%\ 2\ ==\ 0} son ifade olarak dönsün).
4. \texttt{max(a:\ i32,\ b:\ i32)\ -\textgreater{}\ i32} yaz; \texttt{if} ile \texttt{return} kullanarak erken çıkış yap.

\textbf{🔴 Zor}
5. \texttt{fn\ square(x:\ i32)\ -\textgreater{}\ i32\ \{\ x\ *\ x;\ \}} yaz (bilerek sona \texttt{;} koy), \texttt{E0308} hatasını gör ve düzelt. Sonra \texttt{main}`den önce tanımlanmış bir \texttt{helper()} fonksiyonunu \texttt{main} içinden çağır (tanım sırasının önemsizliğini göster). Son olarak 3 parametreli bir fonksiyonu 2 argümanla çağırıp \texttt{E0061}`i gör.

\textbf{Başarı kriteri:} Hepsini hatasız yazarsan → öğrendin ✅

\subsubsection{✅ Kendini Test Et (2-3 saat sonra, sıfırdan)}\label{kendini-test-et-2-3-saat-sonra-sıfırdan-9}

\begin{enumerate}
\def\labelenumi{\arabic{enumi}.}
\tightlist
\item
  Parametre alan, değer döndüren bir fonksiyon yaz.
\item
  Son ifadeyle dönüş ile \texttt{return} ile dönüşü birer örnekte göster.
\item
  \texttt{;} fazlalığının neden hata verdiğini söyle.
\end{enumerate}

\begin{center}\rule{0.5\linewidth}{0.5pt}\end{center}

\subsection{📚 KAYNAKÇA}\label{kaynakça-9}

\subsubsection{Resmi Kaynaklar}\label{resmi-kaynaklar-9}

\begin{itemize}
\tightlist
\item
  \href{https://doc.rust-lang.org/book/ch03-03-how-functions-work.html}{Rust Book --- Functions}
\item
  \href{https://doc.rust-lang.org/reference/items/functions.html}{Rust Reference --- Functions}
\item
  \href{https://doc.rust-lang.org/reference/statements-and-expressions.html}{Rust Reference --- Statements and Expressions}
\end{itemize}

\subsubsection{Sürüm Notları}\label{sürüm-notları-9}

\begin{itemize}
\tightlist
\item
  \href{https://blog.rust-lang.org/2026/05/28/Rust-1.96.0/}{Rust Blog --- Announcing 1.96.0}
\item
  \href{https://doc.rust-lang.org/stable/releases.html\#version-1960-2026-05-28}{Rust 1.96.0 Release Notes}
\end{itemize}

\subsubsection{Hata Kodları}\label{hata-kodları-9}

\begin{itemize}
\tightlist
\item
  \href{https://doc.rust-lang.org/error_codes/E0308.html}{E0308}
\item
  \href{https://doc.rust-lang.org/error_codes/E0061.html}{E0061}
\end{itemize}

\subsubsection{Ek Kaynaklar}\label{ek-kaynaklar-9}

\begin{itemize}
\tightlist
\item
  \href{https://doc.rust-lang.org/rust-by-example/fn.html}{Rust By Example --- Functions}
\item
  \href{https://doc.rust-lang.org/rust-by-example/expression.html}{Rust By Example --- Statements and Expressions}
\end{itemize}

\emph{Not: Bu konu için bu oturumda canlı arama yapılmadı; içerik, Rust Book'un Functions bölümündeki yerleşik kurallara (parametre tipi zorunluluğu, son-ifade-dönüş, tanım sırasının önemsizliği) dayanıyor. Hata mesajlarını (\texttt{E0308}, \texttt{E0061}) ve MIR çıktısını kendi makinende çalıştırarak doğrulaman önerilir.}

\begin{center}\rule{0.5\linewidth}{0.5pt}\end{center}

\section{📘 KONU 11: Ownership (move, Copy, Clone)}\label{konu-11-ownership-move-copy-clone}

\textbf{Kategori:} Bellek \textbar{} \textbf{Sürüm:} Rust 1.96.0

\begin{center}\rule{0.5\linewidth}{0.5pt}\end{center}

\subsection{🟢 KATMAN 1 --- ILK OKUMA}\label{katman-1-—-ilk-okuma-10}

\begin{Shaded}
\begin{Highlighting}[]
\KeywordTok{fn} \FunctionTok{main}\OperatorTok{()} \OperatorTok{\{}
    \KeywordTok{let}\NormalTok{ s1 }\OperatorTok{=} \DataTypeTok{String}\OperatorTok{::}\FunctionTok{from}\OperatorTok{(}\StringTok{"hello"}\OperatorTok{);}
    \KeywordTok{let}\NormalTok{ s2 }\OperatorTok{=}\NormalTok{ s1}\OperatorTok{;}                      \CommentTok{// move: s1 artık geçersiz}
    \CommentTok{// println!("\{s1\}");              // bu satır açılırsa E0382}

    \KeywordTok{let}\NormalTok{ s3 }\OperatorTok{=}\NormalTok{ s2}\OperatorTok{.}\FunctionTok{clone}\OperatorTok{();}              \CommentTok{// clone: gerçek kopya, ikisi de geçerli}
\NormalTok{    println}\OperatorTok{!(}\StringTok{"\{s2\} \{s3\}"}\OperatorTok{);}

    \KeywordTok{let}\NormalTok{ x }\OperatorTok{=} \DecValTok{5}\OperatorTok{;}
    \KeywordTok{let}\NormalTok{ y }\OperatorTok{=}\NormalTok{ x}\OperatorTok{;}                        \CommentTok{// i32 Copy'dir, ikisi de geçerli}
\NormalTok{    println}\OperatorTok{!(}\StringTok{"\{x\} \{y\}"}\OperatorTok{);}

    \FunctionTok{takes\_ownership}\OperatorTok{(}\NormalTok{s3}\OperatorTok{);}              \CommentTok{// s3 fonksiyona taşınır}
    \CommentTok{// println!("\{s3\}");              // burada açılırsa E0382}

    \FunctionTok{makes\_copy}\OperatorTok{(}\NormalTok{x}\OperatorTok{);}                    \CommentTok{// x kopyalanır, burada hâlâ geçerli}
\NormalTok{    println}\OperatorTok{!(}\StringTok{"\{x\}"}\OperatorTok{);}
\OperatorTok{\}}

\KeywordTok{fn} \FunctionTok{takes\_ownership}\OperatorTok{(}\NormalTok{s}\OperatorTok{:} \DataTypeTok{String}\OperatorTok{)} \OperatorTok{\{}
\NormalTok{    println}\OperatorTok{!(}\StringTok{"got: \{s\}"}\OperatorTok{);}
\OperatorTok{\}}                                      \CommentTok{// s burada drop edilir}

\KeywordTok{fn} \FunctionTok{makes\_copy}\OperatorTok{(}\NormalTok{n}\OperatorTok{:} \DataTypeTok{i32}\OperatorTok{)} \OperatorTok{\{}
\NormalTok{    println}\OperatorTok{!(}\StringTok{"got: \{n\}"}\OperatorTok{);}
\OperatorTok{\}}
\CommentTok{// Çıktı:}
\CommentTok{// hello hello}
\CommentTok{// 5 5}
\CommentTok{// got: hello}
\CommentTok{// got: 5}
\CommentTok{// 5}
\end{Highlighting}
\end{Shaded}

Ownership (sahiplik), Rust'ın garbage collector olmadan bellek güvenliği sağlamasını mümkün kılan temel kuraldır: her değerin tam olarak bir owner'ı (sahibi) vardır, owner scope dışına çıkınca değer otomatik drop edilir (serbest bırakılır). \texttt{let\ s2\ =\ s1;} yazınca \texttt{String} gibi Copy olmayan bir tip move (taşınır); artık sadece \texttt{s2} geçerlidir, \texttt{s1}`i kullanmak derleme hatasıdır. \texttt{i32} gibi basit tipler Copy trait'ini (özelliğini) implement ettiği için taşınmaz, bit bit kopyalanır ve ikisi de kullanılabilir kalır. Gerçek bir bağımsız kopya istiyorsan \texttt{.clone()} çağırırsın.

\textbf{1.96.0 notu:} Ownership, move ve Copy/Clone kuralları 1.96.0'da aynı.

\begin{center}\rule{0.5\linewidth}{0.5pt}\end{center}

\subsection{🔵 KATMAN 2 --- RECALL KARTLARI}\label{katman-2-—-recall-kartlari-10}

\begin{itemize}
\item
  Move (varsayılan davranış, Copy olmayan tipler):

\begin{Shaded}
\begin{Highlighting}[]
\KeywordTok{let}\NormalTok{ s1 }\OperatorTok{=} \DataTypeTok{String}\OperatorTok{::}\FunctionTok{from}\OperatorTok{(}\StringTok{"hi"}\OperatorTok{);}
\KeywordTok{let}\NormalTok{ s2 }\OperatorTok{=}\NormalTok{ s1}\OperatorTok{;}
\NormalTok{println}\OperatorTok{!(}\StringTok{"\{s1\}"}\OperatorTok{);}
\end{Highlighting}
\end{Shaded}


  Hata: \texttt{E0382}: borrow of moved value: \texttt{s1} --- note: move occurs because \texttt{s1} has type \texttt{String}, which does not implement the \texttt{Copy} trait
  İlgili: \texttt{.clone()}, referans (Konu 12)
\item
  Copy (basit, stack'te yaşayan tipler): \texttt{i8..i128}, \texttt{u8..u128}, \texttt{f32}, \texttt{f64}, \texttt{bool}, \texttt{char}, bunların tuple'ları

\begin{Shaded}
\begin{Highlighting}[]
\KeywordTok{let}\NormalTok{ x }\OperatorTok{=} \DecValTok{5}\OperatorTok{;}
\KeywordTok{let}\NormalTok{ y }\OperatorTok{=}\NormalTok{ x}\OperatorTok{;}          \CommentTok{// kopya, x hâlâ geçerli}
\end{Highlighting}
\end{Shaded}


  İlgili: \texttt{Copy} trait'i \texttt{Drop} trait'i olan tiplerde implement edilemez
\item
  Clone: gerçek, bağımsız kopya oluşturur

\begin{Shaded}
\begin{Highlighting}[]
\KeywordTok{let}\NormalTok{ s1 }\OperatorTok{=} \DataTypeTok{String}\OperatorTok{::}\FunctionTok{from}\OperatorTok{(}\StringTok{"hi"}\OperatorTok{);}
\KeywordTok{let}\NormalTok{ s2 }\OperatorTok{=}\NormalTok{ s1}\OperatorTok{.}\FunctionTok{clone}\OperatorTok{();}   \CommentTok{// ikisi de geçerli, s1 ayrı bellek}
\end{Highlighting}
\end{Shaded}


  Hata: \texttt{Clone} yoksa → \texttt{E0599}: no method named \texttt{clone} found
\item
  Fonksiyona geçirme move sayılır:

\begin{Shaded}
\begin{Highlighting}[]
\KeywordTok{fn} \FunctionTok{takes}\OperatorTok{(}\NormalTok{s}\OperatorTok{:} \DataTypeTok{String}\OperatorTok{)} \OperatorTok{\{\}}
\KeywordTok{let}\NormalTok{ s }\OperatorTok{=} \DataTypeTok{String}\OperatorTok{::}\FunctionTok{from}\OperatorTok{(}\StringTok{"x"}\OperatorTok{);}
\FunctionTok{takes}\OperatorTok{(}\NormalTok{s}\OperatorTok{);}
\NormalTok{println}\OperatorTok{!(}\StringTok{"\{s\}"}\OperatorTok{);}   \CommentTok{// E0382}
\end{Highlighting}
\end{Shaded}


  İlgili: \texttt{\&s} ile borrow et (Konu 12), taşıma olmaz
\item
  Fonksiyondan dönen değer ownership'i geri verir:

\begin{Shaded}
\begin{Highlighting}[]
\KeywordTok{fn} \FunctionTok{gives}\OperatorTok{()} \OperatorTok{-\textgreater{}} \DataTypeTok{String} \OperatorTok{\{} \DataTypeTok{String}\OperatorTok{::}\FunctionTok{from}\OperatorTok{(}\StringTok{"mine"}\OperatorTok{)} \OperatorTok{\}}
\KeywordTok{let}\NormalTok{ s }\OperatorTok{=} \FunctionTok{gives}\OperatorTok{();}   \CommentTok{// s artık sahip}
\end{Highlighting}
\end{Shaded}
\item
  Kendi struct'ını Copy yapmak:

\begin{Shaded}
\begin{Highlighting}[]
\NormalTok{\#}\OperatorTok{[}\FunctionTok{derive}\OperatorTok{(}\NormalTok{Debug}\OperatorTok{,}\NormalTok{ Clone}\OperatorTok{,}\NormalTok{ Copy}\OperatorTok{)]}
\KeywordTok{struct}\NormalTok{ Point }\OperatorTok{\{}\NormalTok{ x}\OperatorTok{:} \DataTypeTok{i32}\OperatorTok{,}\NormalTok{ y}\OperatorTok{:} \DataTypeTok{i32} \OperatorTok{\}}
\end{Highlighting}
\end{Shaded}


  Hata: içinde \texttt{String} gibi Copy olmayan alan varsa → \texttt{E0204}: the trait Copy may not be implemented for this type
\item
  \texttt{Rc\textless{}T\textgreater{}} paylaşılan sahiplik sağlar (ileri konu, şimdilik sadece adını bil):

\begin{Shaded}
\begin{Highlighting}[]
\KeywordTok{use}\NormalTok{ std}\OperatorTok{::}\NormalTok{rc}\OperatorTok{::}\DataTypeTok{Rc}\OperatorTok{;}
\KeywordTok{let}\NormalTok{ a }\OperatorTok{=} \DataTypeTok{Rc}\OperatorTok{::}\FunctionTok{new}\OperatorTok{(}\DecValTok{5}\OperatorTok{);}
\KeywordTok{let}\NormalTok{ b }\OperatorTok{=}\NormalTok{ a}\OperatorTok{.}\FunctionTok{clone}\OperatorTok{();}   \CommentTok{// veri kopyalanmaz, sayaç artar}
\end{Highlighting}
\end{Shaded}
\item
  \texttt{drop(x);} ile elle erken serbest bırakma:

\begin{Shaded}
\begin{Highlighting}[]
\KeywordTok{let}\NormalTok{ s }\OperatorTok{=} \DataTypeTok{String}\OperatorTok{::}\FunctionTok{from}\OperatorTok{(}\StringTok{"x"}\OperatorTok{);}
\FunctionTok{drop}\OperatorTok{(}\NormalTok{s}\OperatorTok{);}
\CommentTok{// println!("\{s\}");   // E0382}
\end{Highlighting}
\end{Shaded}
\end{itemize}

\textbf{GoT bağlantı haritası:}

\begin{verbatim}
ownership ──┬── her değerin 1 owner'ı var ──► scope bitince drop
            ├── move (varsayılan) ──► Copy olmayan tipler ──► E0382
            ├── Copy ──► basit tipler, stack ──► kopyalanır, taşınmaz
            └── Clone ──► .clone() ──► bağımsız gerçek kopya
fonksiyona geçirme/dönüş ──► ownership transferi
referans (&) ──► Konu 12: Borrowing (taşımadan kullan)
\end{verbatim}

\begin{center}\rule{0.5\linewidth}{0.5pt}\end{center}

\subsection{🟣 KATMAN 3 --- DERIN TEKNIK NOT}\label{katman-3-—-derin-teknik-not-10}

\textbf{Adım 1 --- Bellek Seviyesi}

\begin{itemize}
\tightlist
\item
  Önce soru: ''\texttt{let\ s2\ =\ s1;} anında bellekte ne oluyor?'' Cevap: \texttt{String}`in stack'teki (ptr, len, cap) üçlüsü \texttt{s2}`ye bit bit kopyalanır; heap'teki veri (\texttt{"hello"} metni) hiç taşınmaz, kopyalanmaz. Sadece sahiplik \texttt{s2}`ye geçer ve \texttt{s1} compiler tarafından ``artık geçersiz'' olarak işaretlenir.
\item
  Bu yüzden move O(1)`dir: heap allocation (bellek ayırma) yapılmaz. \texttt{.clone()} ise heap'te yeni bir buffer ayırır ve veriyi byte byte kopyalar; bu yüzden pahalıdır.
\item
  \texttt{i32} gibi Copy tipler tamamen stack'te yaşar; ``move'' kavramı burada anlamsızdır çünkü kopyalamanın maliyeti taşımayla aynıdır (birkaç byte).
\item
  Scope sonunda owner drop edilince, \texttt{String}`in \texttt{Drop} implementasyonu çalışır ve heap buffer serbest bırakılır (\texttt{dealloc}). Taşınmış bir değer için bu çalışmaz --- derleyici ``bu zaten taşındı, drop etme'' diye işaretler (bu konsept MIR'de drop flag olarak bilinir).
\end{itemize}

\textbf{Adım 2 --- Compiler Davranışı}

\begin{itemize}
\tightlist
\item
  Önce soru: ``Derleyici taşınan değeri nasıl takip ediyor?'' Cevap: borrow checker her değişkenin ``geçerli mi, taşınmış mı'' durumunu static olarak (çalıştırmadan) izler; \texttt{s1} taşındıktan sonraki her kullanım derleme hatasıdır.
\item
  \texttt{Copy} trait'i, derleyiciye ``bu tipte \texttt{let\ y\ =\ x;} bir move değil, implicit kopyadır'' der. \texttt{Copy} sadece Rust'ın bildiği basit (stack-only) tipler ve bunları içeren struct/tuple'lar için \texttt{derive} edilebilir; \texttt{Drop} implement eden bir tip asla \texttt{Copy} olamaz (ikisi birbirini dışlar).
\item
  \texttt{.clone()} derin bir kopya garanti etmez; \texttt{Clone} implementasyonu tipin kendisine bağlıdır (örneğin \texttt{Rc::clone} veriyi kopyalamaz, referans sayacını artırır).
\item
  MIR'de taşınan değişkenler için derleyici drop flag'leri ekler; bu da koşullu drop mantığını (hangi dalda taşındı, hangisinde taşınmadı) yönetir. Görmek için: \texttt{cargo\ rustc\ --\ --emit=mir}.
\end{itemize}

\textbf{Adım 3 --- Hata Kodları}

\begin{itemize}
\tightlist
\item
  \textbf{E0382:} taşınmış (veya drop edilmiş) bir değeri kullandın. Çözüm: \texttt{.clone()} çağır, referans kullan (Konu 12) veya değeri tekrar atama.
\item
  \textbf{E0599 (clone yok):} tipte \texttt{Clone} implement edilmemiş. Çözüm: \texttt{\#{[}derive(Clone){]}} ekle (tipi sen tanımlıyorsan).
\item
  \textbf{E0204:} \texttt{Copy} implement edilemeyen bir tipe \texttt{\#{[}derive(Copy){]}} yazdın (içinde \texttt{String}, \texttt{Vec} gibi heap sahibi bir alan var). Çözüm: \texttt{Copy}`yi kaldır, \texttt{Clone} yeterli.
\item
  \textbf{E0507 (ileride detaylanacak):} borrow edilmiş içerikten değer taşımaya çalışmak (\texttt{for\ x\ in\ \&v\ \{\ let\ y\ =\ *x;\ \}} gibi, \texttt{x} Copy değilse).
\item
  Detay için: \texttt{rustc\ --explain\ E0382}.
\end{itemize}

\textbf{Adım 4 --- ``Neden Böyle?''}

\begin{itemize}
\tightlist
\item
  C/C++`ta bir pointer'ı kopyalarsan iki pointer aynı belleği gösterir; biri \texttt{free()} çağırınca diğeri dangling (sarkık) pointer olur --- double-free ve use-after-free hatalarının klasik kaynağı. Rust ownership kuralı bunu derleme zamanında imkânsız kılar: bir değerin her zaman tek bir sorumlusu vardır.
\item
  C++`ın move semantics'i (C++11, \texttt{std::move}) benzer bir fikri opsiyonel olarak sunar; Rust'ta move varsayılan davranıştır ve derleyici tarafından zorunlu kılınır.
\item
  Garbage collector'lı dillerde (Java, Python) bu sorun runtime'da bir GC ile çözülür; Rust bunu ek bir runtime maliyeti olmadan, compile time'da çözer.
\item
  Trade-off: öğrenme eğrisi diktir (E0382 yeni başlayanların en sık gördüğü hatadır), ama karşılığında çalışma zamanı bellek hataları (use-after-free, double-free, dangling pointer) sınıf olarak ortadan kalkar.
\end{itemize}

\textbf{Adım 5 --- Rust 1.96.0 Farkı}

\begin{itemize}
\tightlist
\item
  Rust 1.96.0'daki değişiklikler (yeni Range* tipleri, \texttt{assert\_matches!} makrosu, Wasm linker güncellemeleri) ownership ve move kurallarını değiştirmiyor.
\item
  Yeni Range* tipleri bazı \texttt{core::ops} API'lerini stabilize ediyor; bu, Rust'ın ownership, move ve \texttt{Copy}/\texttt{Clone} kurallarını değiştirmiyor. \texttt{assert\_matches!} ve Wasm linker güncellemeleri de bu kuralları etkilemiyor.
\end{itemize}

\begin{center}\rule{0.5\linewidth}{0.5pt}\end{center}

\subsection{🔴 KATMAN 4 --- SIK UNUTULANLAR + DERINLEŞTIRME}\label{katman-4-—-sik-unutulanlar-derinleştirme-10}

\textbf{Sık Unutulan 1: \texttt{let\ s2\ =\ s1;} taşır, kopyalamaz}

\begin{itemize}
\tightlist
\item
  Neden unutuluyor: başka dillerde (Python, JS) atama genelde ya referans paylaşır ya da örtük kopyalar, hiçbiri ``eskisini geçersiz kılmaz''.
\item
  Yol 1 (analoji): \texttt{s1} bir ev anahtarı; \texttt{s2\ =\ s1} anahtarı devretmektir, artık \texttt{s1} elinde anahtar yok.
\item
  Yol 2 (kod): \texttt{let\ s2\ =\ s1.clone();} yazarsan ikisi de geçerli kalır.
\item
  Yol 3 (hata): \texttt{E0382:\ borrow\ of\ moved\ value:\ s1}
\end{itemize}

\textbf{Sık Unutulan 2: Fonksiyona geçirmek de move'dur}

\begin{itemize}
\tightlist
\item
  Neden unutuluyor: ``sadece değişkenler arası atama taşır'' sanılır.
\item
  Yol 1 (görsel): fonksiyon parametresi de bir \texttt{let} gibidir; \texttt{f(s)} aslında \texttt{let\ param\ =\ s;} demektir.
\item
  Yol 2 (kod): \texttt{takes\_ownership(s3);} sonrası \texttt{s3} kullanılamaz; \texttt{\&s3} ile ver, taşınmasın.
\item
  Yol 3 (hata): \texttt{E0382:\ use\ of\ moved\ value:\ s3}
\end{itemize}

\textbf{Sık Unutulan 3: Hangi tipler Copy, hangileri değil}

\begin{itemize}
\tightlist
\item
  Neden unutuluyor: ``hepsi basit görünüyor'' diye düşünülür.
\item
  Yol 1 (analoji): sabit boyutlu, heap'e dokunmayan değerler (sayılar, bool, char) kartvizit gibidir --- kopyalaması bedava. Heap'te veri taşıyan \texttt{String}/\texttt{Vec} birer dosya dolabıdır --- kopyalamak pahalı, bu yüzden varsayılan davranış taşımaktır.
\item
  Yol 2 (kod): \texttt{i32}, \texttt{bool}, \texttt{char}, \texttt{(i32,\ i32)} → Copy; \texttt{String}, \texttt{Vec\textless{}T\textgreater{}}, \texttt{Box\textless{}T\textgreater{}} → değil.
\item
  Yol 3 (hata): \texttt{\#{[}derive(Copy){]}} bir \texttt{String} alanı olan struct'a eklenirse → \texttt{E0204}
\end{itemize}

\textbf{Hazır Derinleştirme Prompt'u:}

\begin{verbatim}
Rust 1.96.0'da Ownership (move, Copy, Clone) konusunu en detaylı şekilde anlat. Şunları kapsa:

1. BELLEK SEVİYESİ: let s2 = s1; anında stack ve heap'te tam olarak ne oluyor?
   Move neden O(1), clone neden O(n)? Drop flag kavramı nedir? Copy tipler için
   "kopya" ile "move" arasında runtime'da fark var mı?

2. COMPILER DAVRANIŞI: Borrow checker taşınmış değişkenleri nasıl takip ediyor?
   Copy trait derleyiciye ne söylüyor? Copy ile Drop neden birlikte olamıyor?
   MIR'de drop flag'ler nasıl görünür? (--emit=mir)

3. HATA KODLARI: E0382, E0599, E0204. Her biri için: ne zaman, neden, nasıl
   çözülür. `rustc --explain EXXXX` çıktısını da ekle.

4. C/C++ KARŞILAŞTIRMASI: Pointer kopyalama, double-free, use-after-free,
   std::move. Rust'ın zorunlu move semantics'i bu sorunları nasıl önlüyor?

5. ANALOJİ: Günlük hayattan bir benzetme yap (örn: ev anahtarı, fotokopi).

6. RUST 1.96.0: Bu konuda son sürümde değişen bir şey var mı?

7. PRATİK: 3 tane kod örneği ver — kolay, orta, zor.

Açıklama Türkçe, teknik terimler İngilizce, kod yorumları Türkçe olsun.
Her bölümü ayrı başlıkla ver. Kod bloklarını çalıştırılabilir yaz.
\end{verbatim}

\begin{center}\rule{0.5\linewidth}{0.5pt}\end{center}

\subsection{📝 GÖREV (Boş dosyaya yaz)}\label{görev-boş-dosyaya-yaz-9}

\textbf{🟢 Kolay}

\begin{enumerate}
\def\labelenumi{\arabic{enumi}.}
\tightlist
\item
  Bir \texttt{String} oluştur, başka bir değişkene ata (move), eski değişkeni yazdırmayı dene ve \texttt{E0382}`yi gör; sonra kaldırıp düzelt.
\item
  İki \texttt{i32} değişken arasında atama yap ve ikisini de yazdır (Copy davranışını göster).
\end{enumerate}

\textbf{🟡 Orta}
3. Bir \texttt{String}`i \texttt{.clone()} ile kopyala, ikisini de yazdır, sonra birini değiştir (\texttt{push\_str}) ve diğerinin etkilenmediğini göster.
4. \texttt{String} alan bir fonksiyon yaz, bir \texttt{String}`i ona geçir, sonra orijinali kullanmaya çalışıp \texttt{E0382}`yi gör.

\textbf{🔴 Zor}
5. \texttt{struct\ Point\ \{\ x:\ i32,\ y:\ i32\ \}} tanımla, \texttt{\#{[}derive(Debug,\ Clone,\ Copy){]}} ekle; iki \texttt{Point} arasında atama yap, ikisini de kullan. Sonra \texttt{String} alanlı bir struct'a bilerek \texttt{\#{[}derive(Copy){]}} eklemeyi dene, \texttt{E0204}`ü gör ve kaldır. Son olarak \texttt{drop(x);} ile elle bir değeri erken serbest bırak, sonra kullanmaya çalışıp hatayı gör.

\textbf{Başarı kriteri:} Hepsini hatasız yazarsan → öğrendin ✅

\subsubsection{✅ Kendini Test Et (2-3 saat sonra, sıfırdan)}\label{kendini-test-et-2-3-saat-sonra-sıfırdan-10}

\begin{enumerate}
\def\labelenumi{\arabic{enumi}.}
\tightlist
\item
  Move, Copy ve Clone'u birer kod örneğiyle ayırt et.
\item
  \texttt{E0382}`nin nedenini açıkla.
\item
  Hangi tiplerin Copy olduğunu listele.
\end{enumerate}

\begin{center}\rule{0.5\linewidth}{0.5pt}\end{center}

\subsection{📚 KAYNAKÇA}\label{kaynakça-10}

\subsubsection{Resmi Kaynaklar}\label{resmi-kaynaklar-10}

\begin{itemize}
\tightlist
\item
  \href{https://doc.rust-lang.org/book/ch04-01-what-is-ownership.html}{Rust Book --- What is Ownership?}
\item
  \href{https://doc.rust-lang.org/reference/special-types-and-traits.html\#copy}{Rust Reference --- Copy}
\item
  \href{https://doc.rust-lang.org/std/clone/trait.Clone.html}{std::clone::Clone}
\end{itemize}

\subsubsection{Sürüm Notları}\label{sürüm-notları-10}

\begin{itemize}
\tightlist
\item
  \href{https://blog.rust-lang.org/2026/05/28/Rust-1.96.0/}{Rust Blog --- Announcing 1.96.0}
\item
  \href{https://doc.rust-lang.org/stable/releases.html\#version-1960-2026-05-28}{Rust 1.96.0 Release Notes}
\end{itemize}

\subsubsection{Hata Kodları}\label{hata-kodları-10}

\begin{itemize}
\tightlist
\item
  \href{https://doc.rust-lang.org/error_codes/E0382.html}{E0382}
\item
  \href{https://doc.rust-lang.org/error_codes/E0599.html}{E0599}
\item
  \href{https://doc.rust-lang.org/error_codes/E0204.html}{E0204}
\end{itemize}

\subsubsection{Ek Kaynaklar}\label{ek-kaynaklar-10}

\begin{itemize}
\tightlist
\item
  \href{https://learn.microsoft.com/en-gb/training/modules/rust-memory-management/1-what-is-ownership}{Microsoft Learn --- What is ownership?}
\item
  \href{https://itnext.io/rust-ownership-50-code-examples-96203fcf79ea}{itnext.io --- Rust Ownership: 50 Code Examples}
\end{itemize}

\emph{Not: Move semantics, \texttt{E0382} mesajı, \texttt{Copy}/\texttt{Clone} farkı ve \texttt{E0204} davranışı aramada Rust Book ve resmi hata kodu sayfası üzerinden doğrulandı. Heap allocation maliyeti, drop flag kavramı ve MIR görünümü bilgime dayanıyor; \texttt{size\_of} ve panic/hata mesajlarını kendi makinende çalıştırarak kontrol et.}

\begin{center}\rule{0.5\linewidth}{0.5pt}\end{center}

\section{📘 KONU 12: Borrowing (\& ve \&mut)}\label{konu-12-borrowing-ve-mut}

\textbf{Kategori:} Bellek \textbar{} \textbf{Sürüm:} Rust 1.96.0

\begin{center}\rule{0.5\linewidth}{0.5pt}\end{center}

\subsection{🟢 KATMAN 1 --- ILK OKUMA}\label{katman-1-—-ilk-okuma-11}

\begin{Shaded}
\begin{Highlighting}[]
\KeywordTok{fn} \FunctionTok{main}\OperatorTok{()} \OperatorTok{\{}
    \KeywordTok{let}\NormalTok{ s1 }\OperatorTok{=} \DataTypeTok{String}\OperatorTok{::}\FunctionTok{from}\OperatorTok{(}\StringTok{"hello"}\OperatorTok{);}
    \KeywordTok{let}\NormalTok{ len }\OperatorTok{=} \FunctionTok{calculate\_length}\OperatorTok{(\&}\NormalTok{s1}\OperatorTok{);}          \CommentTok{// \&s1: borrow, taşımaz}
\NormalTok{    println}\OperatorTok{!(}\StringTok{"'\{s1\}' has length \{len\}"}\OperatorTok{);}      \CommentTok{// s1 hâlâ geçerli}

    \KeywordTok{let} \KeywordTok{mut}\NormalTok{ s2 }\OperatorTok{=} \DataTypeTok{String}\OperatorTok{::}\FunctionTok{from}\OperatorTok{(}\StringTok{"hello"}\OperatorTok{);}
    \FunctionTok{change}\OperatorTok{(\&}\KeywordTok{mut}\NormalTok{ s2}\OperatorTok{);}                          \CommentTok{// \&mut s2: değiştirmeye izinli borrow}
\NormalTok{    println}\OperatorTok{!(}\StringTok{"\{s2\}"}\OperatorTok{);}

    \KeywordTok{let}\NormalTok{ r1 }\OperatorTok{=} \OperatorTok{\&}\NormalTok{s1}\OperatorTok{;}
    \KeywordTok{let}\NormalTok{ r2 }\OperatorTok{=} \OperatorTok{\&}\NormalTok{s1}\OperatorTok{;}                              \CommentTok{// birden fazla immutable referans OK}
\NormalTok{    println}\OperatorTok{!(}\StringTok{"\{r1\} \{r2\}"}\OperatorTok{);}
\OperatorTok{\}}

\KeywordTok{fn} \FunctionTok{calculate\_length}\OperatorTok{(}\NormalTok{s}\OperatorTok{:} \OperatorTok{\&}\DataTypeTok{String}\OperatorTok{)} \OperatorTok{-\textgreater{}} \DataTypeTok{usize} \OperatorTok{\{}    \CommentTok{// s: referans, sahip değil}
\NormalTok{    s}\OperatorTok{.}\FunctionTok{len}\OperatorTok{()}
\OperatorTok{\}}                                              \CommentTok{// s burada drop EDİLMEZ (ödünç aldı)}

\KeywordTok{fn} \FunctionTok{change}\OperatorTok{(}\NormalTok{s}\OperatorTok{:} \OperatorTok{\&}\KeywordTok{mut} \DataTypeTok{String}\OperatorTok{)} \OperatorTok{\{}
\NormalTok{    s}\OperatorTok{.}\FunctionTok{push\_str}\OperatorTok{(}\StringTok{", world"}\OperatorTok{);}
\OperatorTok{\}}
\CommentTok{// Çıktı:}
\CommentTok{// 'hello' has length 5}
\CommentTok{// hello, world}
\CommentTok{// hello hello}
\end{Highlighting}
\end{Shaded}

Reference (referans), bir değere sahip olmadan ona işaret eden bir pointer'dır; \texttt{\&} ile oluşturulur ve bu işleme borrowing (ödünç alma) denir. \texttt{calculate\_length(\&s1)} çağrısı \texttt{s1}`in sahipliğini taşımaz, fonksiyon bitince \texttt{s1} hâlâ \texttt{main}`de geçerlidir. Varsayılan referanslar immutable'dır (değiştirilemez); değiştirmek istiyorsan hem değişkenin hem referansın \texttt{mut} olması gerekir, bunu \texttt{\&mut} ile yaparsın. Aynı anda birden fazla immutable referans sorun değildir, ama bir \texttt{\&mut} referans varken başka hiçbir referans (ne \texttt{\&} ne \texttt{\&mut}) olamaz --- bu kural Konu 13'te detaylandırılacak.

\textbf{1.96.0 notu:} Borrowing kuralları 1.96.0'da aynı.

\begin{center}\rule{0.5\linewidth}{0.5pt}\end{center}

\subsection{🔵 KATMAN 2 --- RECALL KARTLARI}\label{katman-2-—-recall-kartlari-11}

\begin{itemize}
\item
  Immutable referans: \texttt{\&T} --- okur, değiştiremez

\begin{Shaded}
\begin{Highlighting}[]
\KeywordTok{fn} \FunctionTok{f}\OperatorTok{(}\NormalTok{s}\OperatorTok{:} \OperatorTok{\&}\DataTypeTok{String}\OperatorTok{)} \OperatorTok{\{}\NormalTok{ println}\OperatorTok{!(}\StringTok{"\{s\}"}\OperatorTok{);} \OperatorTok{\}}
\end{Highlighting}
\end{Shaded}
\item
  Mutable referans: \texttt{\&mut\ T} --- hem değişken hem referans \texttt{mut} olmalı

\begin{Shaded}
\begin{Highlighting}[]
\KeywordTok{let} \KeywordTok{mut}\NormalTok{ s }\OperatorTok{=} \DataTypeTok{String}\OperatorTok{::}\FunctionTok{from}\OperatorTok{(}\StringTok{"hi"}\OperatorTok{);}
\KeywordTok{let}\NormalTok{ r }\OperatorTok{=} \OperatorTok{\&}\KeywordTok{mut}\NormalTok{ s}\OperatorTok{;}
\NormalTok{r}\OperatorTok{.}\FunctionTok{push\_str}\OperatorTok{(}\StringTok{"!"}\OperatorTok{);}
\end{Highlighting}
\end{Shaded}


  Hata: \texttt{s} \texttt{mut} değilse → \texttt{E0596}: cannot borrow as mutable, as it is behind a \texttt{\&} reference / not declared as mutable
\item
  Aynı anda tek bir \texttt{\&mut}:

\begin{Shaded}
\begin{Highlighting}[]
\KeywordTok{let} \KeywordTok{mut}\NormalTok{ s }\OperatorTok{=} \DataTypeTok{String}\OperatorTok{::}\FunctionTok{from}\OperatorTok{(}\StringTok{"hi"}\OperatorTok{);}
\KeywordTok{let}\NormalTok{ r1 }\OperatorTok{=} \OperatorTok{\&}\KeywordTok{mut}\NormalTok{ s}\OperatorTok{;}
\KeywordTok{let}\NormalTok{ r2 }\OperatorTok{=} \OperatorTok{\&}\KeywordTok{mut}\NormalTok{ s}\OperatorTok{;}      \CommentTok{// ikinci mutable referans}
\NormalTok{println}\OperatorTok{!(}\StringTok{"\{r1\} \{r2\}"}\OperatorTok{);}
\end{Highlighting}
\end{Shaded}


  Hata: \texttt{E0499}: cannot borrow \texttt{s} as mutable more than once at a time
\item
  \texttt{\&mut} varken \texttt{\&} olamaz (ve tersi):

\begin{Shaded}
\begin{Highlighting}[]
\KeywordTok{let} \KeywordTok{mut}\NormalTok{ s }\OperatorTok{=} \DataTypeTok{String}\OperatorTok{::}\FunctionTok{from}\OperatorTok{(}\StringTok{"hi"}\OperatorTok{);}
\KeywordTok{let}\NormalTok{ r1 }\OperatorTok{=} \OperatorTok{\&}\NormalTok{s}\OperatorTok{;}
\KeywordTok{let}\NormalTok{ r2 }\OperatorTok{=} \OperatorTok{\&}\KeywordTok{mut}\NormalTok{ s}\OperatorTok{;}      \CommentTok{// immutable varken mutable istedi}
\NormalTok{println}\OperatorTok{!(}\StringTok{"\{r1\}"}\OperatorTok{);}
\end{Highlighting}
\end{Shaded}


  Hata: \texttt{E0502}: cannot borrow \texttt{s} as mutable because it is also borrowed as immutable
\item
  NLL (Non-Lexical Lifetimes): referansın ömrü son kullanıldığı yerde biter, süslü parantezin sonunda değil:

\begin{Shaded}
\begin{Highlighting}[]
\KeywordTok{let} \KeywordTok{mut}\NormalTok{ s }\OperatorTok{=} \DataTypeTok{String}\OperatorTok{::}\FunctionTok{from}\OperatorTok{(}\StringTok{"hi"}\OperatorTok{);}
\KeywordTok{let}\NormalTok{ r1 }\OperatorTok{=} \OperatorTok{\&}\NormalTok{s}\OperatorTok{;}
\NormalTok{println}\OperatorTok{!(}\StringTok{"\{r1\}"}\OperatorTok{);}          \CommentTok{// r1'in son kullanımı burada}
\KeywordTok{let}\NormalTok{ r2 }\OperatorTok{=} \OperatorTok{\&}\KeywordTok{mut}\NormalTok{ s}\OperatorTok{;}           \CommentTok{// OK: r1 artık "ölü"}
\NormalTok{println}\OperatorTok{!(}\StringTok{"\{r2\}"}\OperatorTok{);}
\end{Highlighting}
\end{Shaded}
\item
  Dereference (\texttt{*}) ile değere ulaşmak:

\begin{Shaded}
\begin{Highlighting}[]
\KeywordTok{let}\NormalTok{ x }\OperatorTok{=} \DecValTok{5}\OperatorTok{;}
\KeywordTok{let}\NormalTok{ r }\OperatorTok{=} \OperatorTok{\&}\NormalTok{x}\OperatorTok{;}
\NormalTok{println}\OperatorTok{!(}\StringTok{"\{\}"}\OperatorTok{,} \OperatorTok{*}\NormalTok{r}\OperatorTok{);}        \CommentTok{// 5; println! otomatik deref de yapar: \{r\} de olur}
\end{Highlighting}
\end{Shaded}
\item
  Fonksiyondan referans dönmek (dangling referans yasak):

\begin{Shaded}
\begin{Highlighting}[]
\KeywordTok{fn} \FunctionTok{dangle}\OperatorTok{()} \OperatorTok{-\textgreater{}} \OperatorTok{\&}\DataTypeTok{String} \OperatorTok{\{}      \CommentTok{// E0106: missing lifetime specifier}
    \KeywordTok{let}\NormalTok{ s }\OperatorTok{=} \DataTypeTok{String}\OperatorTok{::}\FunctionTok{from}\OperatorTok{(}\StringTok{"x"}\OperatorTok{);}
    \OperatorTok{\&}\NormalTok{s}
\OperatorTok{\}}                              \CommentTok{// s burada drop edilir, referans sarkık kalırdı}
\end{Highlighting}
\end{Shaded}


  Çözüm: \texttt{String}`i doğrudan döndür (ownership transfer), referans değil
\item
  Metot çağrılarında otomatik borrow:

\begin{Shaded}
\begin{Highlighting}[]
\KeywordTok{let}\NormalTok{ s }\OperatorTok{=} \DataTypeTok{String}\OperatorTok{::}\FunctionTok{from}\OperatorTok{(}\StringTok{"hi"}\OperatorTok{);}
\NormalTok{s}\OperatorTok{.}\FunctionTok{len}\OperatorTok{();}            \CommentTok{// aslında (\&s).len() gibi çalışır, taşımaz}
\end{Highlighting}
\end{Shaded}
\end{itemize}

\textbf{GoT bağlantı haritası:}

\begin{verbatim}
borrowing ──┬── &T (immutable) ──► birden fazla OK
            └── &mut T (mutable) ──┬── tek seferde sadece 1 tane (E0499)
                                    └── & ile aynı anda olamaz (E0502)
reference ──► sahip değil ──► scope bitince orijinali drop etmez
NLL ──► referansın ömrü son kullanımda biter
dangling referans yasak ──► E0106 (lifetime, Konu 25'te detaylı)
\end{verbatim}

\begin{center}\rule{0.5\linewidth}{0.5pt}\end{center}

\subsection{🟣 KATMAN 3 --- DERIN TEKNIK NOT}\label{katman-3-—-derin-teknik-not-11}

\textbf{Adım 1 --- Bellek Seviyesi}

\begin{itemize}
\tightlist
\item
  Önce soru: ``Referans bellekte ne?'' Cevap: bir referans, işaret ettiği değerin bellek adresini tutan bir pointer'dır. \texttt{\&i32} 8 byte (64-bit adres); \texttt{\&str} gibi fat pointer'lar 16 byte (adres + uzunluk).
\item
  \texttt{\&s1} oluşturmak hiçbir veri kopyalamaz, hiçbir heap allocation yapmaz; sadece \texttt{s1}`in adresini taşıyan küçük bir değer üretir.
\item
  Fonksiyona \texttt{\&String} geçirmek, çağrılan fonksiyonun stack frame'ine sadece bu adresi kopyalar; \texttt{String}`in kendisi (heap'teki veri) yerinde kalır ve taşınmaz.
\item
  \texttt{\&mut} referans da aynı boyutta bir pointer'dır; fark bellekte değil, derleyicinin bu pointer üzerinden yazmaya izin vermesindedir.
\end{itemize}

\textbf{Adım 2 --- Compiler Davranışı}

\begin{itemize}
\tightlist
\item
  Önce soru: ``Borrow checker ne kontrol ediyor?'' Cevap: her değer için, herhangi bir anda ya (a) sınırsız sayıda \texttt{\&} referans ya da (b) tam olarak bir \texttt{\&mut} referans olabileceğini, ikisinin asla aynı anda olamayacağını doğrular.
\item
  Bu, data race'leri (veri yarışı) compile time'da önler: aynı veriye aynı anda hem okuma hem yazma erişimi (ya da iki yazma erişimi) asla oluşamaz.
\item
  NLL sayesinde borrow checker, referansın söz dizimsel scope'una değil gerçek son kullanım noktasına bakar; bu, eskiden hata veren birçok geçerli kodu artık kabul eder.
\item
  Fonksiyondan \texttt{\&} dönmek, lifetime annotation (yaşam süresi işaretleyici) gerektirir (Konu 25); işaretlenmezse ve derleyici çıkaramazsa \texttt{E0106} verir.
\item
  MIR'de referanslar \texttt{Ref} ve \texttt{\&mut} referanslar \texttt{Mut\ Ref} olarak, borrow checker'ın kendisi ayrı bir analiz aşaması (\texttt{borrowck}) olarak çalışır. Görmek için: \texttt{cargo\ rustc\ --\ --emit=mir} (borrow checker hatalarını görmek için zaten normal derleme yeterli).
\end{itemize}

\textbf{Adım 3 --- Hata Kodları}

\begin{itemize}
\tightlist
\item
  \textbf{E0596:} \texttt{mut} olmayan bir değişkene/referansa mutable borrow almaya çalıştın. Çözüm: değişkeni \texttt{let\ mut} yap.
\item
  \textbf{E0499:} aynı değere aynı anda iki \texttt{\&mut}. Çözüm: ilkini son kullandıktan sonra ikincisini al (NLL buna izin verir) veya scope'ları ayır.
\item
  \textbf{E0502:} \texttt{\&} ve \texttt{\&mut} aynı anda çakıştı. Çözüm: immutable referansların son kullanımını mutable referanstan önceye al.
\item
  \textbf{E0106:} fonksiyon dönüş tipinde referans var ama lifetime belirtilmemiş. Çözüm: ownership döndür veya lifetime annotation ekle (Konu 25).
\item
  Detay için: \texttt{rustc\ --explain\ E0502}.
\end{itemize}

\textbf{Adım 4 --- ``Neden Böyle?''}

\begin{itemize}
\tightlist
\item
  C/C++`ta bir pointer'ı kaç yere istersen kopyalayabilirsin; biri veriyi değiştirirken başka bir thread veya kod yolu aynı anda okursa data race oluşur --- bu klasik eşzamanlılık hatalarının kaynağıdır.
\item
  C++`ın \texttt{const\ T\&} ve \texttt{T\&} ayrımı kavramsal olarak benzer, ama derleyici aynı anda bir \texttt{T\&} ile bir \texttt{const\ T\&}`nin aynı veriye bakmasını engellemez; Rust bunu borrow checker ile zorunlu kılar.
\item
  Rust'ın kuralı ``aynı anda ya çoklu okuyucu ya tek yazıcı'' (reader-writer lock mantığına benzer), ama kilitleme (locking) runtime'da değil, compile time'da gerçekleşir --- çalışma zamanı maliyeti sıfırdır.
\item
  Trade-off: bazı geçerli-görünen desenler (örneğin bir struct'ın iki farklı alanına aynı anda mutable erişim) derleyiciyi ikna etmek için yeniden yazılmalıdır, ama karşılığında data race sınıfı hatalar derleme zamanında elenir.
\end{itemize}

\textbf{Adım 5 --- Rust 1.96.0 Farkı}

\begin{itemize}
\tightlist
\item
  Rust 1.96.0'daki değişiklikler (yeni Range* tipleri, \texttt{assert\_matches!} makrosu, Wasm linker güncellemeleri) bu konuyu etkilemiyor. Borrowing kuralları, NLL davranışı ve E0499/E0502/E0596/E0106 hata kodları 1.96.0'da değişmedi.
\item
  Rust 1.96.0'daki değişiklikler borrowing kurallarını veya borrow checker davranışını değiştirmiyor.
\end{itemize}

\begin{center}\rule{0.5\linewidth}{0.5pt}\end{center}

\subsection{🔴 KATMAN 4 --- SIK UNUTULANLAR + DERINLEŞTIRME}\label{katman-4-—-sik-unutulanlar-derinleştirme-11}

\textbf{Sık Unutulan 1: Referans almak taşımaz}

\begin{itemize}
\tightlist
\item
  Neden unutuluyor: Konu 11'den sonra her şeyin ``taşınacağı'' hissine kapılınabilir.
\item
  Yol 1 (analoji): \texttt{\&s} bir kitabı ödünç vermektir, \texttt{s} taşımak kitabı hediye etmektir. Ödünç veren kitabı geri alır.
\item
  Yol 2 (kod): \texttt{fn\ f(s:\ \&String)\ \{\}} çağrısından sonra orijinal değişken hâlâ kullanılabilir.
\item
  Yol 3 (hata): referans yerine değeri taşırsan (\texttt{fn\ f(s:\ String)}), sonraki kullanımda \texttt{E0382} görürsün (Konu 11).
\end{itemize}

\textbf{Sık Unutulan 2: \texttt{\&mut} tekil kuralı}

\begin{itemize}
\tightlist
\item
  Neden unutuluyor: ``sadece okuyorum, sorun olmaz'' diye iki \texttt{\&mut} birden alınabileceği sanılır.
\item
  Yol 1 (görsel): bir anda sadece bir kişi tahtaya yazabilir; birden fazla kişi aynı anda yazarsa karmaşa çıkar.
\item
  Yol 2 (kod): ilk \texttt{\&mut}`ı kullandıktan sonra (NLL ile ``öldükten'' sonra) ikincisini al.
\item
  Yol 3 (hata): \texttt{E0499:\ cannot\ borrow\ s\ as\ mutable\ more\ than\ once\ at\ a\ time}
\end{itemize}

\textbf{Sık Unutulan 3: \texttt{\&} ve \texttt{\&mut} karışması}

\begin{itemize}
\tightlist
\item
  Neden unutuluyor: ``bir tanesi okuma, diğeri yazma, farklı işler'' diye ayrı sorun sanılmaz.
\item
  Yol 1 (analoji): bir kişi deftere bakarken (okurken) başka biri aynı anda üzerine yazarsa, okuyan kişi yanlış/eksik veri görebilir.
\item
  Yol 2 (kod): immutable referansların \texttt{println!}`ini mutable referanstan önce bitir.
\item
  Yol 3 (hata): \texttt{E0502:\ cannot\ borrow\ s\ as\ mutable\ because\ it\ is\ also\ borrowed\ as\ immutable}
\end{itemize}

\textbf{Hazır Derinleştirme Prompt'u:}

\begin{verbatim}
Rust 1.96.0'da Borrowing (& ve &mut) konusunu en detaylı şekilde anlat. Şunları kapsa:

1. BELLEK SEVİYESİ: Bir referans bellekte kaç byte, ne tutar? &T ile &mut T'nin
   bellek temsili aynı mı? Referans almak neden heap allocation gerektirmiyor?

2. COMPILER DAVRANIŞI: Borrow checker "çoklu & veya tek &mut" kuralını nasıl
   uyguluyor? NLL (Non-Lexical Lifetimes) referansın ömrünü nasıl belirliyor?
   Data race'ler compile time'da nasıl önleniyor? MIR'de referanslar ve borrowck
   nasıl çalışır?

3. HATA KODLARI: E0596, E0499, E0502, E0106. Her biri için: ne zaman, neden,
   nasıl çözülür. `rustc --explain EXXXX` çıktısını da ekle.

4. C/C++ KARŞILAŞTIRMASI: Pointer/referans kopyalama, data race, const T& / T&.
   Rust'ın borrow checker'ı bu sorunları nasıl compile time'da çözüyor?

5. ANALOJİ: Günlük hayattan bir benzetme yap (örn: kitap ödünç verme, ortak defter).

6. RUST 1.96.0: Bu konuda son sürümde değişen bir şey var mı?

7. PRATİK: 3 tane kod örneği ver — kolay, orta, zor.

Açıklama Türkçe, teknik terimler İngilizce, kod yorumları Türkçe olsun.
Her bölümü ayrı başlıkla ver. Kod bloklarını çalıştırılabilir yaz.
\end{verbatim}

\begin{center}\rule{0.5\linewidth}{0.5pt}\end{center}

\subsection{📝 GÖREV (Boş dosyaya yaz)}\label{görev-boş-dosyaya-yaz-10}

\textbf{🟢 Kolay}

\begin{enumerate}
\def\labelenumi{\arabic{enumi}.}
\tightlist
\item
  Bir \texttt{String} oluştur; uzunluğunu \texttt{\&s} alan bir fonksiyonla hesapla ve orijinal değişkeni fonksiyondan sonra da yazdır.
\item
  Bir \texttt{mut\ String} oluştur; \texttt{\&mut} alan bir fonksiyonla ona \texttt{push\_str} ile metin ekle ve sonucu yazdır.
\end{enumerate}

\textbf{🟡 Orta}
3. Aynı \texttt{String}`e iki tane \texttt{\&} referans al, ikisini de aynı \texttt{println!}`de kullan.
4. \texttt{let\ mut\ s} ile başla; önce \texttt{\&s} al ve kullan, kullanımı bitince \texttt{\&mut\ s} al ve kullan (NLL'yi göster).

\textbf{🔴 Zor}
5. Bilerek aynı anda iki \texttt{\&mut\ s} almayı dene ve \texttt{E0499}`u gör; düzelt. Sonra bir \texttt{\&s} dururken \texttt{\&mut\ s} almayı dene ve \texttt{E0502}`yi gör; immutable referansın kullanımını öne alarak düzelt. Son olarak bir fonksiyondan \texttt{\&String} döndürmeyi dene (\texttt{dangle} örneği) ve \texttt{E0106}`yı gör.

\textbf{Başarı kriteri:} Hepsini hatasız yazarsan → öğrendin ✅

\subsubsection{✅ Kendini Test Et (2-3 saat sonra, sıfırdan)}\label{kendini-test-et-2-3-saat-sonra-sıfırdan-11}

\begin{enumerate}
\def\labelenumi{\arabic{enumi}.}
\tightlist
\item
  \texttt{\&} ve \texttt{\&mut} parametreli birer fonksiyon yaz.
\item
  ``Çoklu \texttt{\&} veya tek \texttt{\&mut}'' kuralını bir örnekle göster.
\item
  \texttt{E0499} ve \texttt{E0502}`nin farkını açıkla.
\end{enumerate}

\begin{center}\rule{0.5\linewidth}{0.5pt}\end{center}

\subsection{📚 KAYNAKÇA}\label{kaynakça-11}

\subsubsection{Resmi Kaynaklar}\label{resmi-kaynaklar-11}

\begin{itemize}
\tightlist
\item
  \href{https://doc.rust-lang.org/book/ch04-02-references-and-borrowing.html}{Rust Book --- References and Borrowing}
\item
  \href{https://doc.rust-lang.org/reference/types/pointer.html\#shared-references-}{Rust Reference --- References}
\end{itemize}

\subsubsection{Sürüm Notları}\label{sürüm-notları-11}

\begin{itemize}
\tightlist
\item
  \href{https://blog.rust-lang.org/2026/05/28/Rust-1.96.0/}{Rust Blog --- Announcing 1.96.0}
\item
  \href{https://doc.rust-lang.org/stable/releases.html\#version-1960-2026-05-28}{Rust 1.96.0 Release Notes}
\end{itemize}

\subsubsection{Hata Kodları}\label{hata-kodları-11}

\begin{itemize}
\tightlist
\item
  \href{https://doc.rust-lang.org/error_codes/E0596.html}{E0596}
\item
  \href{https://doc.rust-lang.org/error_codes/E0499.html}{E0499}
\item
  \href{https://doc.rust-lang.org/error_codes/E0502.html}{E0502}
\item
  \href{https://doc.rust-lang.org/error_codes/E0106.html}{E0106}
\end{itemize}

\subsubsection{Ek Kaynaklar}\label{ek-kaynaklar-11}

\begin{itemize}
\tightlist
\item
  \href{https://doc.rust-lang.org/rust-by-example/scope/borrow.html}{Rust By Example --- Borrowing}
\end{itemize}

\emph{Not: ``Çoklu immutable veya tek mutable referans'' kuralı, \texttt{E0499}, \texttt{E0502} mesajları ve NLL davranışı (son kullanıma göre ömür) aramada Rust Book üzerinden doğrulandı. Referans boyutları (8/16 byte), \texttt{E0596}/\texttt{E0106} ve MIR/borrowck detayları bilgime dayanıyor; \texttt{size\_of} ve hata mesajlarını kendi makinende çalıştırarak kontrol et.}

\begin{center}\rule{0.5\linewidth}{0.5pt}\end{center}
\end{document}
